% Arithmétique : division euclidienne et numération — Mathématiques 1ère TI — Cours signé OnBuch+
% Programme officiel MINESEC. Compiler avec : tectonic N21-arithmetique-numeration.tex
\newcommand{\DOCMATIERE}{Mathématiques}
\newcommand{\DOCNIVEAU}{1ère TI}
\newcommand{\DOCMODULE}{Relations et opérations fondamentales dans l'ensemble des nombres réels}
\newcommand{\DOCLECON}{N21}
\newcommand{\DOCTITRE}{Arithmétique : division euclidienne et numération}
% OnBuch+ — préambule « cours signés » ENRICHI (compilé avec tectonic / XeLaTeX)
% Système visuel « Pop » d'OnBuch+ : fond crème (jamais blanc), bordures encre
% épaisses, ombres « dures » décalées, titres Archivo Black en capitales.
% Version MATHÉMATIQUES : graphes (pgfplots/tikz), boîtes pédagogiques
% (définition, propriété, méthode, attention, le savais-tu, exercice résolu,
% exercice, TP, bilan), page de garde et signature OnBuch+ ancrée en bas.
%
% Macros attendues dans main.tex AVANT \input{preamble} :
%   \DOCMATIERE  \DOCNIVEAU  \DOCMODULE  \DOCLECON (numéro)  \DOCTITRE
\documentclass[11pt]{article}

\usepackage{fontspec}
\usepackage[french]{babel}
\usepackage[a4paper,margin=2cm,top=3.4cm,bottom=2.8cm,headheight=1.4cm,footskip=1.3cm]{geometry}
\usepackage{amsmath,amssymb,amsfonts}
\usepackage{mathtools} % \xmapsto, \coloneqq... souvent utilisés par les modèles
\usepackage{cancel} % simplifications barrées
% \abs{z} (module, valeur absolue) et \norm{u} (norme), utilisés par les modèles
\providecommand{\abs}{}\let\abs\relax\DeclarePairedDelimiter{\abs}{\lvert}{\rvert}
\providecommand{\norm}{}\let\norm\relax\DeclarePairedDelimiter{\norm}{\lVert}{\rVert}
\usepackage[table]{xcolor} % [table] : fournit aussi \rowcolors (alternance de lignes)
\usepackage{pagecolor}
\usepackage{tikz}
\usetikzlibrary{arrows.meta,positioning,calc,shapes.geometric,shapes.misc,shapes.arrows,decorations.pathmorphing,decorations.pathreplacing,decorations.markings,patterns,shadows,mindmap,trees,fit,backgrounds,matrix,intersections,angles,quotes}
\usepackage{pgfplots}
\usepackage{pgf-pie} % diagrammes circulaires \pie (statistiques)
\pgfplotsset{compat=1.18}
\usepgfplotslibrary{fillbetween,groupplots} % aires entre courbes (calcul intégral)
\usepackage{enumitem}
\usepackage{fancyhdr}
% [most] charge toutes les bibliothèques tcolorbox (listings, minted, raster,
% documentation...) et gonfle inutilement la mémoire TeX sur un document long
% (~300 boîtes) : on ne charge que ce qui est réellement utilisé.
\usepackage[skins,breakable]{tcolorbox}
\usepackage{graphicx}
\usepackage{colortbl}
\usepackage{booktabs}
\usepackage{tabularx}
\usepackage{diagbox} % case d'en-tête diagonale des tableaux à double entrée
% Colonnes extensibles centrée / gauche / droite (C, L, R), souvent utilisées par les modèles
\newcolumntype{C}{>{\centering\arraybackslash}X}
\newcolumntype{L}{>{\raggedright\arraybackslash}X}
\newcolumntype{R}{>{\raggedleft\arraybackslash}X}
\usepackage{array}
\usepackage{multicol}
\usepackage{siunitx}
% parse-numbers=false : accepte aussi \SI{2,5 \times 10^{-3}}{...}
\sisetup{locale=FR,output-decimal-marker={,},group-minimum-digits=4,parse-numbers=false}
\DeclareSIUnit\ohm{\ensuremath{\symup{\Omega}}} % U+2126 absent de la police
\usepackage{qrcode}
% \degree : commande fréquemment utilisée par les modèles pour un angle ou une
% température en dehors de \SI{}{} (non fournie par les paquets chargés).
\providecommand{\degree}{^{\circ}}
% \qed (fin de démonstration), utilisé par les modèles sans amsthm
\providecommand{\qed}{\hfill$\square$}
% \barre : double barre des valeurs interdites dans un tableau de variations (tkz-tab)
\providecommand{\barre}{\ensuremath{\|}}
% environnement proof (amsthm non chargé) utilisé par les modèles
\ifdefined\proof\else\newenvironment{proof}[1][Démonstration]{\par\noindent\textit{#1.}\ }{\qed\par}\fi
\usepackage{lastpage}
\usepackage{hyperref}
\hypersetup{hidelinks}

% ============================================================
% Palette « Pop » OnBuch+ — mêmes hex que la classe P (plus_ui.dart)
% ============================================================
\definecolor{poporange}{HTML}{F59321}
\definecolor{popdark}{HTML}{E07A0C}
\definecolor{popcream}{HTML}{FFF6E8}
\definecolor{popink}{HTML}{1C1714}
\definecolor{poppurple}{HTML}{7A5AE0}
\definecolor{popgreen}{HTML}{2E9E6A}
\definecolor{poppink}{HTML}{E0457B}
\definecolor{popblue}{HTML}{2F7DE1}
\definecolor{popgold}{HTML}{C98A12}
\definecolor{popmuted}{HTML}{8A7F76}
\definecolor{popline}{HTML}{EADFD2}
\definecolor{popgray}{HTML}{8A7F76}
\colorlet{popgrey}{popgray}
% Alias tolérants (le modèle nomme parfois les couleurs différemment)
\colorlet{popwhite}{white}
\colorlet{popblack}{popink}
\colorlet{popred}{poppink}
\colorlet{popyellow}{popgold}
% Teintes claires pour fonds de tableaux / zones
\colorlet{popblueL}{popblue!10!white}
\colorlet{poporangeL}{poporange!14!white}
\colorlet{popgreenL}{popgreen!12!white}
\colorlet{poppurpleL}{poppurple!12!white}
\colorlet{poppinkL}{poppink!12!white}
\colorlet{popgoldL}{popgold!14!white}

\pagecolor{popcream}

% ============================================================
% Typographie
% ============================================================
\defaultfontfeatures{Ligatures=TeX}
\setmainfont{PlusJakartaSans-Medium.ttf}[Path=../../../fonts/,BoldFont=PlusJakartaSans-Bold.ttf]
% Maths en sans-serif (Fira Math) avec chiffres et lettres droites de Plus
% Jakarta, pour que les formules se fondent dans le texte.
\usepackage[math-style=ISO,bold-style=ISO]{unicode-math}
\setmathfont{FiraMath-Regular.otf}[Path=../../../fonts/]
\setmathfont{PlusJakartaSans-Medium.ttf}[Path=../../../fonts/,range={up/num,up/{Latin,latin},\mathup}]
\usepackage{icomma} % virgule décimale sans espace en mode math
% Caractères mathématiques Unicode absents des polices de texte (Plus Jakarta,
% Archivo Black) : rendus par leur équivalent mathématique, en texte comme en
% formule — sinon ils s'affichent en carré vide (ex. « Arithmétique dans ℤ »).
\usepackage{newunicodechar}
\newunicodechar{ℝ}{\ensuremath{\mathbb{R}}}
\newunicodechar{ℤ}{\ensuremath{\mathbb{Z}}}
\newunicodechar{ℂ}{\ensuremath{\mathbb{C}}}
\newunicodechar{ℕ}{\ensuremath{\mathbb{N}}}
\newunicodechar{ℚ}{\ensuremath{\mathbb{Q}}}
\newunicodechar{𝔻}{\ensuremath{\mathbb{D}}}
\newunicodechar{α}{\ensuremath{\alpha}}
\newunicodechar{β}{\ensuremath{\beta}}
\newunicodechar{γ}{\ensuremath{\gamma}}
\newunicodechar{δ}{\ensuremath{\delta}}
\newunicodechar{ε}{\ensuremath{\varepsilon}}
\newunicodechar{θ}{\ensuremath{\theta}}
\newunicodechar{λ}{\ensuremath{\lambda}}
\newunicodechar{μ}{\ensuremath{\mu}}
\newunicodechar{σ}{\ensuremath{\sigma}}
\newunicodechar{τ}{\ensuremath{\tau}}
\newunicodechar{φ}{\ensuremath{\varphi}}
\newunicodechar{ψ}{\ensuremath{\psi}}
\newunicodechar{ω}{\ensuremath{\omega}}
\newunicodechar{Σ}{\ensuremath{\Sigma}}
% majuscules produites par \MakeUppercase dans les titres (α-aminés -> Α)
\newunicodechar{Α}{\ensuremath{\alpha}}
\newunicodechar{Β}{\ensuremath{\beta}}
\newunicodechar{Γ}{\ensuremath{\Gamma}}
\newunicodechar{Δ}{\ensuremath{\Delta}}
\newunicodechar{Ε}{\ensuremath{\varepsilon}}
\newunicodechar{Θ}{\ensuremath{\Theta}}
\newunicodechar{Λ}{\ensuremath{\Lambda}}
\newunicodechar{Μ}{\ensuremath{\mu}}
\newunicodechar{Τ}{\ensuremath{\tau}}
\newunicodechar{Φ}{\ensuremath{\Phi}}
\newunicodechar{Ψ}{\ensuremath{\Psi}}
\newunicodechar{Ω}{\ensuremath{\Omega}}
\newunicodechar{↦}{\ensuremath{\mapsto}}
\newunicodechar{∈}{\ensuremath{\in}}
\newunicodechar{∉}{\ensuremath{\notin}}
\newunicodechar{∧}{\ensuremath{\wedge}}
\newunicodechar{⇔}{\ensuremath{\Leftrightarrow}}
\newunicodechar{⟺}{\ensuremath{\Longleftrightarrow}}
\newunicodechar{⇒}{\ensuremath{\Rightarrow}}
\newunicodechar{⟹}{\ensuremath{\Longrightarrow}}
\newunicodechar{∘}{\ensuremath{\circ}}
\newunicodechar{∪}{\ensuremath{\cup}}
\newunicodechar{∩}{\ensuremath{\cap}}
\newunicodechar{≡}{\ensuremath{\equiv}}
\newunicodechar{⊂}{\ensuremath{\subset}}
\newunicodechar{⊥}{\ensuremath{\perp}}
\newunicodechar{∃}{\ensuremath{\exists}}
\newunicodechar{∀}{\ensuremath{\forall}}
\newunicodechar{∛}{\ensuremath{\sqrt[3]{\,}}}
\newunicodechar{∜}{\ensuremath{\sqrt[4]{\,}}}
\newunicodechar{⃗}{\ensuremath{{}^{\to}}}
\newunicodechar{ⁿ}{\ifmmode^{n}\else\textsuperscript{\ensuremath{n}}\fi}
\newunicodechar{ˣ}{\ifmmode^{x}\else\textsuperscript{\ensuremath{x}}\fi}
\newunicodechar{ᵇ}{\ifmmode^{b}\else\textsuperscript{\ensuremath{b}}\fi}
\newunicodechar{ʳ}{\ifmmode^{r}\else\textsuperscript{\ensuremath{r}}\fi}
\newunicodechar{ᵘ}{\ifmmode^{u}\else\textsuperscript{\ensuremath{u}}\fi}
\newunicodechar{ʸ}{\ifmmode^{y}\else\textsuperscript{\ensuremath{y}}\fi}
\newunicodechar{ᵃ}{\ifmmode^{a}\else\textsuperscript{\ensuremath{a}}\fi}
\newunicodechar{ᵉ}{\ifmmode^{e}\else\textsuperscript{\ensuremath{e}}\fi}
\newunicodechar{ᵏ}{\ifmmode^{k}\else\textsuperscript{\ensuremath{k}}\fi}
\newunicodechar{ᵖ}{\ifmmode^{p}\else\textsuperscript{\ensuremath{p}}\fi}
\newunicodechar{ᴺ}{\ifmmode^{N}\else\textsuperscript{\ensuremath{N}}\fi}
\newunicodechar{ᶜ}{\ifmmode^{c}\else\textsuperscript{\ensuremath{c}}\fi}
\newunicodechar{ʰ}{\ifmmode^{h}\else\textsuperscript{\ensuremath{h}}\fi}
\newunicodechar{ᵠ}{\ifmmode^{q}\else\textsuperscript{\ensuremath{q}}\fi}
\newunicodechar{ₙ}{\ifmmode_{n}\else\textsubscript{\ensuremath{n}}\fi}
\newunicodechar{ₐ}{\ifmmode_{a}\else\textsubscript{\ensuremath{a}}\fi}
\newunicodechar{ᵢ}{\ifmmode_{i}\else\textsubscript{\ensuremath{i}}\fi}
\newunicodechar{ᵣ}{\ifmmode_{r}\else\textsubscript{\ensuremath{r}}\fi}
\newunicodechar{ₖ}{\ifmmode_{k}\else\textsubscript{\ensuremath{k}}\fi}
\newunicodechar{ᵦ}{\ifmmode_{\beta}\else\textsubscript{\ensuremath{\beta}}\fi}
\newunicodechar{ₓ}{\ifmmode_{x}\else\textsubscript{\ensuremath{x}}\fi}
\newfontfamily\popdisplay{ArchivoBlack-Regular.ttf}[Path=../../../fonts/]
\newfontfamily\poplogo{SpaceGrotesk-Bold.ttf}[Path=../../../fonts/]
\newfontfamily\popsemi{PlusJakartaSans-SemiBold.ttf}[Path=../../../fonts/]
\newfontfamily\popxbold{PlusJakartaSans-ExtraBold.ttf}[Path=../../../fonts/]
\newfontfamily\popxboldls{PlusJakartaSans-ExtraBold.ttf}[Path=../../../fonts/,LetterSpace=3.0]

\setlist[enumerate]{leftmargin=1.7em,itemsep=5pt,topsep=4pt}
\setlist[itemize]{leftmargin=1.7em,itemsep=4pt,topsep=4pt,label={\color{poporange}\textbullet}}
\setlength{\parindent}{0pt}
\setlength{\parskip}{5pt}
\renewcommand{\arraystretch}{1.25}

% pgfplots : style Pop par défaut
\pgfplotsset{
  every axis/.append style={axis line style={popink,line width=1pt},tick style={popink},
    grid style={popline,line width=0.5pt},label style={font=\small},tick label style={font=\footnotesize}},
  popcurve/.style={poporange,line width=1.8pt,no markers,smooth},
}
% Même style disponible dans un \draw TikZ simple (hors pgfplots)
\tikzset{popcurve/.style={poporange,line width=1.8pt}}

% ============================================================
% Pill / squiggle
% ============================================================
\newcommand{\poppill}[3]{%
  \tikz[baseline=-0.55ex]\node[draw=popink,line width=1.2pt,rounded corners=2.6pt,%
    fill=#2,inner xsep=6.5pt,inner ysep=3pt]{%
    \popxboldls\fontsize{7}{7}\selectfont\color{#3}\MakeUppercase{#1}};%
}
\newcommand{\popsquiggle}[1]{%
  \tikz[baseline=-0.4ex]{
    \draw[#1,line width=2.6pt,line cap=round]
      plot [smooth] coordinates {(0,0.09) (0.34,0.01) (0.68,0.21) (1.02,0.01) (1.36,0.10)};
  }%
}
% Mot-clé surligné (marqueur orange sous le texte)
\newcommand{\cle}[1]{{\popxbold\color{popdark}#1}}

% ============================================================
% En-tête / pied
% ============================================================
\pagestyle{fancy}
\fancyhf{}
\renewcommand{\headrulewidth}{0pt}
\renewcommand{\footrulewidth}{0pt}
\lhead{\raisebox{-0.15ex}{\poplogo\fontsize{13}{13}\selectfont\color{popink}OnBuch\color{poporange}+}}
\rhead{\poppill{\DOCMATIERE\ \textbullet\ \DOCNIVEAU\ \textbullet\ Leçon \DOCLECON}{white}{popink}}
\lfoot{\popxbold\fontsize{7.6}{7.6}\selectfont\color{popmuted}\MakeUppercase{Cours signé OnBuch+}\ \textbullet\ {\popsemi\DOCTITRE}}
\rfoot{\tikz[baseline=-0.6ex]\node[rounded corners=8.5pt,fill=popink,minimum height=17pt,minimum width=17pt,inner xsep=5pt,inner sep=0]{\popxbold\fontsize{8}{8}\selectfont\color{popcream}\thepage\,/\,\pageref*{LastPage}};}

% ============================================================
% Titres
% ============================================================
\newcommand{\chaptitre}[2]{%
  \begin{tcolorbox}[enhanced,colback=poporange,colframe=popink,boxrule=2.6pt,arc=11pt,
      drop shadow=popink,left=15pt,right=15pt,top=12pt,bottom=11pt,before skip=3pt,after skip=18pt]
    \poppill{#2}{popink}{popcream}\par\vspace{9pt}
    {\raggedright\hyphenpenalty=10000\exhyphenpenalty=10000\popdisplay\fontsize{22}{24}\selectfont\color{popcream}\MakeUppercase{#1}\par}
  \end{tcolorbox}
}
% Section numérotée : pastille orange avec le numéro + titre Archivo
\newcounter{popsec}
\newcommand{\coursec}[1]{%
  \stepcounter{popsec}\setcounter{popsub}{0}%
  \par\vspace{16pt}\noindent
  \begin{minipage}{\linewidth}
  \tikz[baseline=-0.7ex]\node[draw=popink,line width=1.6pt,fill=poporange,rounded corners=4pt,minimum size=22pt,inner sep=0]{\popdisplay\fontsize{12}{12}\selectfont\color{popcream}\thepopsec};%
  \hspace{8pt}{\popdisplay\fontsize{15}{17}\selectfont\color{popink}\MakeUppercase{#1}}\par
  \vspace{4pt}\noindent{\color{poporange}\rule{36pt}{3.6pt}}
  \end{minipage}\par\nopagebreak\vspace{8pt}
}
\newcounter{popsub}
\newcommand{\courssub}[1]{%
  \stepcounter{popsub}%
  \par\vspace{9pt}\noindent{\popxbold\fontsize{11.5}{13}\selectfont\color{popink}{\color{poporange}\thepopsec.\thepopsub}\hspace{6pt}#1}\par\nopagebreak\vspace{4pt}
}

% ============================================================
% Boîtes pédagogiques — même recette : fond blanc, bordure épaisse colorée,
% ombre dure encre, badge plein. Titre optionnel [#1].
% ============================================================
\tcbset{popbase/.style={breakable,enhanced,colback=white,boxrule=2.3pt,arc=9pt,
  shadow={3pt}{-3pt}{0pt}{popink},left=14pt,right=14pt,top=11pt,bottom=11pt,before skip=11pt,after skip=15pt,
  fontupper=\popsemi\fontsize{10.6}{14.5}\selectfont\color{popink},
  fontlower=\popsemi\fontsize{10.6}{14.5}\selectfont\color{popink}}}
% Coche dessinée (le glyphe ✓ n'existe pas dans les polices du document)
\newcommand{\popcheck}[1]{\tikz[baseline=-0.5ex]\draw[#1,line width=1.6pt,line cap=round,line join=round](-0.13,0.0)--(-0.04,-0.09)--(0.13,0.1);}
\newcommand{\popboxhead}[3]{\poppill{#1}{#2}{white}\ifx\relax#3\relax\else\hspace{6pt}{\popxbold\fontsize{10.6}{12}\selectfont\color{popink}#3}\fi\par\vspace{7pt}}

\newtcolorbox{aretenir}[1][]{popbase,colframe=popgreen,before upper={\popboxhead{À retenir}{popgreen}{#1}}}
\newtcolorbox{exemplebox}[1][]{popbase,colframe=poporange,before upper={\popboxhead{Exemple}{poporange}{#1}}}
\newtcolorbox{definition}[1][]{popbase,colframe=popblue,before upper={\popboxhead{Définition}{popblue}{#1}}}
\newtcolorbox{propriete}[1][]{popbase,colframe=poppurple,before upper={\popboxhead{Propriété}{poppurple}{#1}}}
\newtcolorbox{methode}[1][]{popbase,colframe=popgold,before upper={\popboxhead{Méthode}{popgold}{#1}}}
\newtcolorbox{attention}[1][]{popbase,colframe=poppink,before upper={\popboxhead{Attention}{poppink}{#1}}}
\newtcolorbox{experience}[1][]{popbase,colframe=popblue,colback=popblueL,before upper={\popboxhead{Expérience}{popblue}{#1}}}
% « Le savais-tu ? » : carte violette pleine (contexte, culture, Cameroun)
\newtcolorbox{savaistu}[1][]{popbase,colback=poppurple,colframe=popink,
  fontupper=\popsemi\fontsize{10.4}{14.2}\selectfont\color{white},
  before upper={\poppill{Le savais-tu\,?}{white}{poppurple}\ifx\relax#1\relax\else\hspace{6pt}{\popxbold\color{white}#1}\fi\par\vspace{7pt}}}
% Exercice résolu : énoncé en haut, \tcblower puis la solution sur fond crème
\newcounter{popexo}\newcounter{popexr}
\newtcolorbox{exoresolu}[1][]{popbase,colframe=poporange,colbacklower=popcream,
  segmentation style={popink,line width=1.2pt,dashed},
  before upper={\stepcounter{popexr}\popboxhead{Exercice résolu \thepopexr}{poporange}{#1}},
  before lower={\poppill{Solution}{popink}{popcream}\par\vspace{6pt}}}
% Exercice (non résolu en place) — la correction est dans la partie Corrigés
\newtcolorbox{exercice}[1][]{popbase,colframe=popink,
  before upper={\stepcounter{popexo}\popboxhead{Exercice \thepopexo}{popink}{#1}}}
% Corrigé
\newtcolorbox{corrige}[1][]{popbase,colframe=popgreen,colback=popgreenL,
  before upper={\popboxhead{Corrigé}{popgreen}{#1}}}
% Objectifs / prérequis / bilan
\newtcolorbox{objectifs}{popbase,colframe=popink,colback=poporangeL,
  before upper={\popboxhead{Objectifs de la leçon}{popink}{}}}
\newtcolorbox{prerequis}{popbase,colframe=popink,colback=popblueL,
  before upper={\popboxhead{Prérequis}{popblue}{}}}
\newtcolorbox{bilan}{popbase,colframe=popink,colback=white,boxrule=2.8pt,
  before upper={\popboxhead{Fiche bilan}{poporange}{L'essentiel à connaître pour l'examen}}}
% Figure encadrée avec légende
\newtcolorbox{popfigure}[1][]{enhanced,breakable=false,colback=white,colframe=popline,boxrule=1.4pt,arc=8pt,
  left=10pt,right=10pt,top=10pt,bottom=8pt,before skip=10pt,after skip=12pt,halign=center,
  lower separated=false,halign lower=center,
  fontlower=\popsemi\fontsize{9}{11.5}\selectfont\color{popmuted},
  #1}
\newcommand{\legende}[1]{\par\vspace{6pt}{\popsemi\fontsize{9}{11.5}\selectfont\color{popmuted}#1}}

% ============================================================
% Signature OnBuch+
% ============================================================
\newcommand{\poplogoinline}[1]{{\poplogo\fontsize{#1}{#1}\selectfont\color{popink}OnBuch\color{poporange}+}}
\newcommand{\signatureonbuch}{%
  \begin{tcolorbox}[enhanced,colback=popink,colframe=popink,boxrule=2.2pt,arc=10pt,
      drop shadow=poporange,left=14pt,right=14pt,top=10pt,bottom=10pt,before skip=0pt,after skip=0pt]
    \tikz[baseline=-0.7ex]\node[circle,fill=popgreen,draw=popcream,line width=1.3pt,minimum size=22pt,inner sep=0]{\popcheck{white}};%
    \hspace{9pt}\begin{minipage}[c]{0.62\linewidth}
      {\popxbold\fontsize{12}{13}\selectfont\color{popcream}Signé\ \textbullet\ {\poplogo OnBuch\color{poporange}+}}\par\vspace{2pt}
      {\raggedright\popsemi\fontsize{8}{10.5}\selectfont\color{popline}\MakeUppercase{Équipe pédagogique OnBuch+}\\\MakeUppercase{Conforme au programme officiel MINESEC}\par}
    \end{minipage}\hfill
    {\poplogo\fontsize{22}{22}\selectfont\color{popcream}OnBuch\color{poporange}+}
  \end{tcolorbox}%
}

% ============================================================
% Page de garde
% #1 titre · #2 matière · #3 niveau · #4 module · #5 n° leçon · #6 durée ·
% #7 description / objectifs (texte ou itemize)
% ============================================================
\newcommand{\pagegarde}[7]{%
  \thispagestyle{empty}
  \newgeometry{a4paper,margin=1.7cm,top=1.6cm,bottom=1.4cm}%
  \begin{tikzpicture}[remember picture,overlay]
    \fill[poporange!18] ([xshift=-3.2cm,yshift=-3.6cm]current page.north east) circle (4.6cm);
    \fill[poppurple!14] ([xshift=2.4cm,yshift=7.6cm]current page.south west) circle (3.8cm);
  \end{tikzpicture}%
  {\poplogo\fontsize{30}{30}\selectfont\color{popink}OnBuch\color{poporange}+}\hfill
  \raisebox{0.6ex}{\poppill{Cours signé \textbullet\ Premium}{popink}{popcream}}\par
  \vspace{1pt}\popsquiggle{poppurple}\par
  \vspace{8pt}
  \begin{tcolorbox}[enhanced,colback=poporange,colframe=popink,boxrule=2.8pt,arc=13pt,
      drop shadow=popink,left=18pt,right=18pt,top=12pt,bottom=13pt,after skip=10pt]
    \poppill{#2 \textbullet\ #3}{popink}{popcream}\hspace{5pt}\poppill{#4}{white}{popink}\par\vspace{8pt}
    {\popdisplay\fontsize{14}{14}\selectfont\color{popink}LEÇON #5}\par\vspace{4pt}
    {\raggedright\hyphenpenalty=10000\exhyphenpenalty=10000\popdisplay\fontsize{25}{27}\selectfont\color{popcream}\MakeUppercase{#1}\par}
    \vspace{8pt}
    {\popxbold\fontsize{9.5}{9.5}\selectfont\color{popink}\MakeUppercase{Durée conseillée : #6}}
  \end{tcolorbox}
  \begin{tcolorbox}[enhanced,colback=white,colframe=popink,boxrule=2.2pt,arc=10pt,
      drop shadow=popink,left=16pt,right=16pt,top=10pt,bottom=10pt,after skip=8pt]
    \poppill{Dans cette leçon}{poppurple}{white}\par\vspace{5pt}
    \popsemi\fontsize{9.3}{12.6}\selectfont\color{popink}#7
  \end{tcolorbox}
  \noindent\begin{minipage}[t]{0.487\linewidth}
    \begin{tcolorbox}[enhanced,colback=popgreen,colframe=popink,boxrule=2.2pt,arc=10pt,
        drop shadow=popink,left=11pt,right=11pt,top=10pt,bottom=10pt,height=3.1cm,valign=center]
      \tcbox[colback=white,colframe=popink,boxrule=1.3pt,arc=5pt,boxsep=0pt,nobeforeafter,
        left=3pt,right=3pt,top=3pt,bottom=3pt,valign=center]{\qrcode[height=1.9cm]{https://play.google.com/store/apps/details?id=com.luvvix.onbuch&hl=fr}}%
      \hspace{8pt}\begin{minipage}[c]{0.5\linewidth}
        {\popdisplay\fontsize{11}{12.5}\selectfont\color{white}\MakeUppercase{Télécharge OnBuch}}\par\vspace{4pt}
        {\raggedright\popsemi\fontsize{8.4}{11}\selectfont\color{white}Gratuit \textbullet\ Google Play\\Tuteur IA, annales, concours, résultats\par}
      \end{minipage}
    \end{tcolorbox}
  \end{minipage}\hfill
  \begin{minipage}[t]{0.487\linewidth}
    \begin{tcolorbox}[enhanced,colback=poppurple,colframe=popink,boxrule=2.2pt,arc=10pt,
        drop shadow=popink,left=11pt,right=11pt,top=10pt,bottom=10pt,height=3.1cm,valign=center]
      \tikz[baseline=-0.7ex]\node[circle,fill=white,draw=popink,line width=1.3pt,minimum size=1.6cm,inner sep=0]{\popdisplay\fontsize{24}{24}\selectfont\color{poppurple}+};%
      \hspace{8pt}\begin{minipage}[c]{0.6\linewidth}
        {\popdisplay\fontsize{11}{12.5}\selectfont\color{white}\MakeUppercase{Passe à OnBuch+}}\par\vspace{4pt}
        {\raggedright\popsemi\fontsize{8.4}{11}\selectfont\color{white}Tous les cours signés, Tuteur IA illimité, fiches PDF — onglet OnBuch+ de l'app\par}
      \end{minipage}
    \end{tcolorbox}
  \end{minipage}
  \vspace{8pt}
  \popcontenu
  \vfill
  \signatureonbuch
  \restoregeometry
  \newpage
}

% Rangée « Ce cours contient » (page de garde)
\newcommand{\poptile}[3]{% #1 couleur #2 gros texte #3 légende
  \begin{tcolorbox}[enhanced,colback=#1,colframe=popink,boxrule=2pt,arc=9pt,shadow={2.5pt}{-2.5pt}{0pt}{popink},
      left=6pt,right=6pt,top=9pt,bottom=9pt,halign=center,valign=center,height=1.75cm,width=\linewidth,nobeforeafter]
    {\popdisplay\fontsize{12.5}{13}\selectfont\color{white}\MakeUppercase{#2}}\par\vspace{3pt}
    {\centering\popsemi\fontsize{7.8}{9.5}\selectfont\color{white}#3\par}
  \end{tcolorbox}}
\newcommand{\popcontenu}{%
  \par\noindent{\popxboldls\fontsize{7.5}{7.5}\selectfont\color{popmuted}CE COURS SIGNÉ CONTIENT}\par\vspace{7pt}
  \noindent\begin{minipage}[t]{0.235\linewidth}\poptile{popblue}{Cours}{complet, illustré, pas à pas}\end{minipage}\hfill
  \begin{minipage}[t]{0.235\linewidth}\poptile{poporange}{Méthodes}{et exercices résolus}\end{minipage}\hfill
  \begin{minipage}[t]{0.235\linewidth}\poptile{poppink}{Exercices}{type examen + corrigés}\end{minipage}\hfill
  \begin{minipage}[t]{0.235\linewidth}\poptile{popgreen}{Fiche bilan}{l'essentiel à retenir}\end{minipage}\par}

% Sommaire
\newtcolorbox{sommaire}{enhanced,colback=white,colframe=popink,boxrule=2.2pt,arc=10pt,
  drop shadow=popink,left=18pt,right=18pt,top=13pt,bottom=13pt,before skip=6pt,after skip=20pt,
  before upper={\poppill{Sommaire}{poppurple}{white}\par\vspace{9pt}\popsemi\fontsize{10.6}{14.5}\selectfont\color{popink}}}

% Signature de fin de cours — poussée tout en bas de la dernière page
\newcommand{\findecours}{%
  \par\vfill\nopagebreak
  \begin{center}\popsquiggle{poporange}\end{center}\vspace{4pt}
  \signatureonbuch
}
\AtBeginDocument{\def\llbracket{\mathopen{[\mkern-3.5mu[}}\def\rrbracket{\mathclose{]\mkern-3.5mu]}}} % intervalles d'entiers (glyphe ⟦ absent de la police)
% Glyphes absents de Fira Math : redéfinis à partir de glyphes disponibles
\AtBeginDocument{%
  \def\setminus{\mathbin{\backslash}}\let\smallsetminus\setminus
  \def\vdots{\mathord{\vbox{\baselineskip3.2pt\lineskiplimit0pt\kern4pt\hbox{.}\hbox{.}\hbox{.}}}}%
  \def\ddots{\mathinner{\mkern1mu\raise6.5pt\hbox{.}\mkern2mu\raise3.6pt\hbox{.}\mkern2mu\raise0.7pt\hbox{.}\mkern1mu}}%
  \def\triangle{\mathord{\scalebox{0.9}{$\symup{\Delta}$}}}%
  \def\nabla{\mathord{\raisebox{\depth}{\scalebox{1}[-1]{$\symup{\Delta}$}}}}%
  \def\checkmark{\popcheck{popdark}}%
  \def\star{\mathbin{\ast}}\def\bigstar{\mathord{\ast}}%
  \def\diamond{\mathbin{\scalebox{0.6}{$\Diamond$}}}%
  \def\bigcup{\mathop{\mathchoice{\raisebox{-0.3em}{\scalebox{1.7}{$\cup$}}}{\scalebox{1.25}{$\cup$}}{\cup}{\cup}}\displaylimits}%
  \def\bigcap{\mathop{\mathchoice{\raisebox{-0.3em}{\scalebox{1.7}{$\cap$}}}{\scalebox{1.25}{$\cap$}}{\cap}{\cap}}\displaylimits}%
  \def\lesseqqgtr{\mathrel{\lessgtr}}\def\bigoplus{\mathop{\oplus}\displaylimits}%
}
\providecommand{\ssi}{si et seulement si} % abréviation fréquente
\AtBeginDocument{\providecommand{\wideparen}{\overparen}} % arc de cercle

%% FIGFIX-BEGIN v5 — correctifs globaux des figures (docs/onbuch-generation/tools/figfix)
\usepackage{adjustbox}\usepackage{etoolbox}\tcbuselibrary{raster}
% toute figure TikZ/pgfplots plus large que la ligne est réduite à la largeur disponible
\BeforeBeginEnvironment{tikzpicture}{\begin{adjustbox}{max width=\linewidth}}
\AfterEndEnvironment{tikzpicture}{\end{adjustbox}}
% légendes pgfplots lisibles par-dessus les courbes
\pgfplotsset{every axis/.append style={legend style={fill=white,fill opacity=0.92,text opacity=1,draw=black!15,inner sep=3pt}}}
% étiquettes d'arêtes (syntaxe edge["..."]) avec fond blanc
\tikzset{every edge quotes/.append style={fill=white,inner sep=1.2pt}}
%% FIGFIX-END
\begin{document}

\pagegarde{Arithmétique : division euclidienne et numération}{Mathématiques}{1ère TI}{Relations et opérations fondamentales dans l'ensemble des nombres réels}{N21}{8 h}{Cette leçon étudie la division euclidienne dans ℕ et l'écriture des entiers naturels dans une base quelconque a (2 ≤ a ≤ 10). Elle montre comment passer d'une base à l'autre et comment additionner ou multiplier directement dans une base donnée, compétences clés pour l'informatique et les épreuves du Probatoire TI.
\vspace{4pt}
\begin{multicols}{2}\small
\begin{itemize}
\item À la fin de cette leçon, je sais déterminer le quotient et le reste d'une division euclidienne dans ℕ
\item À la fin de cette leçon, je sais traduire une division euclidienne par la relation a = bq + r avec 0 ≤ r < b
\item À la fin de cette leçon, je sais décomposer un entier naturel en somme de puissances ordonnées d'une base a
\item À la fin de cette leçon, je sais écrire un entier naturel dans un système de numération de base a par divisions successives
\item À la fin de cette leçon, je sais convertir un entier d'une base quelconque vers la base 10 et inversement
\item À la fin de cette leçon, je sais passer directement d'une base à une autre (notamment entre bases 2, 4 et 8 par regroupement de chiffres)
\end{itemize}\end{multicols}}

\begin{sommaire}
\begin{multicols}{2}
\begin{enumerate}
\item Division euclidienne dans ℕ
\item Numération de base a : principe et décomposition
\item Écriture d'un entier en base a : conversions avec la base 10
\item Passage direct d'une base à l'autre
\item Addition en base a
\item Multiplication en base a
\item Activité d'intégration
\item Méthodes et exercices résolus
\item Exercices
\item Corrigés des exercices
\item Fiche bilan
\end{enumerate}
\end{multicols}
\end{sommaire}

\begin{objectifs}
\begin{itemize}
\item À la fin de cette leçon, je sais déterminer le quotient et le reste d'une division euclidienne dans ℕ
\item À la fin de cette leçon, je sais traduire une division euclidienne par la relation a = bq + r avec 0 ≤ r < b
\item À la fin de cette leçon, je sais décomposer un entier naturel en somme de puissances ordonnées d'une base a
\item À la fin de cette leçon, je sais écrire un entier naturel dans un système de numération de base a par divisions successives
\item À la fin de cette leçon, je sais convertir un entier d'une base quelconque vers la base 10 et inversement
\item À la fin de cette leçon, je sais passer directement d'une base à une autre (notamment entre bases 2, 4 et 8 par regroupement de chiffres)
\item À la fin de cette leçon, je sais additionner deux entiers écrits dans la même base avec gestion des retenues
\item À la fin de cette leçon, je sais multiplier deux entiers écrits dans la même base
\item À la fin de cette leçon, je sais vérifier un calcul en base a en le convertissant en base 10
\end{itemize}
\end{objectifs}

\begin{prerequis}
\begin{itemize}
\item Opérations élémentaires dans ℕ (addition, soustraction, multiplication) vues au collège
\item Puissances d'un nombre entier : aⁿ, a⁰ = 1 (a ≠ 0)
\item Division posée et tables de multiplication
\item Notion de multiple et de diviseur d'un entier naturel
\item Développement et factorisation d'expressions littérales simples
\end{itemize}
\end{prerequis}

\newpage

\coursec{Division euclidienne dans \texorpdfstring{$\mathbb{N}$}{N}}

Au marché Mokolo de Yaoundé, maman Élise vend des beignets dorés qui font la fierté du quartier. Ce matin, elle en a préparé 247 et elle les emballe dans des sachets de 12. Deux questions se posent immédiatement : combien de sachets complets obtiendra-t-elle ? Et combien de beignets resteront-ils, invendables par sachet ? Plus loin, dans un local climatisé de Camtel, un technicien configure un routeur qui ne « comprend » que les chiffres 0 et 1 : comment le nombre 13 s'écrit-il pour la machine ?

Ces deux situations, en apparence très différentes, reposent sur la même idée profonde : \cle{regrouper les unités par paquets de taille fixe}. Pour maman Élise, le paquet contient 12 beignets ; pour le routeur, le paquet contient 2 unités. C'est exactement ce que formalisent la \cle{division euclidienne} et la \cle{numération en base $a$}, outils indispensables en informatique, en cryptographie et dans la vie courante. Dans cette première section, nous construisons rigoureusement la division euclidienne dans $\mathbb{N}$ : qu'est-ce qu'elle affirme ? Pourquoi le résultat est-il unique ? Comment déterminer le quotient et le reste sans se tromper ?

\courssub{Rappel : la division posée et son vocabulaire}

Dès le primaire, tu as appris à « poser » une division. Reprenons cette gestuelle avec le vocabulaire exact exigé au Probatoire.

\begin{definition}[Vocabulaire de la division]
Diviser un entier naturel $a$ (le \cle{dividende}) par un entier naturel non nul $b$ (le \cle{diviseur}), c'est chercher deux entiers naturels $q$ (le \cle{quotient}) et $r$ (le \cle{reste}) tels que :
\[ a = b \times q + r \quad \text{avec} \quad 0 \leq r < b. \]
\end{definition}

L'interprétation concrète est celle du \cle{partage équitable} : on dispose de $a$ objets que l'on regroupe en paquets de $b$. On forme $q$ paquets complets, et il reste $r$ objets isolés, forcément moins nombreux que $b$ — sinon on pourrait encore former un paquet de plus !

\begin{popfigure}
\begin{center}
\begin{tikzpicture}[scale=0.62]
% 7 paquets de 6
\foreach \p in {0,...,6}{
  \draw[rounded corners=2pt, thick, popblue, fill=popblueL] ({1.7*\p},0) rectangle ({1.7*\p+1.35},2.1);
  \foreach \i in {0,1,2}{
    \foreach \j in {0,1}{
      \filldraw[popblue] ({1.7*\p+0.4+0.55*\i},{0.55+1.0*\j}) circle (0.18);
    }
  }
}
% reste de 5
\draw[rounded corners=2pt, thick, poporange, dashed, fill=poporangeL] ({1.7*7},0) rectangle ({1.7*7+1.35},2.1);
\foreach \i in {0,1}{
  \foreach \j in {0,1}{
    \filldraw[poporange] ({1.7*7+0.4+0.55*\i},{0.55+1.0*\j}) circle (0.18);
  }
}
\filldraw[poporange] ({1.7*7+0.675},{1.05}) circle (0.18);
\node[popdark, font=\small] at (6.0,-0.7) {7 paquets complets de 6 objets};
\node[popdark, font=\small] at ({1.7*7+0.675},-0.7) {reste : 5};
\end{tikzpicture}
\end{center}
\legende{47 objets regroupés en paquets de 6 : on obtient 7 paquets complets et il reste 5 objets, car $47 = 6 \times 7 + 5$. Le reste 5 est bien strictement inférieur à 6.}
\end{popfigure}

\courssub{Le théorème fondamental : existence et unicité}

Ce que l'expérience du partage suggère, les mathématiques le garantissent : le quotient et le reste existent toujours, et ils sont \emph{uniques}. C'est ce double résultat qui fait toute la puissance de la division euclidienne.

\begin{propriete}[Théorème de la division euclidienne dans $\mathbb{N}$]
Soient $a \in \mathbb{N}$ et $b \in \mathbb{N}^{*}$. Il existe un \textbf{unique} couple $(q, r) \in \mathbb{N}^2$ tel que :
\[ a = bq + r \quad \text{avec} \quad 0 \leq r < b. \]
L'entier $q$ est appelé \cle{quotient}, l'entier $r$ est appelé \cle{reste} de la division euclidienne de $a$ par $b$.
\end{propriete}

L'\emph{existence} du couple $(q,r)$ est \textbf{admise} : elle repose sur une propriété profonde de $\mathbb{N}$ (dite archimédianité) : en ajoutant $b$ à lui-même suffisamment de fois, on finit toujours par dépasser $a$ ; le dernier multiple de $b$ qui ne dépasse pas $a$ fournit le quotient. En revanche, l'\emph{unicité} se démontre, et cette démonstration est formatrice : elle utilise le raisonnement par l'absurde.

\begin{experience}[Démonstration guidée : unicité du couple $(q,r)$]
Supposons, par l'absurde, qu'il existe \textbf{deux} couples $(q, r)$ et $(q', r')$ d'entiers naturels vérifiant tous les deux les conditions du théorème :
\[ a = bq + r \ \text{avec}\ 0 \leq r < b \qquad \text{et} \qquad a = bq' + r' \ \text{avec}\ 0 \leq r' < b. \]
\begin{enumerate}
\item \textbf{Étape 1 : égaliser.} Puisque les deux expressions valent $a$, on a $bq + r = bq' + r'$, donc :
\[ b(q - q') = r' - r. \]
\item \textbf{Étape 2 : encadrer la différence des restes.} De $0 \leq r < b$ et $0 \leq r' < b$, on tire $-b < r' - r < b$, c'est-à-dire :
\[ |r' - r| < b. \]
\item \textbf{Étape 3 : exploiter la divisibilité.} D'après l'étape 1, $r' - r$ est un multiple de $b$ (c'est $b$ fois l'entier relatif $q - q'$). Or le seul multiple de $b$ dont la valeur absolue est \emph{strictement} inférieure à $b$ est $0$. Donc $r' - r = 0$, c'est-à-dire $\boxed{r = r'}$.
\item \textbf{Étape 4 : conclure sur les quotients.} En reportant dans l'étape 1 : $b(q - q') = 0$, et comme $b \neq 0$, on obtient $q = q'$.
\end{enumerate}
Les deux couples sont donc identiques : le couple $(q, r)$ est bien \textbf{unique}. \hfill$\blacksquare$
\end{experience}

\begin{aretenir}[L'idée clé de la démonstration]
Deux restes d'une même division par $b$ diffèrent d'un multiple de $b$, tout en étant tous deux compris entre $0$ et $b - 1$. Le seul multiple de $b$ possible dans cet intervalle étroit est $0$ : les restes coïncident, puis les quotients aussi. Retiens l'enchaînement : \emph{égaliser} $\rightarrow$ \emph{encadrer} $\rightarrow$ \emph{conclure}.
\end{aretenir}

\courssub{Détermination pratique du quotient et du reste}

Sur le plan pratique, deux démarches équivalentes permettent de trouver $q$ et $r$ : la division posée (technique du collège) et l'\cle{encadrement} entre deux multiples consécutifs de $b$.

\begin{methode}[Déterminer $q$ et $r$ par encadrement]
Pour effectuer la division euclidienne de $a$ par $b$ ($b \neq 0$) :
\begin{enumerate}
\item Chercher le \textbf{plus grand multiple de $b$ inférieur ou égal à $a$} ; noter $q$ le facteur correspondant. On obtient l'encadrement :
\[ bq \leq a < b(q+1). \]
\item Calculer le reste par soustraction : $r = a - bq$.
\item \textbf{Vérifier} les deux conditions : $a = bq + r$ et $0 \leq r < b$.
\end{enumerate}
\end{methode}

L'encadrement $bq \leq a < b(q+1)$ se lit très bien sur une demi-droite graduée : $a$ se situe entre deux multiples consécutifs de $b$, et le quotient est exactement le nombre de « sauts » de longueur $b$ que l'on peut faire sans dépasser $a$.

\begin{popfigure}
\begin{center}
\begin{tikzpicture}[scale=0.85]
\draw[->, thick, popdark] (-0.3,0) -- (13.2,0);
% graduations : multiples de 6
\foreach \x/\n in {0/0, 1.5/6, 3/12, 4.5/18, 6/24, 7.5/30, 9/36, 10.5/42, 12/48}{
  \draw[thick, popdark] (\x,0.12) -- (\x,-0.12);
  \node[below, popdark, font=\small] at (\x,-0.15) {$\n$};
}
% multiples remarquables
\filldraw[popblue] (10.5,0) circle (0.09);
\node[above, popblue, font=\small] at (10.5,0.1) {$bq = 6 \times 7$};
\filldraw[popblue] (12,0) circle (0.09);
\node[above, popblue, font=\small] at (12,0.1) {$b(q+1) = 6 \times 8$};
% point a = 47
\filldraw[poporange] ({47/4},0) circle (0.11);
\node[above, poporange, font=\small] at ({47/4},0.12) {$a = 47$};
% accolade du reste
\draw[<->, thick, popgreen] (10.5,-0.9) -- ({47/4},-0.9);
\node[below, popgreen, font=\small] at ({(10.5+47/4)/2},-0.95) {reste $r = 5$};
\end{tikzpicture}
\end{center}
\legende{Encadrement de $a = 47$ entre deux multiples consécutifs de $b = 6$ : $6 \times 7 \leq 47 < 6 \times 8$. Le quotient est $q = 7$ et le reste est la distance $r = 47 - 42 = 5$.}
\end{popfigure}

\begin{exemplebox}[Division euclidienne de 247 par 12]
Revenons aux beignets de maman Élise : on divise $a = 247$ par $b = 12$.
\begin{itemize}
\item Multiples de 12 proches de 247 : $12 \times 20 = 240$ et $12 \times 21 = 252$. Or $240 \leq 247 < 252$, donc $q = 20$.
\item Reste : $r = 247 - 240 = 7$, et l'on a bien $0 \leq 7 < 12$.
\end{itemize}
Conclusion : $247 = 12 \times 20 + 7$. Maman Élise obtient \textbf{20 sachets complets} et il reste \textbf{7 beignets}.
\end{exemplebox}

\begin{attention}[La condition $0 \leq r < b$ n'est pas optionnelle !]
L'erreur la plus fréquente consiste à oublier la condition sur le reste. Par exemple, écrire $247 = 12 \times 19 + 19$ est une égalité \emph{exacte}, mais ce n'est \textbf{pas} la division euclidienne de 247 par 12, car $19 \geq 12$ : cela signifie simplement que la division n'est \emph{pas terminée} — on peut encore former un sachet de 12 ! Un reste supérieur ou égal au diviseur est le signal d'alarme : augmente le quotient. De même, un reste négatif est interdit. Au Probatoire, écris toujours la double condition : $a = bq + r$ \textbf{avec} $0 \leq r < b$.
\end{attention}

\courssub{Cas particuliers à connaître}

\begin{propriete}[Cas particuliers]
Soient $a \in \mathbb{N}$ et $b \in \mathbb{N}^{*}$.
\begin{itemize}
\item \textbf{Si $a < b$} : la division euclidienne de $a$ par $b$ donne $q = 0$ et $r = a$. En effet, $a = b \times 0 + a$ avec $0 \leq a < b$. (On ne peut former aucun paquet complet.)
\item \textbf{Si $r = 0$} : alors $a = bq$, et l'on dit que $b$ \cle{divise} $a$ (ou que $a$ est un \cle{multiple} de $b$, ou que $a$ est \cle{divisible} par $b$). On note $b \mid a$.
\end{itemize}
\end{propriete}

Ainsi, la divisibilité n'est rien d'autre qu'une division euclidienne dont le reste est nul : dire que 12 divise 240, c'est dire que $240 = 12 \times 20 + 0$. Ce lien simple entre reste nul et divisibilité sera exploité en Terminale pour l'étude des nombres premiers.

\begin{exoresolu}[Quotient et reste dans trois situations]
Déterminer le quotient et le reste de la division euclidienne : (a) de 47 par 6 ; (b) de 5 par 9 ; (c) de 132 par 11.
\tcblower
\textbf{(a) Division de 47 par 6.} Les multiples de 6 encadrant 47 sont $6 \times 7 = 42$ et $6 \times 8 = 48$. Comme $42 \leq 47 < 48$, on a $q = 7$ et $r = 47 - 42 = 5$. Vérification : $47 = 6 \times 7 + 5$ avec $0 \leq 5 < 6$. $\checkmark$

\textbf{(b) Division de 5 par 9.} Ici $a = 5 < b = 9$ : aucun paquet de 9 ne peut être formé. Donc $q = 0$ et $r = 5$. Vérification : $5 = 9 \times 0 + 5$ avec $0 \leq 5 < 9$. $\checkmark$

\textbf{(c) Division de 132 par 11.} On remarque que $11 \times 12 = 132$. Donc $q = 12$ et $r = 0$ : 11 divise 132. Vérification : $132 = 11 \times 12 + 0$ avec $0 \leq 0 < 11$. $\checkmark$
\end{exoresolu}

\begin{methode}[Rédaction type au Probatoire]
Pour toute question de division euclidienne, rédige en trois lignes :
\begin{enumerate}
\item L'encadrement : $bq \leq a < b(q+1)$ ;
\item L'égalité : $a = bq + r$ ;
\item La condition : $0 \leq r < b$.
\end{enumerate}
Cette rédaction complète garantit tous les points, même si le calcul est simple.
\end{methode}

\courssub{Complément : vers les congruences (pour la Terminale)}

La division euclidienne possède un prolongement puissant que tu rencontreras en Terminale : les \cle{congruences}. L'idée est que le reste $r$ caractérise complètement la « position » de $a$ par rapport aux multiples de $b$. On dit que deux entiers sont \emph{congrus modulo $b$} lorsqu'ils ont le même reste dans la division euclidienne par $b$, et l'on note $a \equiv r \ [b]$. Par exemple, $247 \equiv 7 \ [12]$ : sur un cadran d'horloge de 12 heures, 247 heures retombent sur la même position que 7 heures. Ce langage, fondé entièrement sur le théorème de cette section, est l'outil de base de l'arithmétique modulaire. \emph{(Notion signalée à titre de complément culturel : elle n'est pas exigible en Première.)}

\begin{savaistu}[De l'Antiquité à ta carte bancaire]
La division euclidienne porte le nom d'Euclide d'Alexandrie (III\textsuperscript{e} siècle av. J.-C.), qui l'utilisa pour construire l'\cle{algorithme d'Euclide} : en répétant des divisions euclidiennes (le diviseur devient dividende, le reste devient diviseur, jusqu'à obtenir un reste nul), on calcule le \textbf{PGCD} de deux entiers sans connaître leurs facteurs premiers. Plus de deux mille ans plus tard, cette même idée est au cœur des codes de sécurité bancaire et du chiffrement RSA : les algorithmes qui protègent tes transactions Mobile Money (MTN MoMo, Orange Money) reposent sur des divisions euclidiennes géantes, effectuées des millions de fois par seconde, avec des restes de plusieurs centaines de chiffres !
\end{savaistu}

\begin{exercice}[À toi de jouer]
\begin{enumerate}
\item Déterminer le quotient et le reste de la division euclidienne de 350 par 13, puis vérifier.
\item Un transporteur de la ligne Douala--Bafoussam charge 200 cartons dans des caisses de 16 cartons. Combien de caisses pleines obtient-il ? Reste-t-il des cartons ?
\item Trouver tous les entiers naturels $a$ tels que la division euclidienne de $a$ par 7 donne un quotient égal au reste.
\end{enumerate}
\end{exercice}

\begin{corrige}[Exercice : À toi de jouer]
\textbf{1.} $13 \times 26 = 338$ et $13 \times 27 = 351$. Comme $338 \leq 350 < 351$, on a $q = 26$ et $r = 350 - 338 = 12$. Vérification : $350 = 13 \times 26 + 12$ avec $0 \leq 12 < 13$. $\checkmark$

\textbf{2.} $16 \times 12 = 192$ et $16 \times 13 = 208 > 200$. Donc $200 = 16 \times 12 + 8$ : il obtient \textbf{12 caisses pleines} et il reste \textbf{8 cartons}.

\textbf{3.} On cherche $a = 7q + r$ avec $q = r$ et $0 \leq r < 7$. Alors $a = 7r + r = 8r$ avec $r \in \{0, 1, 2, 3, 4, 5, 6\}$. Les solutions sont : $a \in \{0, 8, 16, 24, 32, 40, 48\}$.
\end{corrige}

En résumé, la division euclidienne transforme l'idée intuitive du « regroupement par paquets » en un théorème précis : pour $a \in \mathbb{N}$ et $b \in \mathbb{N}^{*}$, l'écriture $a = bq + r$ avec $0 \leq r < b$ existe toujours et est unique. C'est sur cette fondation que nous construirons, dans la prochaine section, les systèmes de numération en base $a$ — et nous comprendrons enfin pourquoi le routeur de Camtel écrit 13 sous la forme 1101.

\coursec{Numération de base a : principe et décomposition}

Dans la section précédente, nous avons établi la division euclidienne dans $\mathbb{N}$ : pour tout entier naturel $a$ et tout entier naturel non nul $b$, il existe un unique couple $(q,r)$ tel que $a = bq + r$ avec $0 \le r < b$. Ce résultat, en apparence très simple, est en réalité la clé qui ouvre la porte de toute l'écriture des nombres. Pourquoi écrivons-nous $245$ pour désigner deux cent quarante-cinq ? Pourquoi les ordinateurs n'utilisent-ils que les chiffres $0$ et $1$ ? La réponse tient en deux mots : \cle{numération positionnelle}. Et c'est précisément en répétant des divisions euclidiennes que l'on construit l'écriture d'un nombre dans n'importe quelle base. C'est l'objet de cette section.

\courssub{Le principe de la numération de position}

\textbf{Une intuition pour commencer.}\par
Imagine un commerçant au marché de Mokolo qui regroupe ses billes de savon par paquets. Il fait des paquets de $10$ billes, puis des cartons de $10$ paquets (soit $100$ billes), puis des caisses de $10$ cartons (soit $1000$ billes). Pour décrire un stock de $245$ billes, il n'a besoin que de trois renseignements : $2$ cartons, $4$ paquets et $5$ billes isolées. Le nombre s'écrit alors $245$ : chaque chiffre indique \emph{combien} de paquets d'une certaine taille, et c'est la \emph{position} du chiffre qui indique la \emph{taille} du paquet.

C'est exactement le principe de notre système décimal : dans $245$, le chiffre $2$ ne vaut pas « deux », il vaut « deux cents » parce qu'il occupe la troisième position à partir de la droite. La valeur d'un chiffre dépend donc de deux choses : le chiffre lui-même, et la place qu'il occupe.

\begin{definition}[Système de numération positionnel de base $a$]
Soit $a$ un entier naturel supérieur ou égal à $2$ ($a \in \{2\,;\,3\,;\,\ldots\,;\,10\}$). Un \cle{système de numération positionnel de base $a$} est un système d'écriture des entiers naturels dans lequel :
\begin{itemize}
\item on utilise exactement $a$ symboles appelés \cle{chiffres} : $0, 1, 2, \ldots, a-1$ ;
\item la valeur représentée par un chiffre dépend de sa \cle{position} dans l'écriture : le chiffre situé en position $i$ (en comptant à partir de $0$ depuis la droite) est multiplié par $a^i$.
\end{itemize}
L'entier $a$ est appelé la \cle{base} du système.
\end{definition}

\begin{attention}[Les chiffres autorisés dépendent de la base]
En base $a$, les seuls chiffres qui \emph{existent} sont $0, 1, \ldots, a-1$. Il est donc \textbf{impossible} d'utiliser un chiffre supérieur ou égal à la base :
\begin{itemize}
\item en base $3$, le chiffre $3$ n'existe pas : une écriture comme $(312)_3$ est \textbf{fausse} ;
\item en base $2$, seuls $0$ et $1$ sont permis : $(1201)_2$ n'a aucun sens ;
\item en base $8$, les chiffres $8$ et $9$ sont interdits.
\end{itemize}
Avant tout calcul, vérifie toujours que chaque chiffre de l'écriture est strictement inférieur à la base. C'est la première chose qu'un correcteur du Probatoire regarde !
\end{attention}

\courssub{Décomposition d'un entier en somme de puissances de $a$}

Revenons à notre exemple : $245 = 2\times 100 + 4\times 10 + 5$, c'est-à-dire $245 = 2\times 10^2 + 4\times 10^1 + 5\times 10^0$. Toute l'écriture décimale repose sur cette décomposition en \cle{puissances de $10$}. Le même principe fonctionne avec n'importe quelle base $a$.

\begin{definition}[Décomposition en base $a$ et écriture d'un entier]
Soit $a \ge 2$ un entier naturel. Tout entier naturel non nul $N$ se décompose en une somme de puissances ordonnées de $a$ :
\[ N = x_n a^n + x_{n-1} a^{n-1} + \cdots + x_2 a^2 + x_1 a + x_0 \]
où les $x_i$ sont des entiers naturels vérifiant $0 \le x_i < a$ pour tout $i$, et $x_n \neq 0$.

On écrit alors $N$ en base $a$ sous la forme :
\[ N = (x_n\, x_{n-1}\, \ldots\, x_1\, x_0)_a \]
c'est-à-dire en juxtaposant les coefficients $x_i$ (appelés les chiffres de l'écriture), du coefficient de la plus grande puissance (à gauche) au coefficient de $a^0$ (à droite). La base est indiquée en indice, entre parenthèses ou en surlignant l'écriture avec l'indice $a$ : $\overline{x_n x_{n-1} \ldots x_1 x_0}^{\,a}$.
\end{definition}

\begin{attention}[Conditions à ne jamais oublier]
Deux conditions sont indispensables dans la définition :
\begin{itemize}
\item chaque coefficient vérifie $0 \le x_i < a$ (ce sont bien des \emph{chiffres} de la base) ;
\item le premier coefficient $x_n$ est \textbf{non nul} : on ne commence jamais une écriture par un $0$ à gauche. On écrit $(1011)_2$ et non $(01011)_2$.
\end{itemize}
Par convention, quand aucune base n'est indiquée, le nombre est écrit en base $10$ : $245$ signifie $(245)_{10}$.
\end{attention}

\courssub{Existence et unicité de la décomposition}

\begin{propriete}[Existence et unicité de l'écriture en base $a$ — admise]
Soit $a$ un entier naturel supérieur ou égal à $2$. Tout entier naturel non nul $N$ s'écrit de manière \textbf{unique} sous la forme
\[ N = x_n a^n + x_{n-1} a^{n-1} + \cdots + x_1 a + x_0 \quad \text{avec } 0 \le x_i < a \text{ et } x_n \neq 0. \]
Autrement dit : tout entier possède une écriture en base $a$, et une seule. Ce résultat est \textbf{admis} au niveau Première ; sa justification repose sur les divisions euclidiennes successives par $a$, que nous détaillons ci-dessous.
\end{propriete}

\begin{experience}[Pourquoi ça marche : les divisions euclidiennes successives]
Prenons $N = 245$ et $a = 10$, et effectuons des divisions euclidiennes répétées :
\begin{itemize}
\item $245 = 10 \times 24 + 5$ : le reste $5$ est le chiffre des unités $x_0$ ;
\item $24 = 10 \times 2 + 4$ : le reste $4$ est le chiffre des dizaines $x_1$ ;
\item $2 = 10 \times 0 + 2$ : le reste $2$ est le chiffre des centaines $x_2$.
\end{itemize}
On obtient $245 = 2\times 10^2 + 4 \times 10 + 5$. Les restes successifs donnent les chiffres, lus \emph{du dernier reste vers le premier} (de gauche à droite dans l'écriture).

Ce procédé fonctionne avec n'importe quelle base $a$ : à chaque étape, la division euclidienne par $a$ fournit un reste compris entre $0$ et $a-1$, c'est-à-dire un chiffre légal de la base. Comme le quotient diminue strictement à chaque étape, le procédé s'arrête. L'unicité de la division euclidienne (vue en section 1) garantit alors l'unicité de la décomposition : à chaque étape, le reste est imposé, donc chaque chiffre est imposé.
\end{experience}

\courssub{Lecture d'une écriture en base $a$ : le poids des positions}

Lire un nombre écrit en base $a$, c'est associer à chaque position une puissance de $a$, en partant de $a^0 = 1$ tout à droite, puis en multipliant par $a$ à chaque déplacement d'une case vers la gauche.

\begin{methode}[Lire un nombre écrit en base $a$]
Pour interpréter l'écriture $(x_n x_{n-1} \ldots x_1 x_0)_a$ :
\begin{enumerate}
\item numéroter les positions de droite à gauche, en commençant à $0$ ;
\item associer à la position $i$ la puissance $a^i$ (le \cle{poids} de la position) ;
\item multiplier chaque chiffre par le poids de sa position ;
\item additionner tous les produits obtenus : on obtient la valeur de $N$ (en base $10$).
\end{enumerate}
\end{methode}

\begin{popfigure}
\begin{center}
\begin{tikzpicture}[scale=1]
% Cases
\foreach \i/\c in {0/1, 1/0, 2/1, 3/1}{
  \draw[fill=popblueL, draw=popblue, thick] ({1.6*\i},0) rectangle ({1.6*\i+1.4},1.1);
  \node[font=\Large\bfseries, popdark] at ({1.6*\i+0.7},0.55) {\c};
}
% Poids
\foreach \i/\p in {0/$2^3$, 1/$2^2$, 2/$2^1$, 3/$2^0$}{
  \node[poporange, font=\bfseries] at ({1.6*\i+0.7},1.5) {\p};
}
% Valeurs
\foreach \i/\v in {0/8, 1/4, 2/2, 3/1}{
  \node[popmuted, font=\small] at ({1.6*\i+0.7},-0.45) {\v};
}
\node[popmuted, font=\small, anchor=east] at (-0.2,-0.45) {valeurs :};
\node[poporange, font=\small\bfseries, anchor=east] at (-0.2,1.5) {poids :};
\node[popdark, font=\small\bfseries, anchor=east] at (-0.2,0.55) {chiffres :};
% Accolade
\draw[decorate, decoration={brace, amplitude=5pt, mirror}, popdark, thick] (0,-0.85) -- (6.2,-0.85);
\node[popdark] at (3.1,-1.35) {$(1011)_2 = 1\times 8 + 0\times 4 + 1\times 2 + 1\times 1 = 11$};
\end{tikzpicture}
\end{center}
\legende{Schéma en « poids de position » : chaque case contient un chiffre, et le poids de la case est une puissance de la base. Ici, lecture de $(1011)_2$.}
\end{popfigure}

\begin{exemplebox}[Lecture de $(1011)_2$]
Décomposons le nombre binaire $(1011)_2$ en puissances de $2$ :
\begin{align*}
(1011)_2 &= 1\times 2^3 + 0\times 2^2 + 1\times 2^1 + 1\times 2^0 \\
&= 1\times 8 + 0\times 4 + 1\times 2 + 1\times 1 \\
&= 8 + 0 + 2 + 1 = 11.
\end{align*}
Donc $(1011)_2 = 11$. Remarque que le chiffre $0$ en deuxième position contribue pour $0$ à la somme, mais il est \emph{indispensable} : sans lui, $(111)_2 = 7$, qui est un tout autre nombre !
\end{exemplebox}

\courssub{Bases usuelles et tableau des puissances}

Certaines bases reviennent constamment, en mathématiques comme dans la vie courante :

\begin{itemize}
\item la \cle{base $10$} (système \emph{décimal}) : celle de notre vie quotidienne, héritée de nos dix doigts ; chiffres $0$ à $9$ ;
\item la \cle{base $2$} (système \emph{binaire}) : celle de l'informatique ; chiffres $0$ et $1$ ;
\item la \cle{base $8$} (système \emph{octal}) : utilisée en informatique comme raccourci du binaire ; chiffres $0$ à $7$.
\end{itemize}

Pour travailler vite et sans erreur, il faut connaître par cœur les premières puissances des bases usuelles. Les voici réunies :

\begin{center}
\rowcolors{1}{popcream}{white}
\begin{tabularx}{\linewidth}{|c|X|X|X|X|}
\hline
\rowcolor{popblueL} \textbf{Exposant $k$} & \textbf{$2^k$} & \textbf{$3^k$} & \textbf{$5^k$} & \textbf{$8^k$} \\
\hline
$0$ & $1$ & $1$ & $1$ & $1$ \\
\hline
$1$ & $2$ & $3$ & $5$ & $8$ \\
\hline
$2$ & $4$ & $9$ & $25$ & $64$ \\
\hline
$3$ & $8$ & $27$ & $125$ & $512$ \\
\hline
$4$ & $16$ & $81$ & $625$ & $4096$ \\
\hline
$5$ & $32$ & $243$ & $3125$ & \\
\hline
$6$ & $64$ & $729$ & & \\
\hline
$7$ & $128$ & & & \\
\hline
$8$ & $256$ & & & \\
\hline
\end{tabularx}
\end{center}

La base $2$ étant la plus utilisée en exercices, voici ses puissances en détail, à mémoriser absolument :

\begin{popfigure}
\begin{center}
\begin{tikzpicture}[scale=1]
\foreach \i/\p/\v in {0/{$2^0$}/1, 1/{$2^1$}/2, 2/{$2^2$}/4, 3/{$2^3$}/8, 4/{$2^4$}/16, 5/{$2^5$}/32, 6/{$2^6$}/64, 7/{$2^7$}/128, 8/{$2^8$}/256}{
  \draw[fill=popgreenL, draw=popgreen, thick] ({1.35*\i},0) rectangle ({1.35*\i+1.15},1.5);
  \node[popdark, font=\bfseries] at ({1.35*\i+0.575},1.1) {\p};
  \draw[popgreen] ({1.35*\i},0.75) -- ({1.35*\i+1.15},0.75);
  \node[popdark, font=\large] at ({1.35*\i+0.575},0.35) {\v};
}
\end{tikzpicture}
\end{center}
\legende{Les puissances de $2$ de $2^0$ à $2^8$ : chaque case vaut le double de la précédente.}
\end{popfigure}

\begin{aretenir}[L'essentiel de la section]
\begin{itemize}
\item En base $a$ ($a \ge 2$), on utilise les chiffres $0, 1, \ldots, a-1$, et la valeur d'un chiffre dépend de sa position.
\item Tout entier naturel non nul $N$ s'écrit de manière \textbf{unique} : $N = x_n a^n + \cdots + x_1 a + x_0$ avec $0 \le x_i < a$ et $x_n \neq 0$ ; on note $N = (x_n x_{n-1} \ldots x_1 x_0)_a$.
\item L'écriture s'obtient par des \textbf{divisions euclidiennes successives} par $a$ : les restes donnent les chiffres.
\item Pour lire un nombre en base $a$ : multiplier chaque chiffre par la puissance de $a$ de sa position (en partant de $a^0$ à droite), puis additionner.
\item Puissances de $2$ à connaître : $1, 2, 4, 8, 16, 32, 64, 128, 256$.
\end{itemize}
\end{aretenir}

\begin{exoresolu}[Lecture d'écritures en différentes bases]
\textbf{Énoncé.} Écrire en base $10$ chacun des nombres suivants : $(2102)_3$ ; $(345)_8$ ; $(11010)_2$.

\tcblower
\textbf{Solution détaillée.} On applique la méthode des poids de position : chaque chiffre est multiplié par la puissance de la base correspondant à sa position, en partant de l'exposant $0$ à droite.

\textbf{1.} Pour $(2102)_3$, les chiffres sont tous inférieurs à $3$, l'écriture est valide :
\begin{align*}
(2102)_3 &= 2\times 3^3 + 1\times 3^2 + 0\times 3^1 + 2\times 3^0 \\
&= 2\times 27 + 1\times 9 + 0\times 3 + 2\times 1 \\
&= 54 + 9 + 0 + 2 = 65.
\end{align*}
Donc $(2102)_3 = 65$.

\textbf{2.} Pour $(345)_8$, les chiffres $3, 4, 5$ sont tous inférieurs à $8$ :
\begin{align*}
(345)_8 &= 3\times 8^2 + 4\times 8^1 + 5\times 8^0 \\
&= 3\times 64 + 4\times 8 + 5\times 1 \\
&= 192 + 32 + 5 = 229.
\end{align*}
Donc $(345)_8 = 229$.

\textbf{3.} Pour $(11010)_2$ :
\begin{align*}
(11010)_2 &= 1\times 2^4 + 1\times 2^3 + 0\times 2^2 + 1\times 2^1 + 0\times 2^0 \\
&= 16 + 8 + 0 + 2 + 0 = 26.
\end{align*}
Donc $(11010)_2 = 26$.

\textbf{Vérification mentale utile :} en base $2$, un nombre qui se termine par $0$ est pair (comme $26$), un nombre qui se termine par $1$ est impair. C'est un bon réflexe de contrôle !
\end{exoresolu}

\begin{attention}[Pièges fréquents à la lecture d'une base]
\begin{itemize}
\item \textbf{Se tromper de sens dans les positions.} La position $0$ (poids $a^0 = 1$) est toujours le chiffre le plus \emph{à droite}. Numéroter depuis la gauche est l'erreur la plus répandue.
\item \textbf{Oublier que $a^0 = 1$.} Le dernier chiffre à droite est toujours multiplié par $1$, pas par $a$.
\item \textbf{Oublier les chiffres nuls.} Dans $(1001)_2$, les deux zéros centraux comptent pour les positions : $(1001)_2 = 2^3 + 2^0 = 9$, et non $2^1 + 2^0 = 3$.
\item \textbf{Confondre le chiffre et sa valeur.} Dans $(345)_8$, le chiffre $3$ vaut $3 \times 64 = 192$, pas $3$.
\end{itemize}
\end{attention}

\begin{savaistu}[Le monde tourne en base 2]
Ton téléphone portable, l'ordinateur du lycée, le réseau MTN ou Orange qui transporte tes messages : tous ne connaissent que la base $2$. Au cœur de ces machines, l'information est codée par des courants électriques qui n'ont que deux états stables : le courant passe (état $1$) ou ne passe pas (état $0$). Un chiffre binaire s'appelle un \emph{bit} (contraction de l'anglais \emph{binary digit}, « chiffre binaire »). Un groupe de $8$ bits forme un \emph{octet}, qui peut représenter $2^8 = 256$ valeurs différentes, de $(00000000)_2 = 0$ à $(11111111)_2 = 255$. Quand tu achètes un téléphone de « $128$ Go » à Douala, tu achètes en réalité des milliards de petites cases qui ne savent stocker que des $0$ et des $1$ ! Et ce n'est pas nouveau : le mathématicien Leibniz (fin du XVII\textsuperscript{e} siècle) avait déjà décrit le système binaire et pressenti son intérêt pour les machines à calculer, trois siècles avant l'ère du numérique.
\end{savaistu}

\begin{exercice}[Pour t'entraîner]
\begin{enumerate}
\item Les écritures suivantes sont-elles valides ? Justifier : $(2031)_4$ ; $(182)_8$ ; $(1021)_2$.
\item Écrire en base $10$ : $(1111)_2$ ; $(1201)_3$ ; $(507)_8$.
\item Quel est le plus grand nombre que l'on peut écrire avec exactement $4$ chiffres en base $2$ ? En base $3$ ?
\end{enumerate}
\end{exercice}

\begin{corrige}[Exercice : pour t'entraîner]
\textbf{1.} $(2031)_4$ : valide, car $0, 1, 2, 3 < 4$. $(182)_8$ : invalide, car le chiffre $8$ n'existe pas en base $8$. $(1021)_2$ : invalide, car le chiffre $2$ n'existe pas en base $2$.

\textbf{2.} $(1111)_2 = 8 + 4 + 2 + 1 = 15$. $(1201)_3 = 1\times 27 + 2\times 9 + 0\times 3 + 1 = 27 + 18 + 1 = 46$. $(507)_8 = 5\times 64 + 0\times 8 + 7 = 320 + 7 = 327$.

\textbf{3.} En base $2$ : $(1111)_2 = 15$. En base $3$ : $(2222)_3 = 2\times 27 + 2\times 9 + 2\times 3 + 2 = 54 + 18 + 6 + 2 = 80$. On remarque au passage que le plus grand nombre à $n$ chiffres en base $a$ vaut $a^n - 1$ : $2^4 - 1 = 15$ et $3^4 - 1 = 80$.
\end{corrige}

Maintenant que le principe de la décomposition est bien compris, la question suivante se pose naturellement : comment passer concrètement de la base $10$ à une base $a$, et inversement ? C'est l'objet de la section suivante, où les divisions euclidiennes successives deviendront ton outil de travail quotidien.

\coursec{Écriture d'un entier en base a : conversions avec la base 10}

Dans la section précédente, nous avons appris à décomposer un entier naturel $N$ en somme de puissances ordonnées d'une base $a$ :
\[ N = x_n a^n + x_{n-1} a^{n-1} + \cdots + x_1 a + x_0, \quad \text{avec } 0 \le x_i < a \text{ et } x_n \ne 0, \]
ce qui nous a permis d'écrire $N = \overline{x_n x_{n-1} \ldots x_1 x_0}^{\,a}$. Mais une question pratique demeure : comment passer \emph{concrètement} de l'écriture décimale d'un nombre à son écriture en base $a$, et inversement ? C'est toute l'affaire de cette section. Tu vas découvrir deux méthodes complémentaires : le \cle{développement en puissances} pour aller de la base $a$ vers la base $10$, et la méthode des \cle{divisions euclidiennes successives} pour aller de la base $10$ vers la base $a$. La seconde repose entièrement sur la division euclidienne étudiée en début de leçon : tu vois déjà que tout se tient !

\courssub{Conversion de la base a vers la base 10 : développer en puissances de a}

\textbf{Intuition.} L'écriture $\overline{x_n x_{n-1} \ldots x_1 x_0}^{\,a}$ n'est qu'une notation abrégée. Pour retrouver la valeur décimale du nombre, il suffit de « déplier » cette notation, c'est-à-dire d'écrire explicitement la somme des puissances de $a$ qu'elle cache, puis de calculer dans le système décimal que nous connaissons bien.

\begin{methode}[Convertir un nombre de la base $a$ vers la base $10$]
Pour convertir $N = \overline{x_n x_{n-1} \ldots x_1 x_0}^{\,a}$ en base $10$ :
\begin{enumerate}
\item Repérer la position de chaque chiffre : le chiffre le plus à droite a le rang $0$, le suivant le rang $1$, et ainsi de suite jusqu'au rang $n$ pour le chiffre le plus à gauche.
\item Écrire le développement $N = x_n \times a^n + x_{n-1} \times a^{n-1} + \cdots + x_1 \times a + x_0$.
\item Calculer les puissances de $a$, effectuer les produits, puis additionner : le résultat est l'écriture décimale de $N$.
\end{enumerate}
\end{methode}

\begin{exemplebox}[De la base $2$ et de la base $8$ vers la base $10$]
Convertissons $\overline{1001011}^{\,2}$ et $\overline{113}^{\,8}$ en base $10$.
\begin{align*}
\overline{1001011}^{\,2} &= 1\times 2^6 + 0\times 2^5 + 0\times 2^4 + 1\times 2^3 + 0\times 2^2 + 1\times 2 + 1 \\
&= 64 + 0 + 0 + 8 + 0 + 2 + 1 = 75. \\[4pt]
\overline{113}^{\,8} &= 1\times 8^2 + 1\times 8 + 3 = 64 + 8 + 3 = 75.
\end{align*}
Les deux écritures désignent donc le même entier : $75$.
\end{exemplebox}

\begin{attention}[Bien compter les rangs à partir de zéro]
L'erreur classique consiste à attribuer le rang $1$ au chiffre des unités. Non : le chiffre le plus à droite est toujours multiplié par $a^0 = 1$. Un bon réflexe : un nombre de $n+1$ chiffres en base $a$ fait intervenir les puissances de $a^0$ jusqu'à $a^n$. Ainsi $\overline{1001011}^{\,2}$ a $7$ chiffres, donc la plus grande puissance est $2^6$, pas $2^7$.
\end{attention}

\courssub{Conversion de la base 10 vers la base a : les divisions euclidiennes successives}

\textbf{Intuition.} Écrivons $75$ en base $2$. Dire que $75 = \overline{x_n \ldots x_1 x_0}^{\,2}$ signifie que $75 = 2 \times (\text{un certain entier}) + x_0$ avec $x_0 \in \{0 ; 1\}$. Autrement dit, le \emph{dernier chiffre} $x_0$ de l'écriture binaire n'est rien d'autre que le \emph{reste} de la division euclidienne de $75$ par $2$ ! En divisant ensuite le quotient obtenu par $2$, on obtient de la même façon l'avant-dernier chiffre, et ainsi de suite. Les chiffres de l'écriture en base $a$ apparaissent donc naturellement comme les restes successifs des divisions par $a$.

\begin{methode}[Convertir un entier de la base $10$ vers la base $a$]
Pour écrire un entier naturel non nul $N$ en base $a$ :
\begin{enumerate}
\item Effectuer la division euclidienne de $N$ par $a$ : on obtient un quotient $q_1$ et un reste $r_0$.
\item Effectuer la division euclidienne de $q_1$ par $a$ : on obtient $q_2$ et $r_1$. Recommencer avec $q_2$, et ainsi de suite.
\item S'arrêter dès que le quotient obtenu est nul.
\item Lire les restes \textbf{du dernier au premier} (de bas en haut dans la disposition en escalier) : ils forment, dans cet ordre, les chiffres de l'écriture de $N$ en base $a$.
\end{enumerate}
\end{methode}

\begin{experience}[Pourquoi ça marche ? Justification sur un nombre à trois chiffres]
Supposons qu'un entier $N$ s'écrive avec trois chiffres en base $a$ : $N = \overline{x_2 x_1 x_0}^{\,a}$, c'est-à-dire
\[ N = x_2 a^2 + x_1 a + x_0, \quad \text{avec } 0 \le x_i < a \text{ et } x_2 \ne 0. \]
\textbf{Étape 1 : première division.} Factorisons $a$ dans les deux premiers termes :
\[ N = a \underbrace{(x_2 a + x_1)}_{q_1} + \underbrace{x_0}_{r_0}. \]
Comme $0 \le x_0 < a$, c'est exactement la division euclidienne de $N$ par $a$ : le reste $r_0$ est le chiffre des unités $x_0$, et le quotient est $q_1 = x_2 a + x_1$.

\textbf{Étape 2 : deuxième division.} Divisons le quotient $q_1$ par $a$ :
\[ q_1 = a \underbrace{x_2}_{q_2} + \underbrace{x_1}_{r_1}. \]
Comme $0 \le x_1 < a$, le reste de cette division est le chiffre $x_1$, et le nouveau quotient est $q_2 = x_2$.

\textbf{Étape 3 : troisième division.} Divisons $q_2 = x_2$ par $a$. Puisque $0 < x_2 < a$, on obtient :
\[ q_2 = a \times \underbrace{0}_{q_3} + \underbrace{x_2}_{r_2}. \]
Le quotient est nul : on s'arrête. Le dernier reste est $x_2$.

\textbf{Conclusion.} Les restes successifs sont $x_0$, puis $x_1$, puis $x_2$ : en les lisant du dernier au premier, on retrouve $\overline{x_2 x_1 x_0}^{\,a} = N$. Le raisonnement se généralise à un nombre quelconque de chiffres : chaque division par $a$ « décolle » le chiffre des unités de ce qui reste à écrire. C'est pourquoi la méthode des divisions successives donne toujours l'écriture correcte.
\end{experience}

\begin{attention}[Le piège du sens de lecture]
C'est l'erreur la plus fréquente au Probatoire : lire les restes dans l'ordre où ils apparaissent. Retiens bien : le \textbf{premier reste obtenu est le chiffre des unités} (le plus à droite), et le \textbf{dernier reste est le chiffre de plus haut rang} (le plus à gauche). On lit donc les restes \textbf{de bas en haut} dans la disposition en escalier. Une vérification en reconvertissant vers la base $10$ permet de détecter immédiatement une inversion.
\end{attention}

\courssub{Exemple complet : écrire 75 en base 2, en base 3 et en base 8}

\textbf{Disposition en escalier.} On présente les divisions successives en cascade : chaque quotient devient le dividende de la division suivante, et les restes s'alignent à droite. Voici le cas de $75$ en base $2$ :

\begin{popfigure}
\begin{center}
\begin{tikzpicture}[scale=1, every node/.style={font=\normalsize}]
% Escalier des divisions de 75 par 2
\node (n1) at (0,0) {$75$};
\node (r1) at (1.6,0) {$\mathbf{1}$};
\node (n2) at (0,-0.7) {$37$};
\node (r2) at (1.6,-0.7) {$\mathbf{1}$};
\node (n3) at (0,-1.4) {$18$};
\node (r3) at (1.6,-1.4) {$\mathbf{0}$};
\node (n4) at (0,-2.1) {$9$};
\node (r4) at (1.6,-2.1) {$\mathbf{1}$};
\node (n5) at (0,-2.8) {$4$};
\node (r5) at (1.6,-2.8) {$\mathbf{0}$};
\node (n6) at (0,-3.5) {$2$};
\node (r6) at (1.6,-3.5) {$\mathbf{0}$};
\node (n7) at (0,-4.2) {$1$};
\node (r7) at (1.6,-4.2) {$\mathbf{1}$};
\node (n8) at (0,-4.9) {$0$};
% traits de division
\foreach \y in {0,-0.7,-1.4,-2.1,-2.8,-3.5,-4.2} {
  \draw[popline] (0.45,\y) -- (0.45,\y-0.7);
  \draw[popline] (0.45,\y) -- (1.15,\y);
}
% étiquettes
\node[popmuted, anchor=west] at (2.3,0) {restes};
\node[popmuted, anchor=east] at (-0.3,0) {quotients};
% flèche sens de lecture (CORRIGÉE : du bas vers le haut)
\draw[-{Stealth}, very thick, poporange] (2.6,-4.4) -- (2.6,-0.2) node[midway, right, poporange, align=center]{sens de\\lecture};
\end{tikzpicture}
\end{center}
\legende{Divisions euclidiennes successives de $75$ par $2$. Les restes (en gras) se lisent de bas en haut : $75 = \overline{1001011}^{\,2}$.}
\end{popfigure}

Détaillons chaque ligne : $75 = 2 \times 37 + 1$ ; $37 = 2 \times 18 + 1$ ; $18 = 2 \times 9 + 0$ ; $9 = 2 \times 4 + 1$ ; $4 = 2 \times 2 + 0$ ; $2 = 2 \times 1 + 0$ ; $1 = 2 \times 0 + 1$. Le quotient est nul : on s'arrête. En lisant les restes du dernier au premier : $1\,0\,0\,1\,0\,1\,1$, donc $75 = \overline{1001011}^{\,2}$.

L'arbre suivant visualise le cheminement complet, du nombre de départ jusqu'à la lecture finale :

\begin{popfigure}
\begin{center}
\begin{tikzpicture}[
  level distance=1.05cm,
  every node/.style={font=\small},
  level 1/.style={sibling distance=2.6cm},
  level 2/.style={sibling distance=2.6cm},
  edge from parent/.style={draw=popblue, -{Stealth}},
  qnode/.style={circle, draw=popblue, fill=popblueL, inner sep=2.5pt},
  rnode/.style={rectangle, rounded corners=2pt, draw=poporange, fill=poporangeL, inner sep=2.5pt}
]
\node[qnode]{$75$}
  child { node[qnode]{$37$}
    child { node[qnode]{$18$}
      child { node[qnode]{$9$}
        child { node[qnode]{$4$}
          child { node[qnode]{$2$}
            child { node[qnode]{$1$}
              child { node[qnode]{$0$}
                edge from parent node[rnode, right]{$1$} }
              edge from parent node[rnode, right]{$0$} }
            edge from parent node[rnode, right]{$0$} }
          edge from parent node[rnode, right]{$1$} }
        edge from parent node[rnode, right]{$0$} }
      edge from parent node[rnode, right]{$1$} }
    edge from parent node[rnode, right]{$1$} };
\end{tikzpicture}
\end{center}
\legende{Chaîne des quotients (cercles bleus) et restes successifs (étiquettes orange) de la division de $75$ par $2$. En remontant les restes du bas vers le haut : $\overline{1001011}^{\,2}$.}
\end{popfigure}

\begin{exemplebox}[75 en base 3 et en base 8]
\textbf{En base 3 :} $75 = 3 \times 22 + 0$ ; $22 = 3 \times 7 + 1$ ; $7 = 3 \times 2 + 1$ ; $2 = 3 \times 0 + 2$. Lecture des restes de bas en haut : $2\,2\,1\,0$, donc $75 = \overline{2210}^{\,3}$.

\textbf{En base 8 :} $75 = 8 \times 9 + 3$ ; $9 = 8 \times 1 + 1$ ; $1 = 8 \times 0 + 1$. Lecture des restes de bas en haut : $1\,1\,3$, donc $75 = \overline{113}^{\,8}$.
\end{exemplebox}

\begin{exoresolu}[Conversions et vérifications systématiques]
Écrire $75$ en base $2$, en base $3$ et en base $8$, puis vérifier chaque résultat par reconversion en base $10$.
\tcblower
\textbf{Conversions} (divisions successives, détaillées ci-dessus) :
\[ 75 = \overline{1001011}^{\,2} = \overline{2210}^{\,3} = \overline{113}^{\,8}. \]
\textbf{Vérifications} par développement en puissances :
\begin{align*}
\overline{1001011}^{\,2} &= 2^6 + 2^3 + 2 + 1 = 64 + 8 + 2 + 1 = 75. \checkmark \\
\overline{2210}^{\,3} &= 2\times 3^3 + 2\times 3^2 + 1\times 3 + 0 = 54 + 18 + 3 + 0 = 75. \checkmark \\
\overline{113}^{\,8} &= 8^2 + 8 + 3 = 64 + 8 + 3 = 75. \checkmark
\end{align*}
Les trois vérifications aboutissent à $75$ : les conversions sont correctes. Prends l'habitude de toujours effectuer cette vérification, surtout à l'examen : elle coûte une minute et peut sauver des points.
\end{exoresolu}

\begin{aretenir}[Les deux sens de conversion]
\begin{itemize}
\item \textbf{Base $a$ vers base $10$} : développer en puissances de $a$ et calculer :
\[ \overline{x_n x_{n-1} \ldots x_1 x_0}^{\,a} = x_n a^n + x_{n-1} a^{n-1} + \cdots + x_1 a + x_0. \]
\item \textbf{Base $10$ vers base $a$} : divisions euclidiennes successives par $a$ jusqu'à un quotient nul ; les restes, lus \textbf{du dernier au premier}, donnent l'écriture en base $a$.
\item \textbf{Vérification} : toujours reconvertir le résultat en base $10$.
\end{itemize}
\end{aretenir}

\begin{savaistu}[L'algorithme des ordinateurs]
La méthode des divisions successives n'est pas qu'un exercice de classe : c'est exactement l'algorithme qu'utilisent les ordinateurs et les logiciels pour afficher les nombres en binaire ! En machine, tout nombre est stocké sous forme de bits ($0$ ou $1$), et la conversion décimal-binaire s'effectue par cette cascade de divisions par $2$, car elle n'exige que des opérations simples et rapides. Au Cameroun, chaque fois qu'un technicien de MTN ou d'Orange configure une adresse IP (par exemple $192.168.1.1$), les nombres manipulés sont convertis en binaire par ce procédé : $192 = \overline{11000000}^{\,2}$. Tu pratiques donc, sans le savoir, un algorithme au cœur de toutes les télécommunications modernes.
\end{savaistu}

\begin{exercice}[À toi de jouer]
\begin{enumerate}
\item Écrire $123$ en base $2$, en base $5$ et en base $8$, puis vérifier chaque résultat.
\item Convertir $\overline{101101}^{\,2}$, $\overline{2102}^{\,3}$ et $\overline{456}^{\,7}$ en base $10$.
\end{enumerate}
\end{exercice}

\begin{corrige}[Exercice]
\textbf{1.} Base $2$ : $123 = 2\times 61 + 1$ ; $61 = 2\times 30 + 1$ ; $30 = 2\times 15 + 0$ ; $15 = 2\times 7 + 1$ ; $7 = 2\times 3 + 1$ ; $3 = 2\times 1 + 1$ ; $1 = 2\times 0 + 1$. D'où $123 = \overline{1111011}^{\,2}$. Vérification : $64 + 32 + 16 + 8 + 0 + 2 + 1 = 123$. \checkmark

Base $5$ : $123 = 5\times 24 + 3$ ; $24 = 5\times 4 + 4$ ; $4 = 5\times 0 + 4$. D'où $123 = \overline{443}^{\,5}$. Vérification : $4\times 25 + 4\times 5 + 3 = 100 + 20 + 3 = 123$. \checkmark

Base $8$ : $123 = 8\times 15 + 3$ ; $15 = 8\times 1 + 7$ ; $1 = 8\times 0 + 1$. D'où $123 = \overline{173}^{\,8}$. Vérification : $64 + 7\times 8 + 3 = 64 + 56 + 3 = 123$. \checkmark

\textbf{2.} $\overline{101101}^{\,2} = 32 + 8 + 4 + 1 = 45$ ; $\overline{2102}^{\,3} = 2\times 27 + 1\times 9 + 0\times 3 + 2 = 54 + 9 + 2 = 65$ ; $\overline{456}^{\,7} = 4\times 49 + 5\times 7 + 6 = 196 + 35 + 6 = 237$.
\end{corrige}

Tu maîtrises désormais les allers-retours entre la base $10$ et n'importe quelle base $a$. Dans la section suivante, nous irons plus loin : nous apprendrons à passer \emph{directement} d'une base à une autre, sans repasser par la base $10$ lorsque ce n'est pas nécessaire.

\coursec{Passage direct d'une base à l'autre}

Dans la section précédente, tu as appris à convertir un entier de la base $10$ vers une base $a$ quelconque (par divisions successives) et, inversement, à retrouver l'écriture décimale d'un nombre donné en base $a$ (par développement en puissances de $a$). Une question naturelle se pose alors : comment passer \emph{directement} d'une base $a$ à une base $b$, par exemple de la base $2$ à la base $8$, sans revenir à la base $10$ ? C'est l'objet de cette section. Nous verrons d'abord la méthode générale, qui fonctionne toujours, puis un cas particulier remarquable et très rapide : les conversions entre les bases $2$, $4$ et $8$, fondées sur un simple \cle{regroupement de chiffres binaires}.

\courssub{Méthode générale : passage par la base 10}

Soient $a$ et $b$ deux entiers naturels supérieurs ou égaux à $2$, et $N$ un entier naturel dont on connaît l'écriture en base $a$. Pour obtenir son écriture en base $b$, on ne connaît a priori qu'un seul « pont » fiable entre les deux mondes : la base $10$, dans laquelle nous savons calculer.

\begin{methode}[Conversion d'une base $a$ vers une base $b$ quelconques]
Pour écrire en base $b$ un entier $N$ donné en base $a$ :
\begin{enumerate}
\item \textbf{Étape 1 :} convertir $N$ de la base $a$ vers la base $10$ en développant en puissances de $a$ :
\[ N = x_n a^n + x_{n-1}a^{n-1} + \cdots + x_1 a + x_0. \]
\item \textbf{Étape 2 :} convertir le résultat obtenu de la base $10$ vers la base $b$ par divisions euclidiennes successives par $b$, puis lire les restes \emph{du dernier au premier}.
\end{enumerate}
\end{methode}

\begin{exemplebox}[De la base $5$ à la base $7$]
Convertissons $N = (342)_5$ en base $7$.

\textbf{Étape 1 : base $5 \to$ base $10$.}
\[ N = 3\times 5^2 + 4\times 5 + 2 = 3\times 25 + 20 + 2 = 97. \]

\textbf{Étape 2 : base $10 \to$ base $7$.} On effectue les divisions successives par $7$ :
\[ 97 = 7\times 13 + 6 \qquad ; \qquad 13 = 7\times 1 + 6 \qquad ; \qquad 1 = 7\times 0 + 1. \]
En lisant les restes du dernier au premier : $97 = (166)_7$.

\textbf{Conclusion :} $(342)_5 = (166)_7$.
\end{exemplebox}

Cette méthode est universelle, mais elle demande deux conversions successives. Existe-t-il des situations où l'on peut faire mieux, sans aucun calcul ? Oui : lorsque les deux bases sont des puissances d'un même nombre. C'est le cas des bases $2$, $4$ et $8$, omniprésentes en informatique.

\courssub{Le cas particulier fondamental : bases 2, 4 et 8}

\textbf{Intuition.} Observons les égalités suivantes :
\[ 4 = 2^2 \qquad \text{et} \qquad 8 = 2^3. \]
Ces relations signifient qu'\emph{un} chiffre en base $4$ « pèse » exactement autant qu'un bloc de \emph{deux} chiffres binaires, et qu'\emph{un} chiffre en base $8$ pèse exactement autant qu'un bloc de \emph{trois} chiffres binaires. Autrement dit, la base $8$ est une écriture « compressée » du binaire : trois bits se regroupent en un seul chiffre octal.

\begin{propriete}[Regroupement des chiffres binaires]
Tout entier naturel $N$ s'écrivant en base $2$ peut s'écrire en base $8$ en regroupant ses chiffres binaires (bits) par \cle{blocs de 3} en partant de la droite, puis en remplaçant chaque bloc par le chiffre octal correspondant. De même, on obtient l'écriture en base $4$ en regroupant les bits par \cle{blocs de 2}.
\end{propriete}

\begin{experience}[Pourquoi le regroupement par blocs de 3 fonctionne-t-il ?]
Écrivons le développement binaire d'un entier $N$ et regroupons les termes trois par trois. Par exemple, avec $N = (x_5 x_4 x_3 x_2 x_1 x_0)_2$ :
\begin{align*}
N &= x_5 2^5 + x_4 2^4 + x_3 2^3 + x_2 2^2 + x_1 2^1 + x_0 2^0 \\
&= \underbrace{\left(x_5 2^2 + x_4 2^1 + x_3\right)}_{\text{un chiffre entre } 0 \text{ et } 7}\times 2^3 + \underbrace{\left(x_2 2^2 + x_1 2^1 + x_0\right)}_{\text{un chiffre entre } 0 \text{ et } 7} \\
&= y_1 \times 8^1 + y_0 \times 8^0,
\end{align*}
où $y_1 = x_5 2^2 + x_4 2 + x_3$ et $y_0 = x_2 2^2 + x_1 2 + x_0$ sont bien des entiers compris entre $0$ et $7$, donc des \emph{chiffres} licites en base $8$. La clé est que $2^3 = 8$ : chaque bloc de trois bits se factorise exactement en une puissance de $8$. Le même raisonnement avec $2^2 = 4$ justifie les blocs de deux bits pour la base $4$.
\end{experience}

\begin{aretenir}[Table de correspondance chiffres octaux et blocs de 3 bits]
\begin{center}
\begin{tabular}{|c|c|c|c|c|c|c|c|c|}
\hline
\rowcolor{popblueL} Chiffre octal & 0 & 1 & 2 & 3 & 4 & 5 & 6 & 7 \\
\hline
Bloc binaire & 000 & 001 & 010 & 011 & 100 & 101 & 110 & 111 \\
\hline
\end{tabular}
\end{center}
Pour la base $4$, on utilise les blocs de 2 bits : $0 = (00)_2$, $1 = (01)_2$, $2 = (10)_2$, $3 = (11)_2$.
\end{aretenir}

\begin{popfigure}
\begin{center}
\begin{tikzpicture}[scale=0.9, every node/.style={font=\small}]
% Table de correspondance octal <-> binaire
\foreach \i/\b in {0/000, 1/001, 2/010, 3/011, 4/100, 5/101, 6/110, 7/111}{
  \node[draw=popblue, fill=popblueL, minimum width=1.05cm, minimum height=0.75cm, thick] at (\i*1.15, 0.6) {\textbf{\i}};
  \node[draw=popdark, fill=popcream, minimum width=1.05cm, minimum height=0.75cm] at (\i*1.15, -0.25) {\texttt{\b}};
}
\node[font=\small\itshape, popmuted] at (3.9, 1.35) {Chiffre octal};
\node[font=\small\itshape, popmuted] at (3.9, -1.0) {Bloc de 3 bits};
\end{tikzpicture}
\end{center}
\legende{Table de correspondance entre les chiffres de la base $8$ et les blocs de trois bits.}
\end{popfigure}

\courssub{Du binaire vers l'octale : regroupement par blocs de 3}

\begin{methode}[Conversion binaire $\to$ octale]
\begin{enumerate}
\item Partir de l'écriture binaire de $N$.
\item Regrouper les bits par \cle{blocs de 3} \textbf{en partant de la droite}.
\item Si le bloc le plus à gauche est incomplet, le compléter par des zéros à gauche (cela ne change pas la valeur du nombre).
\item Remplacer chaque bloc par le chiffre octal correspondant, à l'aide de la table de correspondance.
\end{enumerate}
\end{methode}

\begin{exemplebox}[Conversion de $(101101)_2$ en base $8$]
On regroupe les bits par blocs de 3 en partant de la droite :
\[ (101101)_2 = \underbrace{101}_{5}\ \underbrace{101}_{5} \]
Le bloc de gauche $101$ vaut $1\times 4 + 0\times 2 + 1 = 5$ ; le bloc de droite $101$ vaut aussi $5$. Donc :
\[ (101101)_2 = (55)_8. \]
\textbf{Vérification par la base 10 :} $(101101)_2 = 32 + 8 + 4 + 1 = 45$ et $(55)_8 = 5\times 8 + 5 = 45$. C'est cohérent.
\end{exemplebox}

\begin{popfigure}
\begin{center}
\begin{tikzpicture}[scale=1, every node/.style={font=\small}]
% Bits
\foreach \i/\bit in {0/1, 1/0, 2/1, 3/1, 4/0, 5/1}{
  \node[draw=popdark, fill=popcream, minimum width=0.8cm, minimum height=0.8cm, thick] at (\i*0.9, 0) {\textbf{\bit}};
}
% Accolades de regroupement (blocs de 3)
\draw[decorate, decoration={brace, amplitude=6pt, mirror}, popblue, thick] (-0.45, -0.5) -- (2.25, -0.5);
\draw[decorate, decoration={brace, amplitude=6pt, mirror}, poporange, thick] (2.35, -0.5) -- (5.05, -0.5);
\node[popblue, font=\small] at (0.9, -1.15) {bloc $101 = 5$};
\node[poporange, font=\small] at (3.7, -1.15) {bloc $101 = 5$};
% Flèches vers chiffres octaux
\draw[-{Stealth}, popblue, thick] (0.9, -1.45) -- (0.9, -2.1);
\draw[-{Stealth}, poporange, thick] (3.7, -1.45) -- (3.7, -2.1);
\node[draw=popblue, fill=popblueL, minimum width=0.8cm, minimum height=0.8cm, thick] at (0.9, -2.55) {\textbf{5}};
\node[draw=poporange, fill=poporangeL, minimum width=0.8cm, minimum height=0.8cm, thick] at (3.7, -2.55) {\textbf{5}};
\node[font=\small\itshape, popmuted] at (2.3, -3.3) {$(101101)_2 = (55)_8$};
\end{tikzpicture}
\end{center}
\legende{Regroupement de $(101101)_2$ en blocs de trois bits, en partant de la droite : chaque bloc donne un chiffre octal.}
\end{popfigure}

\begin{attention}[On regroupe toujours en partant de la droite !]
L'erreur la plus fréquente est de regrouper les bits \emph{en partant de la gauche}. C'est faux, car les puissances de $2$ croissent de la droite vers la gauche : c'est le chiffre des unités (tout à droite) qui porte le poids $2^0$. Par exemple, pour $(11011)_2$ :
\begin{itemize}
\item \textbf{Correct} (depuis la droite) : $(11011)_2 = \underbrace{011}_{3}\ \underbrace{011}_{3} = (33)_8$ après complétion par un zéro à gauche.
\item \textbf{Faux} (depuis la gauche) : on obtiendrait les blocs $110$ et $11$, mal alignés sur les puissances de $8$.
\end{itemize}
Vérification : $(11011)_2 = 16 + 8 + 2 + 1 = 27$ et $(33)_8 = 3\times 8 + 3 = 27$. Le regroupement depuis la droite est le bon.
\end{attention}

\courssub{De l'octale vers le binaire : éclatement de chaque chiffre}

\begin{methode}[Conversion octale $\to$ binaire]
Remplacer \textbf{chaque chiffre} de l'écriture octale par son écriture binaire sur \cle{exactement 3 bits} (en complétant par des zéros à gauche si besoin), puis juxtaposer les blocs. On peut ensuite supprimer les éventuels zéros inutiles tout à gauche du résultat.
\end{methode}

\begin{exemplebox}[Conversion de $(231)_8$ en base $2$]
On éclate chaque chiffre octal en un bloc de 3 bits :
\[ 2 = (010)_2 \qquad ; \qquad 3 = (011)_2 \qquad ; \qquad 1 = (001)_2. \]
En juxtaposant : $(231)_8 = (010\,011\,001)_2 = (10011001)_2$ après suppression du zéro initial.

\textbf{Vérification par la base 10 :} $(231)_8 = 2\times 64 + 3\times 8 + 1 = 128 + 24 + 1 = 153$ et $(10011001)_2 = 128 + 16 + 8 + 1 = 153$. Parfait.
\end{exemplebox}

\courssub{Binaire et base 4 : blocs de 2 bits}

Le même principe s'applique entre les bases $2$ et $4$, puisque $4 = 2^2$ : on regroupe les bits par \cle{blocs de 2} en partant de la droite.

\begin{exemplebox}[Conversion de $(110101)_2$ en base $4$]
Regroupement par blocs de 2 depuis la droite :
\[ (110101)_2 = \underbrace{11}_{3}\ \underbrace{01}_{1}\ \underbrace{01}_{1} = (311)_4. \]
Vérification : $(110101)_2 = 32 + 16 + 4 + 1 = 53$ et $(311)_4 = 3\times 16 + 1\times 4 + 1 = 53$.
\end{exemplebox}

\begin{exoresolu}[Conversion complète avec vérification]
Énoncé. On donne $N = (1100101)_2$.
\begin{enumerate}
\item Convertir $N$ en base $8$ par regroupement de bits.
\item Vérifier le résultat en passant par la base $10$.
\end{enumerate}
\tcblower
\textbf{1. Conversion en base $8$.} Le nombre comporte $7$ bits ; on regroupe par blocs de 3 en partant de la droite et on complète le bloc de gauche par des zéros :
\[ N = (1\,100\,101)_2 = (001\,100\,101)_2. \]
Or $001 = 1$, $100 = 4$ et $101 = 5$. Donc :
\[ N = (145)_8. \]

\textbf{2. Vérification par la base $10$.}
\begin{align*}
(1100101)_2 &= 2^6 + 2^5 + 2^2 + 2^0 = 64 + 32 + 4 + 1 = 101, \\
(145)_8 &= 1\times 8^2 + 4\times 8 + 5 = 64 + 32 + 5 = 101.
\end{align*}
Les deux écritures désignent le même entier $101$ : la conversion est correcte.
\end{exoresolu}

\begin{attention}[Zéros de complétion : à gauche seulement]
Compléter un bloc par des zéros ne se fait \textbf{jamais qu'à gauche} de l'écriture binaire (ou à gauche d'un bloc), car un zéro placé à gauche d'un entier ne change pas sa valeur : $(101)_2 = (0101)_2$. Ajouter un zéro à droite reviendrait à multiplier le nombre par $2$ !
\end{attention}

\begin{exercice}[À toi de jouer]
\begin{enumerate}
\item Convertir $(111010011)_2$ en base $8$, puis en base $4$.
\item Convertir $(672)_8$ en base $2$.
\item Convertir $(312)_4$ en base $2$, puis vérifier par la base $10$.
\end{enumerate}
\end{exercice}

\begin{corrige}[Exercice : conversions entre bases 2, 4 et 8]
\begin{enumerate}
\item Regroupement par blocs de 3 depuis la droite : $(111\,010\,011)_2 = (723)_8$. Par blocs de 2 : $(0111\,0100\,11)_2$ se réécrit $(01\,11\,01\,00\,11)_2 = (13103)_4$.
\item $6 = (110)_2$, $7 = (111)_2$, $2 = (010)_2$, donc $(672)_8 = (110111010)_2$.
\item $3 = (11)_2$, $1 = (01)_2$, $2 = (10)_2$, donc $(312)_4 = (110110)_2$. Vérification : $(312)_4 = 3\times 16 + 4 + 2 = 54$ et $(110110)_2 = 32 + 16 + 4 + 2 = 54$.
\end{enumerate}
\end{corrige}

\begin{savaistu}[La base 16 : le langage des informaticiens]
Les informaticiens utilisent aussi la \cle{base 16}, dite \emph{hexadécimale}. Comme il faut seize chiffres, on emploie les dix chiffres usuels puis les lettres : $A = 10$, $B = 11$, $C = 12$, $D = 13$, $E = 14$, $F = 15$. Puisque $16 = 2^4$, le passage binaire $\leftrightarrow$ hexadécimal se fait par blocs de \textbf{quatre} bits, exactement comme nous l'avons fait pour l'octale. Par exemple, $(11110010)_2 = (F2)_{16}$ car $1111 = 15 = F$ et $0010 = 2$. C'est cette écriture compacte que tu rencontres dans les codes couleurs des pages web (par exemple \texttt{FF8000} pour l'orange) ou dans les adresses mémoire des ordinateurs. Tu retrouveras cette base en Terminale et dans toutes les études d'informatique : la technique des blocs de bits que tu viens de maîtriser est déjà la bonne !
\end{savaistu}

\begin{aretenir}[L'essentiel de la section]
\begin{itemize}
\item Entre deux bases quelconques $a$ et $b$, on passe par la \cle{base 10} : développement en puissances de $a$, puis divisions successives par $b$.
\item Entre les bases $2$, $4$ et $8$, on utilise le \cle{regroupement de bits} : blocs de $2$ pour la base $4$ (car $4 = 2^2$), blocs de $3$ pour la base $8$ (car $8 = 2^3$).
\item Le regroupement se fait \textbf{toujours en partant de la droite}, avec complétion par des zéros à gauche si nécessaire.
\item Dans le sens octale $\to$ binaire, chaque chiffre octal est remplacé par son écriture sur exactement $3$ bits.
\end{itemize}
\end{aretenir}

\coursec{Addition en base $a$}

Dans la section précédente, tu as appris à passer d'une base à l'autre : tu sais désormais transformer $(1011)_2$ en $11_{10}$, ou encore $(456)_8$ en son équivalent décimal. Une question naturelle se pose alors : est-il possible de calculer \emph{directement} dans une base $a$, sans repasser par la base 10 ? La réponse est oui, et c'est même ainsi que fonctionnent les machines : un ordinateur ne connaît que le binaire, et pourtant il additionne, multiplie, compare à une vitesse vertigineuse. Dans cette section, tu vas découvrir que l'addition en base $a$ obéit exactement aux mêmes règles que celle que tu pratiques depuis le primaire en base 10, à une différence près, minuscule en apparence mais fondamentale : \textbf{la retenue ne vaut plus 10, elle vaut $a$}.

\courssub{Le principe : additionner comme en base 10, avec une retenue qui vaut $a$}

\textbf{Intuition.} Réfléchis à ce que tu fais quand tu poses $47 + 38$ en base 10. Tu additionnes les unités : $7 + 8 = 15$. Comme $15$ dépasse $10$, tu écris $5$ et tu « retiens 1 ». Pourquoi ? Parce que $15 = 1 \times 10 + 5$ : dix unités forment une dizaine, que tu vas reporter dans la colonne des dizaines. La retenue traduit exactement la division euclidienne de la somme de la colonne par la base : $15 = 10 \times 1 + 5$, on pose le reste $5$, on retient le quotient $1$.

En base $a$, le mécanisme est identique : $a$ unités d'un rang forment une unité du rang supérieur. En base 8, huit unités forment une « huitaine » ; en base 2, deux unités forment une « paire ». La retenue apparaît donc dès que la somme d'une colonne \emph{atteint ou dépasse $a$}, et elle se calcule par division euclidienne par $a$.

\begin{propriete}[Règle de la retenue en base $a$]
Pour additionner deux entiers écrits en base $a$, on additionne colonne par colonne, de droite à gauche. Si $S$ désigne la somme des chiffres d'une colonne (retenue éventuelle incluse), on effectue la division euclidienne de $S$ par $a$ :
\[ S = a \times q + r \quad \text{avec} \quad 0 \le r < a. \]
On \textbf{pose} le reste $r$ dans la colonne courante et on \textbf{retient} le quotient $q$ pour la colonne suivante. En pratique, comme la somme de deux chiffres plus une retenue ne dépasse pas $2a - 1$, le quotient $q$ vaut toujours $0$ ou $1$ : la retenue est donc toujours $1$ dès que $S \ge a$.
\end{propriete}

\begin{experience}[Pourquoi la retenue vaut-elle $1$ ?]
Justifions proprement ce fait. Les chiffres en base $a$ sont compris entre $0$ et $a-1$. La somme de deux chiffres d'une colonne, augmentée d'une retenue de $1$ venant de la colonne précédente, vaut au plus :
\[ (a-1) + (a-1) + 1 = 2a - 1. \]
Or la division euclidienne de $2a - 1$ par $a$ donne $2a - 1 = a \times 1 + (a - 1)$, donc le quotient est au plus $1$. Conclusion : en additionnant \emph{deux} nombres, la retenue ne peut valoir que $0$ ou $1$, et elle vaut $1$ exactement lorsque la somme de la colonne atteint ou dépasse $a$. C'est la généralisation directe de ce que tu fais en base 10 : $9 + 9 + 1 = 19 = 10 \times 1 + 9$.
\end{experience}

\begin{attention}[L'erreur classique : retenir 10 au lieu de $a$]
L'erreur la plus fréquente au Probatoire est de continuer à « penser en base 10 ». Par exemple, en base 8, beaucoup d'élèves écrivent $5 + 4 = 9$. C'est faux \emph{dans l'écriture en base 8} : le chiffre $9$ n'existe pas en base 8 ! On a $5 + 4 = 9_{10} = 8 + 1 = 1 \times 8 + 1$, donc en base 8 :
\[ 5 + 4 = (11)_8 \qquad \text{on pose } 1, \text{ on retient } 1. \]
De même, en base 2 : $1 + 1 = (10)_2$ (on pose $0$, on retient $1$), et $1 + 1 + 1 = (11)_2$ (on pose $1$, on retient $1$). Avant chaque colonne, demande-toi : « ma somme atteint-elle $a$ ? » Si oui, retranche $a$ et retiens $1$.
\end{attention}

\courssub{Tables d'addition en base 2 et en base 8}

Pour calculer vite et sans erreur, il est utile de connaître les petites tables d'addition. En base 2, il n'y a que deux chiffres, $0$ et $1$ : la table tient en quatre lignes.

\begin{aretenir}[Table d'addition en base 2]
\[ 0 + 0 = 0 \qquad 0 + 1 = 1 \qquad 1 + 0 = 1 \qquad 1 + 1 = 10 \]
Autrement dit : $1 + 1$ s'écrit $(10)_2$ : on pose $0$ et on retient $1$. Et lorsqu'une retenue s'ajoute : $1 + 1 + 1 = (11)_2$, on pose $1$ et on retient $1$.
\end{aretenir}

En base 8, les chiffres vont de $0$ à $7$. Voici les résultats qui font apparaître une retenue (ceux où la somme atteint $8$) :

\begin{center}
\begin{tabularx}{\linewidth}{|l|X|X|X|X|}
\hline
\rowcolor{popblueL} Somme & $4+4$ & $5+4$ & $6+6$ & $7+7$ \\
\hline
En base 10 & $8$ & $9$ & $12$ & $14$ \\
\hline
En base 8 & $10$ & $11$ & $14$ & $16$ \\
\hline
On pose / on retient & pose $0$, retient $1$ & pose $1$, retient $1$ & pose $4$, retient $1$ & pose $6$, retient $1$ \\
\hline
\end{tabularx}
\end{center}

La règle est toujours la même : dès que la somme atteint $8$, on retranche $8$ pour obtenir le chiffre posé, et on retient $1$.

\courssub{Technique opératoire : l'addition posée en colonnes}

\begin{methode}[Additionner deux nombres en base $a$]
\begin{enumerate}
\item \textbf{Poser} les deux nombres en colonnes, alignés à droite (chiffres des unités face à face), comme en base 10.
\item \textbf{Additionner} la colonne la plus à droite, en comptant mentalement en base 10 si cela t'aide.
\item \textbf{Comparer à $a$} : si la somme $S$ est strictement inférieure à $a$, on pose $S$ et il n'y a pas de retenue ; si $S \ge a$, on pose $S - a$ et on retient $1$.
\item \textbf{Recommencer} colonne par colonne vers la gauche, sans oublier d'ajouter la retenue éventuelle.
\item Si une retenue subsiste après la dernière colonne, \textbf{l'écrire} comme nouveau chiffre à gauche du résultat.
\item \textbf{Vérifier} en convertissant les deux nombres et le résultat en base 10.
\end{enumerate}
\end{methode}

\begin{exemplebox}[Addition en base 2 : $(1011)_2 + (110)_2$]
Posons l'opération en alignant à droite, et traitons les colonnes de droite à gauche :
\begin{itemize}
\item Colonne des unités : $1 + 0 = 1$. On pose $1$, retenue $0$.
\item Colonne des « deuxaines » : $1 + 1 = 2 = (10)_2$. On pose $0$, on retient $1$.
\item Colonne des « quatraines » : $0 + 1 + 1$ (retenue) $= 2 = (10)_2$. On pose $0$, on retient $1$.
\item Colonne des « huitaines » : $1 + 0 + 1$ (retenue) $= 2 = (10)_2$. On pose $0$, on retient $1$.
\item La retenue finale devient un nouveau chiffre : on pose $1$.
\end{itemize}
Résultat : $(1011)_2 + (110)_2 = (10001)_2$.
\end{exemplebox}

\begin{popfigure}
\begin{center}
\begin{tikzpicture}[scale=1.05, every node/.style={font=\large}]
% Retenues (corrigées : colonnes 2^1, 2^2, 2^3)
\node[poporange, font=\small\bfseries] at (-0.3,1.15) {1};
\node[poporange, font=\small\bfseries] at (-1.1,1.15) {1};
\node[poporange, font=\small\bfseries] at (-1.9,1.15) {1};
% Premier nombre (1011)
\node at (-1.9,0.5) {1};
\node at (-1.1,0.5) {0};
\node at (-0.3,0.5) {1};
\node at (0.5,0.5) {1};
% Deuxième nombre (110)
\node[popblue] at (-2.7,-0.3) {$+$};
\node at (-1.1,-0.3) {1};
\node at (-0.3,-0.3) {1};
\node at (0.5,-0.3) {0};
% Ligne
\draw[popdark, thick] (-2.9,-0.75) -- (1.1,-0.75);
% Résultat (10001)
\node[popgreen, font=\bfseries] at (-2.7,-1.25) {1};
\node[popgreen, font=\bfseries] at (-1.9,-1.25) {0};
\node[popgreen, font=\bfseries] at (-1.1,-1.25) {0};
\node[popgreen, font=\bfseries] at (-0.3,-1.25) {0};
\node[popgreen, font=\bfseries] at (0.5,-1.25) {1};
% Légendes colonnes
\node[popmuted, font=\footnotesize] at (0.5,2.0) {$2^0$};
\node[popmuted, font=\footnotesize] at (-0.3,2.0) {$2^1$};
\node[popmuted, font=\footnotesize] at (-1.1,2.0) {$2^2$};
\node[popmuted, font=\footnotesize] at (-1.9,2.0) {$2^3$};
\node[popmuted, font=\footnotesize] at (-2.7,2.0) {$2^4$};
\end{tikzpicture}
\end{center}
\legende{Addition posée $(1011)_2 + (110)_2 = (10001)_2$. En orange, les retenues de $1$ : elles apparaissent dès qu'une colonne atteint $2$.}
\end{popfigure}

\begin{exemplebox}[Addition en base 8 : $(456)_8 + (267)_8$]
Posons et calculons colonne par colonne :
\begin{itemize}
\item Unités : $6 + 7 = 13_{10} = 8 + 5$. On pose $5$, on retient $1$.
\item Huitaines : $5 + 6 + 1 = 12_{10} = 8 + 4$. On pose $4$, on retient $1$.
\item Soixante-quatraines : $4 + 2 + 1 = 7 < 8$. On pose $7$, pas de retenue.
\end{itemize}
Résultat : $(456)_8 + (267)_8 = (745)_8$.
\end{exemplebox}

\courssub{La table d'addition en base 5}

Pour s'entraîner efficacement, rien ne vaut une table complète. Voici la table d'addition en base 5, où chaque case contient le résultat écrit \emph{en base 5} : dès que la somme atteint $5$, on écrit un nombre à deux chiffres.

\begin{popfigure}
\begin{center}
\begin{tikzpicture}[scale=0.95]
% Grille 6x6 (en-têtes + 5 lignes)
\foreach \i in {0,...,5}{
  \draw[popline] (\i,0) -- (\i,-6);
  \draw[popline] (0,-\i) -- (6,-\i);
}
\draw[popdark, thick] (0,0) rectangle (6,-6);
% En-têtes
\node at (0.5,-0.5) {\textbf{+}};
\node[popblue] at (1.5,-0.5) {\textbf{0}};
\node[popblue] at (2.5,-0.5) {\textbf{1}};
\node[popblue] at (3.5,-0.5) {\textbf{2}};
\node[popblue] at (4.5,-0.5) {\textbf{3}};
\node[popblue] at (5.5,-0.5) {\textbf{4}};
\node[popblue] at (0.5,-1.5) {\textbf{0}};
\node[popblue] at (0.5,-2.5) {\textbf{1}};
\node[popblue] at (0.5,-3.5) {\textbf{2}};
\node[popblue] at (0.5,-4.5) {\textbf{3}};
\node[popblue] at (0.5,-5.5) {\textbf{4}};
% Ligne 0
\node at (1.5,-1.5) {0}; \node at (2.5,-1.5) {1}; \node at (3.5,-1.5) {2}; \node at (4.5,-1.5) {3}; \node at (5.5,-1.5) {4};
% Ligne 1
\node at (1.5,-2.5) {1}; \node at (2.5,-2.5) {2}; \node at (3.5,-2.5) {3}; \node at (4.5,-2.5) {4}; \node[poporange, font=\bfseries] at (5.5,-2.5) {10};
% Ligne 2
\node at (1.5,-3.5) {2}; \node at (2.5,-3.5) {3}; \node at (3.5,-3.5) {4}; \node[poporange, font=\bfseries] at (4.5,-3.5) {10}; \node[poporange, font=\bfseries] at (5.5,-3.5) {11};
% Ligne 3
\node at (1.5,-4.5) {3}; \node at (2.5,-4.5) {4}; \node[poporange, font=\bfseries] at (3.5,-4.5) {10}; \node[poporange, font=\bfseries] at (4.5,-4.5) {11}; \node[poporange, font=\bfseries] at (5.5,-4.5) {12};
% Ligne 4
\node at (1.5,-5.5) {4}; \node[poporange, font=\bfseries] at (2.5,-5.5) {10}; \node[poporange, font=\bfseries] at (3.5,-5.5) {11}; \node[poporange, font=\bfseries] at (4.5,-5.5) {12}; \node[poporange, font=\bfseries] at (5.5,-5.5) {13};
\end{tikzpicture}
\end{center}
\legende{Table d'addition en base 5. Les cases orange correspondent aux sommes qui atteignent ou dépassent $5$ : le résultat s'écrit avec deux chiffres, ce qui correspond à une retenue dans une addition posée.}
\end{popfigure}

Lis par exemple la case $4 + 4 = (13)_5$ : en effet $4 + 4 = 8_{10} = 5 + 3 = 1 \times 5 + 3$. Dans une addition posée, cela signifie : on pose $3$, on retient $1$.

\begin{exoresolu}[Addition en base 5 avec vérification]
Calculer en base 5 la somme $(243)_5 + (134)_5$, puis vérifier le résultat en base 10.
\tcblower
\textbf{Addition posée en base 5}, colonne par colonne, de droite à gauche :
\begin{itemize}
\item Unités : $3 + 4 = 7_{10} = 5 + 2$. On pose $2$, on retient $1$.
\item Cinquaines : $4 + 3 + 1 = 8_{10} = 5 + 3$. On pose $3$, on retient $1$.
\item Vingt-cinquaines : $2 + 1 + 1 = 4 < 5$. On pose $4$, pas de retenue.
\end{itemize}
Donc $(243)_5 + (134)_5 = (432)_5$.

\textbf{Vérification en base 10.} Convertissons chaque nombre :
\begin{align*}
(243)_5 &= 2 \times 5^2 + 4 \times 5 + 3 = 50 + 20 + 3 = 73,\\
(134)_5 &= 1 \times 5^2 + 3 \times 5 + 4 = 25 + 15 + 4 = 44,\\
(432)_5 &= 4 \times 5^2 + 3 \times 5 + 2 = 100 + 15 + 2 = 117.
\end{align*}
Or $73 + 44 = 117$ : le résultat est confirmé. Cette vérification par conversion est ton meilleur filet de sécurité le jour de l'examen.
\end{exoresolu}

\begin{exoresolu}[Addition en base 8 avec vérification]
Vérifier par conversion en base 10 que $(456)_8 + (267)_8 = (745)_8$.
\tcblower
Convertissons les trois écritures :
\begin{align*}
(456)_8 &= 4 \times 8^2 + 5 \times 8 + 6 = 256 + 40 + 6 = 302,\\
(267)_8 &= 2 \times 8^2 + 6 \times 8 + 7 = 128 + 48 + 7 = 183,\\
(745)_8 &= 7 \times 8^2 + 4 \times 8 + 5 = 448 + 32 + 5 = 485.
\end{align*}
Or $302 + 183 = 485$ : l'addition en base 8 est correcte.
\end{exoresolu}

\begin{attention}[Pièges de rédaction à éviter]
\begin{itemize}
\item \textbf{Ne jamais écrire un chiffre interdit} : en base 5, les seuls chiffres autorisés sont $0, 1, 2, 3, 4$. Si tu poses un $5$ ou un $7$ dans un résultat en base 5, c'est que tu as oublié une retenue.
\item \textbf{Ne pas oublier la retenue finale} : si la dernière colonne génère une retenue, elle devient un chiffre supplémentaire à gauche. Ainsi $(1)_2 + (1)_2 = (10)_2$, et non $0$.
\item \textbf{Toujours indiquer la base} en indice : $(432)_5$ et $432$ ne désignent pas le même nombre. Sans indice, le correcteur suppose la base 10.
\end{itemize}
\end{attention}

\begin{savaistu}[Des milliards d'additions binaires par seconde]
L'addition binaire que tu viens d'apprendre est littéralement l'opération la plus exécutée au monde. Au cœur de ton téléphone portable, le processeur est construit à partir de millions de minuscules circuits appelés \emph{additionneurs}, qui ne savent faire qu'une chose : additionner des $0$ et des $1$ avec la règle $1 + 1 = 10$, retenue incluse, exactement comme tu viens de le faire. Un additionneur « complet » (\emph{full adder}) calcule pour chaque colonne une somme et une retenue à partir de trois entrées : les deux chiffres et la retenue précédente. Enchaînés en cascade, ces circuits additionnent des nombres de 32 ou 64 chiffres binaires en une fraction de nanoseconde. Quand tu envoies un message Orange Money, que tu écoutes de la musique ou que tu photographies un paysage de l'Ouest camerounais, ton téléphone effectue des \textbf{milliards} d'additions binaires par seconde. La soustraction, la multiplication et même la division des machines sont d'ailleurs ramenées, en dernière analyse, à des suites d'additions binaires. Tu maîtrises donc désormais le geste fondamental de toute l'informatique moderne !
\end{savaistu}

\begin{exercice}[À toi de jouer]
\begin{enumerate}
\item Calculer en base 2 : $(1101)_2 + (1011)_2$, puis vérifier en base 10.
\item Calculer en base 5 : $(304)_5 + (222)_5$, puis vérifier en base 10.
\item Calculer en base 8 : $(375)_8 + (146)_8$, puis vérifier en base 10.
\end{enumerate}
\end{exercice}

\begin{corrige}[Exercice : additions en bases 2, 5 et 8]
\begin{enumerate}
\item Unités : $1+1 = 2$, pose $0$, retient $1$. Deuxaines : $0+1+1 = 2$, pose $0$, retient $1$. Quatraines : $1+0+1 = 2$, pose $0$, retient $1$. Huitaines : $1+1+1 = 3 = (11)_2$, pose $1$, retient $1$ ; la retenue finale s'écrit à gauche. Résultat : $(11000)_2$. Vérification : $(1101)_2 = 13$, $(1011)_2 = 11$, $(11000)_2 = 24$, et $13 + 11 = 24$. Correct.
\item Unités : $4+2 = 6 = 5+1$, pose $1$, retient $1$. Cinquaines : $0+2+1 = 3$, pose $3$. Vingt-cinquaines : $3+2 = 5 = (10)_5$, pose $0$, retient $1$ ; la retenue finale s'écrit à gauche. Résultat : $(1031)_5$. Vérification : $(304)_5 = 79$, $(222)_5 = 62$, $(1031)_5 = 141$, et $79 + 62 = 141$. Correct.
\item Unités : $5+6 = 11_{10} = 8+3$, pose $3$, retient $1$. Huitaines : $7+4+1 = 12_{10} = 8+4$, pose $4$, retient $1$. Soixante-quatraines : $3+1+1 = 5$, pose $5$. Résultat : $(543)_8$. Vérification : $(375)_8 = 253$, $(146)_8 = 102$, $(543)_8 = 355$, et $253 + 102 = 355$. Correct.
\end{enumerate}
\end{corrige}

\begin{aretenir}[L'essentiel sur l'addition en base $a$]
\begin{itemize}
\item On pose l'addition en colonnes, alignées à droite, exactement comme en base 10.
\item Une \cle{retenue} de $1$ apparaît dès que la somme d'une colonne \textbf{atteint ou dépasse $a$} : on pose $S - a$ et on retient $1$.
\item En base 2 : $1 + 1 = (10)_2$ et $1 + 1 + 1 = (11)_2$.
\item Les chiffres du résultat doivent tous être strictement inférieurs à $a$.
\item La \cle{vérification} par conversion en base 10 est systématique : elle sécurise ta copie.
\end{itemize}
\end{aretenir}

\coursec{Multiplication en base $a$}

Dans la section précédente, nous avons appris à \cle{additionner} deux entiers écrits dans une même base $a$ : la technique est la même qu'en base 10, à condition de remplacer le réflexe « dès qu'on atteint dix, je retiens un » par « dès qu'on atteint $a$, je retiens un ». Nous allons maintenant franchir une étape supplémentaire : \cle{multiplier} directement en base $a$, sans repasser par la base 10. Tu vas découvrir que la multiplication n'est rien d'autre qu'une addition répétée, organisée grâce à une \emph{table de multiplication} dressée dans la base considérée. Et en base 2, tu verras que tout devient étonnamment simple : c'est le secret de la rapidité des ordinateurs.

\courssub{Principe : la table de multiplication en base $a$}

\courssub{L'idée : tout se fait dans la base}

Pour poser une multiplication en base 10, tu utilises sans y penser la table de Pythagore apprise au primaire : $7 \times 8 = 56$. Mais cette égalité est écrite \emph{en base 10}. En base $a$, le produit de deux chiffres doit être \textbf{écrit en base $a$} : par exemple, en base 5, $4 \times 3 = 12$ (en base 10), et $12 = 2 \times 5 + 2$, donc $4 \times 3 = \overline{22}^{5}$. Toute la difficulté — et toute la richesse — de la multiplication en base $a$ tient dans cette remarque : \emph{les résultats de la table sont exprimés dans la base $a$}.

\begin{definition}[Table de multiplication en base $a$]
La \cle{table de multiplication} en base $a$ est le tableau donnant, pour tous les chiffres $x$ et $y$ de la base ($0 \le x < a$ et $0 \le y < a$), le produit $x \times y$ \textbf{écrit en base $a$}.
\end{definition}

\begin{exemplebox}[Table de multiplication en base 2]
En base 2, les seuls chiffres sont $0$ et $1$. La table est triviale :
\begin{center}
\begin{tabular}{|c|c|c|}
\hline
\rowcolor{popblueL} $\times$ & $0$ & $1$ \\
\hline
$0$ & $0$ & $0$ \\
\hline
$1$ & $0$ & $1$ \\
\hline
\end{tabular}
\end{center}
Multiplier par $1$ recopie le nombre ; multiplier par $0$ donne $0$. Il n'y a \textbf{aucune retenue} à gérer dans les produits partiels : c'est le rêve du calculateur !
\end{exemplebox}

\begin{exemplebox}[Construction de la table de multiplication en base 5]
Les chiffres de la base 5 sont $0, 1, 2, 3, 4$. Pour remplir la table, on calcule chaque produit en base 10, puis on le convertit en base 5. Par exemple :
\begin{itemize}
\item $2 \times 3 = 6 = 1 \times 5 + 1$, donc $2 \times 3 = \overline{11}^{5}$ ;
\item $3 \times 3 = 9 = 1 \times 5 + 4$, donc $3 \times 3 = \overline{14}^{5}$ ;
\item $3 \times 4 = 12 = 2 \times 5 + 2$, donc $3 \times 4 = \overline{22}^{5}$ ;
\item $4 \times 4 = 16 = 3 \times 5 + 1$, donc $4 \times 4 = \overline{31}^{5}$.
\end{itemize}
La table complète est présentée dans la figure ci-dessous.
\end{exemplebox}

\begin{popfigure}
\begin{center}
\begin{tikzpicture}[scale=0.95]
\foreach \i/\v in {0/0,1/1,2/2,3/3,4/4}{
  \node[fill=popblueL, minimum width=1.15cm, minimum height=0.75cm] at (\i*1.2,0) {\textbf{\v}};
  \node[fill=popblueL, minimum width=1.15cm, minimum height=0.75cm] at (-1.2,-\i*0.8) {\textbf{\v}};
}
\node[fill=poporangeL, minimum width=1.15cm, minimum height=0.75cm] at (-1.2,0) {$\times$};
\foreach \i/\row in {0/{0,0,0,0,0},1/{0,1,2,3,4},2/{0,2,4,11,13},3/{0,3,11,14,22},4/{0,4,13,22,31}}{
  \foreach [count=\j from 0] \c in \row {
    \node[draw=popline, fill=white, minimum width=1.15cm, minimum height=0.75cm] at (\j*1.2,-\i*0.8) {\c};
  }
}
\end{tikzpicture}
\end{center}
\legende{Table de multiplication en base 5 : chaque produit est écrit \emph{en base 5}. Par exemple, la case $(4,3)$ contient $22$ car $4 \times 3 = 12 = \overline{22}^{5}$.}
\end{popfigure}

\begin{attention}[Ne pas confondre les deux écritures]
Quand tu remplis ou utilises une table en base $a$, le contenu des cases est un nombre \textbf{écrit en base $a$}. Ainsi, en base 5, la case $4 \times 4$ contient $31$, qui se lit « trois-un » et vaut $3 \times 5 + 1 = 16$ en base 10. Ne lis jamais « trente et un » !
\end{attention}

\courssub{Technique opératoire : la multiplication posée}

La disposition est exactement celle que tu connais en base 10 : on multiplie le premier facteur par chaque chiffre du second, en \textbf{décalant} chaque produit partiel d'un rang vers la gauche, puis on \textbf{additionne} les produits partiels. Deux points exigent toute ta vigilance : les produits partiels utilisent la table en base $a$ (avec retenues), et l'addition finale se fait \emph{en base $a$}.

\begin{methode}[Multiplier deux entiers écrits en base $a$]
Pour calculer $\overline{x}^{a} \times \overline{y}^{a}$ :
\begin{enumerate}
\item \textbf{Construire (ou avoir sous les yeux) la table de multiplication en base $a$}, dont tous les résultats sont écrits en base $a$.
\item \textbf{Poser les produits partiels} : multiplier $\overline{x}^{a}$ par chaque chiffre de $\overline{y}^{a}$, en commençant par le chiffre des unités, et décaler chaque produit partiel d'un rang vers la gauche par rapport au précédent. Les retenues se gèrent comme en base 10, mais on retient dès qu'on atteint $a$.
\item \textbf{Additionner les produits partiels en base $a$}, colonne par colonne, avec les retenues de l'addition en base $a$.
\item \textbf{Vérifier} le résultat en convertissant les deux facteurs et le produit en base 10.
\end{enumerate}
\end{methode}

\begin{exemplebox}[Multiplication en base 2 : $\overline{101}^{2} \times \overline{11}^{2}$]
Posons l'opération. La table de base 2 ne contient que des $0$ et des $1$ : multiplier par $1$ recopie le nombre, multiplier par $0$ donne $0$.
\begin{itemize}
\item Premier produit partiel : $\overline{101}^{2} \times 1 = \overline{101}^{2}$.
\item Deuxième produit partiel : $\overline{101}^{2} \times 1 = \overline{101}^{2}$, décalé d'un rang vers la gauche, ce qui donne $\overline{1010}^{2}$.
\item Addition finale en base 2 : $\overline{101}^{2} + \overline{1010}^{2} = \overline{1111}^{2}$ (colonne par colonne : $1+0=1$, $0+1=1$, $1+0=1$, $0+1=1$).
\end{itemize}
Donc $\overline{101}^{2} \times \overline{11}^{2} = \overline{1111}^{2}$. La figure ci-dessous montre l'opération posée.
\end{exemplebox}

\begin{popfigure}
\begin{center}
\begin{tikzpicture}[scale=1.05, every node/.style={font=\large}]
\node at (0,0) {$\overline{101}^{2}$};
\node at (0,-0.7) {$\times\ \overline{11}^{2}$};
\draw[popdark, thick] (-1.6,-1.05) -- (1.6,-1.05);
\node[popblue] at (-0.35,-1.5) {$101$};
\node[popblue] at (0.35,-2.1) {$101\,$};
\draw[popdark, thick] (-1.6,-2.45) -- (1.6,-2.45);
\node[poporange] at (0,-2.95) {$\overline{1111}^{2}$};
\node[popmuted, font=\small] at (2.9,-1.5) {$\leftarrow\ \overline{101}^{2}\times 1$};
\node[popmuted, font=\small] at (3.4,-2.1) {$\leftarrow\ \overline{101}^{2}\times 1$, décalé d'un rang};
\node[popmuted, font=\small] at (2.9,-2.95) {$\leftarrow$ somme en base 2};
\end{tikzpicture}
\end{center}
\legende{Multiplication posée de $\overline{101}^{2}$ par $\overline{11}^{2}$ : deux produits partiels décalés, puis une addition en base 2.}
\end{popfigure}

\begin{exoresolu}[Multiplication en base 5 : $\overline{234}^{5} \times \overline{3}^{5}$]
Énoncé. Calculer $\overline{234}^{5} \times \overline{3}^{5}$ en posant l'opération en base 5, puis vérifier le résultat en base 10.
\tcblower
Solution détaillée. On multiplie chaque chiffre de $\overline{234}^{5}$ par $3$, de droite à gauche, en utilisant la table de base 5 :
\begin{itemize}
\item $4 \times 3 = 12 = \overline{22}^{5}$ : on écrit $2$ et on retient $2$ ;
\item $3 \times 3 + 2 = 9 + 2 = 11 = \overline{21}^{5}$ : on écrit $1$ et on retient $2$ ;
\item $2 \times 3 + 2 = 6 + 2 = 8 = \overline{13}^{5}$ : on écrit $3$ et on retient $1$ ;
\item on abaisse la dernière retenue : $1$.
\end{itemize}
Le résultat s'écrit donc, en lisant de gauche à droite :
\[ \overline{234}^{5} \times \overline{3}^{5} = \overline{1312}^{5}. \]
Vérification en base 10 : $\overline{234}^{5} = 2 \times 25 + 3 \times 5 + 4 = 69$ et $69 \times 3 = 207$. Or
\[ \overline{1312}^{5} = 1 \times 125 + 3 \times 25 + 1 \times 5 + 2 = 125 + 75 + 5 + 2 = 207. \]
Les deux résultats coïncident : le calcul est correct.
\end{exoresolu}

\begin{attention}[L'erreur classique : additionner les produits partiels en base 10]
Lorsqu'une multiplication posée en base $a$ comporte plusieurs produits partiels, l'\textbf{addition finale doit se faire en base $a$}, avec la règle « dès que la somme d'une colonne atteint $a$, on retient $1$ ». Additionner les produits partiels avec les réflexes de la base 10 (retenue à partir de dix) est l'erreur la plus fréquente au Probatoire. Exemple : en base 2, $1 + 1 = \overline{10}^{2}$ : on écrit $0$ et on retient $1$, et non pas « $2$ » !
\end{attention}

\courssub{Le cas particulier de la base 2 : décalages et additions}

En base 2, le second facteur ne contient que des chiffres $0$ et $1$. Chaque produit partiel est donc soit nul (multiplication par $0$), soit une simple \textbf{copie décalée} du premier facteur (multiplication par $1$). La multiplication binaire se ramène ainsi à une suite de \cle{décalages} suivie d'une \cle{addition} binaire.

\begin{exemplebox}[Décalages en action]
Reprenons $\overline{101}^{2} \times \overline{11}^{2}$ : le multiplicateur $\overline{11}^{2}$ possède deux chiffres $1$, donc on additionne deux copies de $\overline{101}^{2}$, l'une décalée de $0$ rang, l'autre de $1$ rang : $\overline{101}^{2} + \overline{1010}^{2} = \overline{1111}^{2}$. De même, pour calculer $\overline{1011}^{2} \times \overline{101}^{2}$, on additionnerait $\overline{1011}^{2}$ (décalage $0$) et $\overline{101100}^{2}$ (décalage $2$), le chiffre du milieu étant $0$.
\end{exemplebox}

\begin{savaistu}[Pourquoi les processeurs calculent-ils si vite ?]
Ton téléphone et ton ordinateur représentent tous les nombres en binaire. Or tu viens de le voir : en base 2, la multiplication se réduit à des \emph{décalages} (déplacer les bits vers la gauche) et des \emph{additions}. Ce sont les deux opérations les plus simples à réaliser avec des circuits électroniques ! Un processeur effectue des milliards de telles opérations par seconde. Les ingénieurs parlent d'ailleurs de « décalage de bits » (\emph{bit shifting}) : décaler un registre d'un cran vers la gauche revient exactement à multiplier par $2$. C'est la même idée que ton réflexe « multiplier par dix, c'est ajouter un zéro » en base 10.
\end{savaistu}

\courssub{Synthèse : multiplier par $a$ en base $a$}

En base 10, tu sais depuis longtemps que $37 \times 10 = 370$ : multiplier par dix revient à \textbf{ajouter un zéro à droite}. Ce phénomène n'a rien de magique : il vient de ce que $10$ est la \emph{base} du système décimal. Le même phénomène se produit dans \textbf{toute} base $a$.

\begin{propriete}[Multiplication par la base]
Soit $a \ge 2$ un entier et $N$ un entier naturel dont l'écriture en base $a$ est
\[ N = \overline{x_n x_{n-1} \ldots x_1 x_0}^{a}. \]
Alors l'écriture de $a \times N$ en base $a$ s'obtient en \textbf{ajoutant un zéro à droite} de l'écriture de $N$ :
\[ a \times N = \overline{x_n x_{n-1} \ldots x_1 x_0 0}^{a}. \]
Autrement dit, en base $a$, multiplier par $a$ décale tous les chiffres d'un rang vers la gauche. Plus généralement, multiplier par $a^k$ revient à ajouter $k$ zéros à droite.
\end{propriete}

\begin{experience}[Démonstration guidée]
Écrivons $N$ sous forme décomposée : $N = x_n a^n + x_{n-1} a^{n-1} + \cdots + x_1 a + x_0$, avec $0 \le x_i < a$ et $x_n \neq 0$.
\begin{enumerate}
\item Multiplions les deux membres par $a$ :
\[ a \times N = x_n a^{n+1} + x_{n-1} a^{n} + \cdots + x_1 a^{2} + x_0 a. \]
\item Observons : chaque exposant a augmenté de $1$, et le terme constant (en $a^0$) a disparu, c'est-à-dire qu'il vaut $0$.
\item L'expression obtenue est exactement la décomposition en base $a$ du nombre $\overline{x_n x_{n-1} \ldots x_1 x_0 0}^{a}$, puisque les coefficients $x_i$ vérifient toujours $0 \le x_i < a$ et que le coefficient dominant $x_n$ est non nul.
\end{enumerate}
Conclusion : multiplier par $a$ ajoute un zéro à droite. Le raisonnement se répète $k$ fois pour $a^k$, ce qui ajoute $k$ zéros.
\end{experience}

\begin{exemplebox}[Applications directes]
\begin{itemize}
\item En base 2 : $\overline{101}^{2} \times \overline{10}^{2} = \overline{1010}^{2}$. Vérification : $5 \times 2 = 10$ et $\overline{1010}^{2} = 8 + 2 = 10$.
\item En base 5 : $\overline{234}^{5} \times \overline{10}^{5} = \overline{2340}^{5}$. Vérification : $69 \times 5 = 345$ et $\overline{2340}^{5} = 2 \times 125 + 3 \times 25 + 4 \times 5 = 250 + 75 + 20 = 345$.
\item En base 10 : $57 \times 100 = 5700$ (deux zéros, car $100 = 10^2$).
\end{itemize}
\end{exemplebox}

\begin{aretenir}[L'essentiel de la multiplication en base $a$]
\begin{itemize}
\item La \textbf{table de multiplication} en base $a$ donne tous les produits de chiffres \emph{écrits en base $a$}.
\item La multiplication posée suit la même disposition qu'en base 10 : \textbf{produits partiels décalés}, puis \textbf{addition finale en base $a$}.
\item En \textbf{base 2}, tout se ramène à des décalages et des additions : c'est le fondement du calcul informatique.
\item Multiplier par $a$ revient à \textbf{ajouter un zéro à droite} ; multiplier par $a^k$ ajoute $k$ zéros.
\item On \textbf{vérifie toujours} le résultat par conversion en base 10.
\end{itemize}
\end{aretenir}

\begin{exercice}[Pour t'entraîner]
\begin{enumerate}
\item Calculer en base 2 : $\overline{110}^{2} \times \overline{11}^{2}$, puis vérifier en base 10.
\item Calculer en base 5 : $\overline{43}^{5} \times \overline{4}^{5}$, puis vérifier en base 10.
\item Sans poser d'opération, écrire en base 5 le produit $\overline{3204}^{5} \times \overline{100}^{5}$.
\end{enumerate}
\end{exercice}

\begin{corrige}[Exercice]
\begin{enumerate}
\item Produits partiels : $\overline{110}^{2}$ et $\overline{1100}^{2}$. Somme en base 2 : $\overline{110}^{2} + \overline{1100}^{2} = \overline{10010}^{2}$. Vérification : $\overline{110}^{2} = 6$, $\overline{11}^{2} = 3$, $6 \times 3 = 18$ et $\overline{10010}^{2} = 16 + 2 = 18$.
\item $3 \times 4 = 12 = \overline{22}^{5}$ : on écrit $2$, retenue $2$ ; $4 \times 4 + 2 = 18 = \overline{33}^{5}$ : on écrit $3$, retenue $3$ ; on abaisse $3$. Résultat : $\overline{43}^{5} \times \overline{4}^{5} = \overline{332}^{5}$. Vérification : $\overline{43}^{5} = 23$, $23 \times 4 = 92$ et $\overline{332}^{5} = 3 \times 25 + 3 \times 5 + 2 = 92$.
\item Multiplier par $\overline{100}^{5} = 5^2$ ajoute deux zéros à droite : $\overline{3204}^{5} \times \overline{100}^{5} = \overline{320400}^{5}$.
\end{enumerate}
\end{corrige}

\newpage

\coursec{Activité d'intégration}

\coursec{Arithmétique : division euclidienne et numération}

\courssub{Objectifs de l'activité}

À l'issue de cette activité d'intégration, tu devras être capable de :
\begin{itemize}
\item déterminer le quotient et le reste d'une division euclidienne et traduire le résultat par la relation $a = bq + r$ avec $0 \leq r < b$ ;
\item écrire un entier naturel dans un système de numération de base $a$ par la méthode des divisions successives ;
\item convertir un nombre d'une base à une autre, en passant par la base $10$ ou par regroupement direct de bits (entre les bases $2$ et $8$) ;
\item additionner et multiplier des entiers écrits dans le même système de numération, en gérant correctement les retenues ;
\item vérifier un calcul effectué en base $a$ en le traduisant en base $10$.
\end{itemize}

\courssub{Situation}

Le Lycée Bilingue de Nkongsamba vient de recevoir, grâce à un partenariat avec une entreprise de télécommunications, un lot de $200$ nouveaux casiers destinés aux élèves de Première. Le proviseur, M.~Etoundi, souhaite moderniser la gestion de ces casiers de trois façons :

\begin{itemize}
\item les casiers seront rangés dans le couloir par \textbf{rangées de $12$} ; il faut savoir combien de rangées complètes on pourra former et combien de casiers resteront sur la dernière rangée ;
\item un futur \textbf{système électronique d'ouverture} (financé l'année prochaine) exigera que chaque casier porte son numéro écrit en \textbf{base $2$} (binaire), car le circuit électronique ne connaît que deux états : $0$ et $1$ ;
\item en attendant, chaque élève recevra un \textbf{badge} sur lequel le numéro du casier sera imprimé en \textbf{base $8$} (octal), un format plus court que le binaire et facile à graver.
\end{itemize}

Le gardien du lycée, M.~Bekono, a déjà reçu trois badges gravés en binaire : $(1100100)_2$, $(1011011)_2$ et $(11111001)_2$. Il doit retrouver les numéros correspondants en base $10$ pour les reporter dans le registre, puis en base $8$ pour les badges définitifs.

Par ailleurs, le gardien signale que le casier de numéro binaire $(1101)_2$ est suivi d'une zone provisoire de $(1011)_2$ casiers supplémentaires installés dans l'annexe. Il veut connaître le numéro binaire du \emph{dernier} casier de l'annexe, obtenu en ajoutant $(1011)_2$ au numéro $(1101)_2$.

Enfin, pour s'entraîner, chaque groupe d'élèves devra fabriquer son propre « code secret » : choisir un numéro de casier, l'écrire en binaire et en octal, puis le faire décoder par un autre groupe.

\courssub{Consignes}

Travail en groupes de quatre élèves. Chaque groupe rédige ses réponses sur une feuille qui sera corrigée par un autre groupe lors de la restitution.

\begin{enumerate}
\item \textbf{Rangement des casiers.} On range les $200$ casiers par rangées de $12$.
\begin{enumerate}
\item Effectue la division euclidienne de $200$ par $12$.
\item Traduis le résultat par une égalité de la forme $a = bq + r$ en précisant la condition sur $r$.
\item Réponds à la question du proviseur : combien de rangées complètes et combien de casiers sur la rangée incomplète ?
\end{enumerate}
\item \textbf{Numérotation des dix premiers casiers.} En utilisant la méthode des divisions successives, écris les entiers de $1$ à $10$ en base $2$, puis en base $8$. Présente les résultats dans un tableau à trois colonnes : base $10$, base $2$, base $8$.
\item \textbf{Décodage des badges.}
\begin{enumerate}
\item Convertis le numéro $(1100100)_2$ en base $10$ en détaillant la décomposition en puissances de $2$.
\item Sans repasser par la base $10$, convertis $(1100100)_2$ en base $8$ par regroupement des bits par paquets de trois (en partant de la droite). Justifie pourquoi cette méthode fonctionne.
\item Procède de même pour $(1011011)_2$ et $(11111001)_2$ : donne pour chacun l'écriture en base $10$ puis en base $8$.
\end{enumerate}
\item \textbf{L'annexe du gardien.}
\begin{enumerate}
\item Pose et effectue l'addition $(1101)_2 + (1011)_2$ en base $2$, en indiquant les retenues.
\item Vérifie le résultat en convertissant les trois nombres en base $10$.
\item Le gardien affirme : « le dernier casier de l'annexe porte le numéro octal $(30)_8$ ». A-t-il raison ?
\end{enumerate}
\item \textbf{Fabrication d'un code secret.} Choisis un numéro de casier compris entre $50$ et $200$. Écris-le en base $2$ et en base $8$, puis multiplie-le par $(11)_2$ directement en binaire. Remets ton code (uniquement les écritures en base $2$ et $8$) à un autre groupe, qui devra retrouver le numéro en base $10$ et vérifier ta multiplication.
\end{enumerate}

\courssub{Ressources à mobiliser}

\begin{aretenir}[Boîte à outils de l'activité]
\begin{itemize}
\item \textbf{Division euclidienne :} pour $a \in \mathbb{N}$ et $b \in \mathbb{N}^*$, il existe un unique couple $(q, r)$ tel que $a = bq + r$ avec $0 \leq r < b$.
\item \textbf{Décomposition en base $a$ :} tout entier $N$ s'écrit de manière unique
\[ N = x_n a^n + x_{n-1} a^{n-1} + \cdots + x_1 a + x_0, \quad 0 \leq x_i < a,\ x_n \neq 0, \]
et on note $N = \overline{x_n x_{n-1} \ldots x_1 x_0}^{\,a}$ ou $(x_n x_{n-1} \ldots x_0)_a$.
\item \textbf{De la base $10$ vers la base $a$ :} divisions successives par $a$ ; le résultat se lit en remontant les restes, du dernier au premier.
\item \textbf{De la base $a$ vers la base $10$ :} on développe la somme des puissances de $a$.
\item \textbf{Lien base $2$ -- base $8$ :} comme $8 = 2^3$, chaque chiffre octal correspond à un bloc de $3$ bits ; on regroupe les bits par paquets de trois à partir de la droite.
\item \textbf{Addition en base $a$ :} on additionne chiffre à chiffre ; dès qu'une somme atteint ou dépasse $a$, on retranche $a$ et on reporte une retenue de $1$.
\item \textbf{Multiplication en base $a$ :} on utilise la technique usuelle (produits partiels décalés et addition) en effectuant tous les calculs intermédiaires en base $a$.
\end{itemize}
\end{aretenir}

\courssub{Démarche et corrigé commenté}

\begin{exoresolu}[Corrigé commenté de l'activité]
\textbf{Question 1.} Division euclidienne de $200$ par $12$.

\textbf{Question 2.} Écriture des entiers de $1$ à $10$ en bases $2$ et $8$.

\textbf{Question 3.} Décodage des trois badges binaires.

\textbf{Question 4.} Addition binaire de l'annexe et vérification.

\textbf{Question 5.} Exemple de code secret fabriqué par un groupe.

\tcblower

\textbf{Question 1 : le rangement des casiers.}

On cherche combien de fois $12$ « rentre » dans $200$. Comme $12 \times 16 = 192$ et $12 \times 17 = 204 > 200$, le quotient est $q = 16$ et le reste est $r = 200 - 192 = 8$. La division euclidienne s'écrit :
\[ 200 = 12 \times 16 + 8, \quad \text{avec } 0 \leq 8 < 12. \]
La condition $0 \leq r < b$ est bien vérifiée : $8 < 12$. Le proviseur peut donc former \cle{16 rangées complètes} de $12$ casiers, et il restera \cle{8 casiers} sur une dix-septième rangée incomplète.

\textbf{Question 2 : les dix premiers numéros.}

Pour passer de la base $10$ à la base $2$, on effectue des divisions successives par $2$ et on lit les restes de bas en haut. Par exemple pour $6$ : $6 = 2 \times 3 + 0$, puis $3 = 2 \times 1 + 1$, puis $1 = 2 \times 0 + 1$. En remontant les restes : $6 = (110)_2$. Pour la base $8$, une seule division suffit pour les nombres inférieurs à $8$ ; pour $9$ et $10$ : $9 = 8 \times 1 + 1$ donc $9 = (11)_8$, et $10 = 8 \times 1 + 2$ donc $10 = (12)_8$.

\begin{center}
\begin{tabular}{|c|c|c|}
\hline
\rowcolor{popblueL} Base $10$ & Base $2$ & Base $8$ \\ \hline
$1$ & $(1)_2$ & $(1)_8$ \\ \hline
$2$ & $(10)_2$ & $(2)_8$ \\ \hline
$3$ & $(11)_2$ & $(3)_8$ \\ \hline
$4$ & $(100)_2$ & $(4)_8$ \\ \hline
$5$ & $(101)_2$ & $(5)_8$ \\ \hline
$6$ & $(110)_2$ & $(6)_8$ \\ \hline
$7$ & $(111)_2$ & $(7)_8$ \\ \hline
$8$ & $(1000)_2$ & $(10)_8$ \\ \hline
$9$ & $(1001)_2$ & $(11)_8$ \\ \hline
$10$ & $(1010)_2$ & $(12)_8$ \\ \hline
\end{tabular}
\end{center}

\textbf{Question 3 : décodage des badges.}

\emph{a) Conversion de $(1100100)_2$ en base $10$.} On développe en puissances de $2$, de droite à gauche :
\begin{align*}
(1100100)_2 &= 1 \times 2^6 + 1 \times 2^5 + 0 \times 2^4 + 0 \times 2^3 + 1 \times 2^2 + 0 \times 2^1 + 0 \times 2^0 \\
&= 64 + 32 + 0 + 0 + 4 + 0 + 0 = 100.
\end{align*}
Le badge $(1100100)_2$ correspond donc au casier numéro $100$.

\emph{b) Passage direct en base $8$.} Puisque $8 = 2^3$, un bloc de trois bits représente exactement un chiffre entre $0$ et $7$, c'est-à-dire un chiffre octal. On regroupe les bits par paquets de trois \emph{à partir de la droite}, en complétant éventuellement par des zéros à gauche :
\[ (1100100)_2 = (\underbrace{001}_{1}\ \underbrace{100}_{4}\ \underbrace{100}_{4})_2 = (144)_8. \]
Vérification en base $10$ : $(144)_8 = 1 \times 8^2 + 4 \times 8 + 4 = 64 + 32 + 4 = 100$. C'est cohérent.

\emph{c) Les deux autres badges.} Pour $(1011011)_2$ :
\[ (1011011)_2 = 2^6 + 2^4 + 2^3 + 2^1 + 2^0 = 64 + 16 + 8 + 2 + 1 = 91. \]
Regroupement par blocs de trois : $(001\ 011\ 011)_2 = (133)_8$. Vérification : $1 \times 64 + 3 \times 8 + 3 = 91$.

Pour $(11111001)_2$ :
\[ (11111001)_2 = 2^7 + 2^6 + 2^5 + 2^4 + 2^3 + 2^0 = 128 + 64 + 32 + 16 + 8 + 1 = 249. \]
Regroupement : $(011\ 111\ 001)_2 = (371)_8$. Vérification : $3 \times 64 + 7 \times 8 + 1 = 192 + 56 + 1 = 249$.

\textbf{Question 4 : l'annexe du gardien.}

\emph{a) Addition en base $2$.} On pose l'addition en colonnes, de droite à gauche, avec les retenues :
\[
\begin{array}{c@{}c@{}c@{}c@{}c@{}c}
& & 1 & 1 & 0 & 1 \\
+ & & 1 & 0 & 1 & 1 \\ \cline{2-6}
& 1 & 1 & 0 & 0 & 0
\end{array}
\]
Détail des retenues :
\begin{itemize}
\item colonne des unités ($2^0$) : $1 + 1 = 10_2$, on écrit $0$, on retient $1$ ;
\item colonne des $2^1$ : $0 + 1 + 1 = 10_2$, on écrit $0$, on retient $1$ ;
\item colonne des $2^2$ : $1 + 0 + 1 = 10_2$, on écrit $0$, on retient $1$ ;
\item colonne des $2^3$ : $1 + 1 + 1 = 11_2$, on écrit $1$, on retient $1$ que l'on pose devant.
\end{itemize}
Résultat : $(1101)_2 + (1011)_2 = (11000)_2$.

\emph{b) Vérification en base $10$.} $(1101)_2 = 8 + 4 + 1 = 13$ et $(1011)_2 = 8 + 2 + 1 = 11$. Or $13 + 11 = 24$, et $(11000)_2 = 16 + 8 = 24$. Le calcul binaire est correct.

\emph{c) L'affirmation du gardien.} Convertissons $(11000)_2$ en octal : $(011\ 000)_2 = (30)_8$. Le gardien a raison : le dernier casier de l'annexe porte bien le numéro octal $(30)_8$, c'est-à-dire $24$ en base $10$.

\textbf{Question 5 : exemple de code secret.}

Un groupe choisit le numéro $75$. Divisions successives par $2$ :
$75 = 2 \times 37 + 1$,
$37 = 2 \times 18 + 1$,
$18 = 2 \times 9 + 0$,
$9 = 2 \times 4 + 1$,
$4 = 2 \times 2 + 0$,
$2 = 2 \times 1 + 0$,
$1 = 2 \times 0 + 1$.
En remontant les restes : $75 = (1001011)_2$.
Par regroupement par blocs de trois : $(001\ 001\ 011)_2 = (113)_8$.

Multiplication par $(11)_2$ en binaire : multiplier par $(11)_2$ revient à additionner le nombre avec lui-même décalé d'un rang vers la gauche :
\[
\begin{array}{r@{}c@{}c@{}c@{}c@{}c@{}c@{}c@{}c@{}c}
& & & 1 & 0 & 0 & 1 & 0 & 1 & 1 \\
\times & & & & & & & & 1 & 1 \\ \cline{2-10}
& & & 1 & 0 & 0 & 1 & 0 & 1 & 1 \\
+ & & 1 & 0 & 0 & 1 & 0 & 1 & 1 & 0 \\ \cline{2-10}
& & 1 & 1 & 1 & 0 & 0 & 0 & 0 & 1
\end{array}
\]
Résultat : $(1001011)_2 \times (11)_2 = (11100001)_2$.

Vérification en base $10$ : $(11)_2 = 3$ et $75 \times 3 = 225$ ; or $(11100001)_2 = 128 + 64 + 32 + 1 = 225$. Le groupe décodeur retrouve $75$ à partir de $(1001011)_2$ ou de $(113)_8$ et confirme la multiplication.
\end{exoresolu}

\courssub{Bilan de l'activité}

Cette activité a permis de mettre en évidence plusieurs idées essentielles de la leçon :

\begin{itemize}
\item La \cle{division euclidienne} n'est pas qu'un calcul scolaire : elle répond à une question concrète de rangement (combien de paquets complets, combien de restants), et c'est elle qui fonde la méthode des \cle{divisions successives} pour écrire un nombre en base $a$.
\item L'écriture d'un nombre dépend de la \cle{base} choisie, mais le nombre lui-même ne change pas : $(1100100)_2$, $(144)_8$ et $100$ désignent le même casier. La base $10$ sert de « pont » universel pour vérifier toutes les conversions.
\item Entre les bases $2$ et $8$, le passage est \emph{direct} grâce à la relation $8 = 2^3$ : le regroupement des bits par blocs de trois évite tout calcul lourd. C'est exactement le principe utilisé en informatique, où l'octal et l'hexadécimal servent à écrire les données binaires de façon compacte.
\item Les opérations en base $a$ suivent les mêmes règles qu'en base $10$, à condition de se rappeler que la retenue vaut $1$ dès que la somme atteint $a$ (et non $10$). La vérification en base $10$ reste le réflexe le plus sûr pour contrôler un résultat.
\item Enfin, la fabrication et le décodage croisé des codes secrets ont montré qu'un codage fiable exige une convention claire (la base doit être connue du destinataire) et une vérification systématique : compétence précieuse dans les métiers du numérique, des télécommunications et de la banque, au Cameroun comme ailleurs.
\end{itemize}

\newpage

\coursec{Méthodes et exercices résolus}

\coursec{Leçon 21 : Arithmétique -- Division euclidienne et numération}

\noindent Bienvenue dans cette leçon fondamentale. L'arithmétique des entiers est le socle sur lequel repose toute l'algèbre et l'informatique. Que ce soit pour partager équitablement une récolte de cacao entre coopérateurs, coder un message pour le transmettre via le réseau MTN ou Orange, ou simplement comprendre le fonctionnement interne de ton smartphone, la division euclidienne et les changements de base sont des outils indispensables. Nous allons les maîtriser avec rigueur et méthode, exactement comme l'exige le Probatoire.

\courssub{Division euclidienne dans $\mathbb{N}$}

\begin{definition}[Théorème de la division euclidienne]
Soient $a \in \mathbb{N}$ (le \cle{dividende}) et $b \in \mathbb{N}^{*}$ (le \cle{diviseur}). Il existe un \cle{couple unique} $(q, r)$ d'entiers naturels tel que :
\[ a = bq + r \qquad \text{avec} \qquad 0 \leq r < b. \]
$q$ est le \cle{quotient} et $r$ le \cle{reste}.
Si $r = 0$, on dit que $b$ \cle{divise} $a$ (ou que $a$ est un \cle{multiple} de $b$), ce que l'on note $b \mid a$.
\end{definition}

\begin{propriete}[Critère de divisibilité]
$b \mid a \quad \Longleftrightarrow \quad \text{le reste de la division euclidienne de } a \text{ par } b \text{ est nul } (r=0)$.
\end{propriete}

\begin{methode}[Quotient et reste d'une division euclidienne]
Pour effectuer la division euclidienne de $a \in \mathbb{N}$ par $b \in \mathbb{N}^{*}$ :
\begin{enumerate}
\item \textbf{Poser} la division (ou encadrer $a$ entre deux multiples consécutifs de $b$) : on cherche l'unique entier $q$ tel que
\[ bq \leq a < b(q+1). \]
\item \textbf{Calculer} le reste : $r = a - bq$.
\item \textbf{Vérifier impérativement} la double condition :
\[ a = bq + r \quad \text{et} \quad 0 \leq r < b. \]
\item \textbf{Rédiger la conclusion} : « Le quotient est $q$ et le reste est $r$ ».
\end{enumerate}
\textbf{Réflexes :} Si $a < b$, alors $q = 0$ et $r = a$. Si $b \mid a$, alors $r = 0$. La condition $0 \leq r < b$ est ce qui garantit l'\cle{unicité} du couple $(q, r)$ : ne jamais l'oublier dans la rédaction.
\end{methode}

\begin{exoresolu}[Division euclidienne de 2\,847 par 13]
Effectuer la division euclidienne de $2\,847$ par $13$. Écrire l'égalité euclidienne correspondante et vérifier.
\tcblower
\textbf{Étape 1 : recherche du quotient.} On encadre $2\,847$ entre deux multiples consécutifs de $13$.
\begin{itemize}
\item $13 \times 200 = 2\,600$ ; reste $2\,847 - 2\,600 = 247$.
\item $13 \times 19 = 247$ exactement.
\end{itemize}
Donc $13 \times 219 = 2\,600 + 247 = 2\,847$. Le quotient est $q = 219$.

\textbf{Étape 2 : calcul du reste.}
\[ r = a - bq = 2\,847 - 13 \times 219 = 2\,847 - 2\,847 = 0. \]

\textbf{Étape 3 : vérification.}
\[ 2\,847 = 13 \times 219 + 0 \quad \text{avec} \quad 0 \leq 0 < 13. \]
La condition sur le reste est satisfaite.

\textbf{Conclusion :} Dans la division euclidienne de $2\,847$ par $13$, le quotient est $\boxed{219}$ et le reste est $\boxed{0}$ ; on dit que $2\,847$ est divisible par $13$ (ou que $13$ divise $2\,847$).
\end{exoresolu}

\begin{exoresolu}[Un problème concret de partage]
Un transporteur de la ligne Douala--Bafoussam dispose de $1\,000$ billets de \SI{500}{\text{F}} à regrouper en liasses de $24$ billets pour ses caisses. Combien de liasses complètes peut-il former ? Combien de billets resteront isolés ?
\tcblower
Il s'agit de la division euclidienne de $a = 1\,000$ par $b = 24$.

\textbf{Étape 1 : encadrement.} On cherche $q$ tel que $24q \leq 1\,000 < 24(q+1)$.
\[ 24 \times 40 = 960 \quad ; \quad 24 \times 41 = 984 \quad ; \quad 24 \times 42 = 1\,008 > 1\,000. \]
Donc $24 \times 41 = 984 \leq 1\,000 < 1\,008 = 24 \times 42$, d'où $q = 41$.

\textbf{Étape 2 : reste.}
\[ r = 1\,000 - 24 \times 41 = 1\,000 - 984 = 16. \]

\textbf{Étape 3 : vérification.}
\[ 1\,000 = 24 \times 41 + 16 \quad \text{avec} \quad 0 \leq 16 < 24. \quad \checkmark \]

\textbf{Conclusion :} Le transporteur peut former $\boxed{41}$ liasses complètes (soit $41 \times 24 \times 500 = 492\,000$ F) et il restera $\boxed{16}$ billets isolés (soit $8\,000$ F).
\end{exoresolu}

\begin{methode}[Exploiter l'égalité euclidienne $a = bq + r$ avec $0 \leq r < b$]
Devant un énoncé donnant des informations sur une division euclidienne :
\begin{enumerate}
\item \textbf{Traduire} chaque information par une équation : $a = bq + r$, avec la contrainte $0 \leq r < b$.
\item \textbf{Assembler} les équations en un système et le résoudre (souvent par substitution).
\item \textbf{Filtrer} les solutions trouvées grâce à la condition $0 \leq r < b$ (c'est elle qui élimine les solutions parasites).
\item \textbf{Vérifier} en remplaçant dans l'égalité de départ.
\end{enumerate}
\textbf{Piège classique :} Une égalité $a = bq + r$ seule ne suffit pas à affirmer que $q$ et $r$ sont le quotient et le reste : il faut \textbf{en plus} $0 \leq r < b$. Par exemple, $100 = 7 \times 14 + 2$ est l'égalité euclidienne ($0 \leq 2 < 7$), mais $100 = 7 \times 13 + 9$ n'en est pas une car le reste $9$ n'est pas strictement inférieur au diviseur $7$ ($9 \geq 7$).
\end{methode}

\begin{exoresolu}[Retrouver le dividende et le diviseur]
Dans une division euclidienne, le quotient est $12$ et le reste est $7$. De plus, la somme du dividende et du diviseur est $228$. Déterminer le dividende $a$ et le diviseur $b$.
\tcblower
\textbf{Étape 1 : traduction.} Les données de l'énoncé s'écrivent :
\[ \begin{cases} a = 12b + 7 \\ a + b = 228 \end{cases} \quad \text{avec la condition } 0 \leq 7 < b. \]

\textbf{Étape 2 : résolution par substitution.} Remplaçons $a$ par $12b + 7$ dans la deuxième équation :
\begin{align*}
(12b + 7) + b &= 228 \\
13b + 7 &= 228 \\
13b &= 221 \\
b &= \frac{221}{13} = 17.
\end{align*}
Alors $a = 12 \times 17 + 7 = 204 + 7 = 211$.

\textbf{Étape 3 : filtrage par la condition.} On doit avoir $0 \leq r < b$, c'est-à-dire $7 < 17$ : c'est vérifié. La solution est donc acceptable.

\textbf{Étape 4 : vérification.} $a + b = 211 + 17 = 228$ \checkmark{} et $211 = 17 \times 12 + 7$ avec $0 \leq 7 < 17$ \checkmark.

\textbf{Conclusion :} Le dividende est $\boxed{a = 211}$ et le diviseur est $\boxed{b = 17}$.
\end{exoresolu}

\courssub{Numération de base $a$ : principe et décomposition}

\begin{definition}[Système de numération de base $a$]
Soit $a \in \mathbb{N}$, $a \geq 2$. Tout entier naturel $N$ peut s'écrire de manière \cle{unique} sous la forme :
\[ N = x_n a^n + x_{n-1} a^{n-1} + \cdots + x_2 a^2 + x_1 a + x_0 \]
où les \cle{chiffres} $x_i$ vérifient $x_i \in \mathbb{N}$, $0 \leq x_i < a$ pour tout $i$, et $x_n \neq 0$ (sauf si $N=0$).
On note alors $N = \overline{x_n x_{n-1} \ldots x_1 x_0}^{\,a}$.
La base $10$ (usuelle) s'écrit sans exposant : $N = \overline{x_n \ldots x_0}$.
\end{definition}

\begin{aretenir}[Validité d'une écriture]
Dans une écriture en base $a$, \textbf{chaque chiffre doit être strictement inférieur à $a$}.
\begin{itemize}
\item En base $2$ (binaire) : chiffres autorisés $\{0, 1\}$.
\item En base $3$ : chiffres autorisés $\{0, 1, 2\}$.
\item En base $5$ : chiffres autorisés $\{0, 1, 2, 3, 4\}$.
\item \ldots
\item En base $10$ : chiffres autorisés $\{0, 1, \ldots, 9\}$.
\end{itemize}
Exemple : $\overline{253}^{\,5}$ est \textbf{invalide} (chiffre $5 \not< 5$) ; $\overline{1023}^{\,2}$ est \textbf{invalide} (chiffres $2, 3 \not< 2$).
\end{aretenir}

\begin{savaistu}[Le binaire, langage des machines]
Tout traitement informatique (téléphone, ordinateur, carte SIM, box internet CAMTEL) repose sur la base $2$. Un \cle{bit} (chiffre binaire) vaut $0$ ou $1$. Un octet ($8$ bits) permet de coder $2^8 = 256$ valeurs (caractères, couleurs, sons). Les adresses IP (ex: $192.168.1.1$) sont des nombres en base $2$ écrits en base $10$ par paquets de $8$ bits pour être lisibles par l'homme. Maîtriser le passage base $2 \leftrightarrow$ base $10$, c'est comprendre la langue secrète du numérique !
\end{savaistu}

\courssub{Écriture d'un entier en base $a$ : conversions avec la base 10}

\begin{methode}[Conversion base $a$ vers base 10 (par décomposition)]
Pour convertir en base 10 le nombre $N = \overline{x_n x_{n-1} \ldots x_1 x_0}^{\,a}$ :
\begin{enumerate}
\item \textbf{Numéroter} les chiffres de droite à gauche à partir du rang $0$ : $x_0$ (rang 0), $x_1$ (rang 1), \ldots, $x_n$ (rang $n$).
\item \textbf{Écrire la décomposition} en somme de puissances ordonnées de $a$ :
\[ N = x_n a^n + x_{n-1} a^{n-1} + \cdots + x_2 a^2 + x_1 a + x_0. \]
\item \textbf{Calculer} les puissances de $a$ ($a^0 = 1$, $a^1 = a$, $a^2$, $a^3$, \ldots), effectuer les produits puis la somme, \textbf{en base 10}.
\item \textbf{Conclure} : « $N = \ldots$ en base 10 ».
\end{enumerate}
\textbf{Contrôle préalable :} Vérifier que chaque chiffre $x_i$ satisfait $0 \leq x_i < a$. Sinon, l'écriture est invalide.
\end{methode}

\begin{exoresolu}[Conversion de la base 2 et de la base 5 vers la base 10]
\begin{enumerate}
\item Convertir en base 10 le nombre $N = \overline{101101}^{\,2}$.
\item L'écriture $\overline{3425}^{\,5}$ est-elle valide ? Si oui, convertir ce nombre en base 10.
\end{enumerate}
\tcblower
\textbf{1. Conversion de $N = \overline{101101}^{\,2}$.}

\textbf{Numérotation des rangs} (de droite à gauche, en partant de $0$) :
\[ N = 1 \times 2^5 + 0 \times 2^4 + 1 \times 2^3 + 1 \times 2^2 + 0 \times 2^1 + 1 \times 2^0. \]

\textbf{Calcul des puissances de 2 :} $2^0 = 1$, $2^1 = 2$, $2^2 = 4$, $2^3 = 8$, $2^4 = 16$, $2^5 = 32$.

\textbf{Somme :}
\begin{align*}
N &= 1 \times 32 + 0 \times 16 + 1 \times 8 + 1 \times 4 + 0 \times 2 + 1 \times 1 \\
&= 32 + 0 + 8 + 4 + 0 + 1 = 45.
\end{align*}
\textbf{Conclusion :} $\overline{101101}^{\,2} = 45$ en base 10.

\textbf{2. Étude de $\overline{3425}^{\,5}$.}

\textbf{Validité :} En base $5$, les seuls chiffres autorisés sont $0, 1, 2, 3, 4$. Or l'écriture contient le chiffre $5$, qui vérifie $5 \geq 5$. L'écriture $\overline{3425}^{\,5}$ est donc \textbf{invalide} : ce nombre n'existe pas en base $5$. Il n'y a rien à convertir.
\end{exoresolu}

\begin{methode}[Conversion base 10 vers base $a$ (divisions successives)]
Pour écrire en base $a$ un entier $N$ donné en base 10 :
\begin{enumerate}
\item \textbf{Effectuer} la division euclidienne de $N$ par $a$ : $N = a q_1 + r_0$. Le reste $r_0$ est le chiffre de rang $0$ (le plus à droite).
\item \textbf{Recommencer} avec le quotient : $q_1 = a q_2 + r_1$, puis $q_2 = a q_3 + r_2$, \ldots{} jusqu'à obtenir un quotient nul ($q_{k+1} = 0$).
\item \textbf{Lire les restes en sens inverse} de leur obtention (du dernier reste $r_k$ vers le premier $r_0$) :
\[ N = \overline{r_k \, r_{k-1} \ldots r_1 \, r_0}^{\,a}. \]
\item \textbf{Vérifier} en reconvertissant le résultat en base 10 (décomposition en puissances de $a$).
\end{enumerate}
\textbf{Disposition pratique (en cascade) :} On aligne les divisions verticalement, le quotient de chaque ligne devenant le dividende de la ligne suivante. Les restes se lisent de bas en haut.
\end{methode}

\begin{exoresolu}[Écrire 2024 en base 2 et 173 en base 3]
\begin{enumerate}
\item Écrire le nombre $2024$ en base $2$.
\item Écrire le nombre $173$ en base $3$, puis vérifier le résultat.
\end{enumerate}
\tcblower
\textbf{1. Écriture de 2024 en base 2.} Divisions euclidiennes successives par $2$ (disposition en cascade) :
\[
\begin{array}{r|l}
2024 & 2 \\
1012 & 2 \\
506  & 2 \\
253  & 2 \\
126  & 2 \\
63   & 2 \\
31   & 2 \\
15   & 2 \\
7    & 2 \\
3    & 2 \\
1    & 2 \\
0    &
\end{array}
\quad
\begin{array}{l}
\text{Restes} \\
0 \\
0 \\
0 \\
1 \\
0 \\
1 \\
1 \\
1 \\
1 \\
1 \\
1 \leftarrow \text{dernier reste (poids fort)}
\end{array}
\]
Lecture des restes \textbf{du dernier vers le premier} (de bas en haut) :
\[ 2024 = \overline{11111101000}^{\,2}. \]

\textbf{2. Écriture de 173 en base 3.} Divisions successives par $3$ :
\[
\begin{array}{r|l}
173 & 3 \\
57  & 3 \\
19  & 3 \\
6   & 3 \\
2   & 3 \\
0   &
\end{array}
\quad
\begin{array}{l}
\text{Restes} \\
2 \\
0 \\
1 \\
0 \\
2 \leftarrow \text{dernier reste}
\end{array}
\]
Lecture des restes de bas en haut : $173 = \overline{20102}^{\,3}$.

\textbf{Vérification} par décomposition en puissances de $3$ :
\begin{align*}
\overline{20102}^{\,3} &= 2 \times 3^4 + 0 \times 3^3 + 1 \times 3^2 + 0 \times 3 + 2 \\
&= 2 \times 81 + 0 + 9 + 0 + 2 = 162 + 11 = 173. \quad \checkmark
\end{align*}
\textbf{Conclusion :} $2024 = \overline{11111101000}^{\,2}$ et $173 = \overline{20102}^{\,3}$.
\end{exoresolu}

\courssub{Passage direct d'une base à l'autre}

\begin{methode}[Conversion directe d'une base $a$ à une base $b$ ($a \neq 10$, $b \neq 10$)]
\begin{enumerate}
\item \textbf{Étape 1 (Pivot base 10) :} Convertir le nombre de la base $a$ vers la \textbf{base 10} par décomposition en puissances de $a$ :
\[ N = x_n a^n + \cdots + x_1 a + x_0. \]
\item \textbf{Étape 2 :} Convertir le résultat (en base 10) vers la base $b$ par divisions euclidiennes successives par $b$.
\item \textbf{Étape 3 : Vérifier} en reconvertissant le résultat final en base 10 : on doit retrouver la même valeur qu'à l'étape 1.
\end{enumerate}
\textbf{Remarque :} Au programme de Première, on ne sait pas passer directement de la base $a$ à la base $b$ sans transiter par la base 10. La base 10 est le « pivot » obligatoire de toute conversion.
\end{methode}

\begin{exoresolu}[De la base 2 à la base 5]
Un technicien de CAMTEL relève sur un registre le nombre binaire $N = \overline{110101}^{\,2}$. Écrire $N$ en base $5$.
\tcblower
\textbf{Étape 1 : passage de la base 2 à la base 10.}
\begin{align*}
N &= 1 \times 2^5 + 1 \times 2^4 + 0 \times 2^3 + 1 \times 2^2 + 0 \times 2^1 + 1 \times 2^0 \\
&= 32 + 16 + 0 + 4 + 0 + 1 = 53.
\end{align*}
Donc $N = 53$ en base 10.

\textbf{Étape 2 : passage de la base 10 à la base 5.} Divisions successives par $5$ :
\[
\begin{array}{r|l}
53 & 5 \\
10 & 5 \\
2  & 5 \\
0  &
\end{array}
\quad
\begin{array}{l}
\text{Restes} \\
3 \\
0 \\
2 \leftarrow
\end{array}
\]
Lecture des restes de bas en haut : $53 = \overline{203}^{\,5}$.

\textbf{Étape 3 : vérification.}
\[ \overline{203}^{\,5} = 2 \times 5^2 + 0 \times 5 + 3 = 2 \times 25 + 0 + 3 = 53. \quad \checkmark \]

\textbf{Conclusion :} $\overline{110101}^{\,2} = \overline{203}^{\,5}$.
\end{exoresolu}

\courssub{Addition en base $a$}

\begin{methode}[Addition posée en base $a$]
Pour additionner deux nombres écrits en base $a$, on pose l'addition en colonnes \textbf{comme en base 10}, avec une seule modification : la \cle{retenue} se déclenche dès qu'une somme de colonne atteint ou dépasse $a$ (et non $10$).
\begin{enumerate}
\item Additionner les chiffres d'une colonne (plus la retenue éventuelle) : soit $S$ cette somme.
\item \textbf{Effectuer la division euclidienne de $S$ par $a$} : $S = a q + r$. On écrit $r$ dans la colonne et on retient $q$ pour la colonne suivante.
\item Recommencer de colonne en colonne, de droite à gauche.
\item \textbf{Vérifier} en convertissant les deux nombres et le résultat en base 10.
\end{enumerate}
\textbf{Phrase-clé :} « En base $a$, on retient $1$ dès que la colonne atteint $a$ ».
\end{methode}

\begin{exoresolu}[Addition en base 2 et en base 5]
\begin{enumerate}
\item Calculer en base $2$ : $\overline{1011}^{\,2} + \overline{110}^{\,2}$.
\item Calculer en base $5$ : $\overline{243}^{\,5} + \overline{134}^{\,5}$. Vérifier en base 10.
\end{enumerate}
\tcblower
\textbf{1. Addition en base 2.} On pose l'opération en colonnes (rappel : en base 2, $1+1 = \overline{10}^{\,2}$, c'est-à-dire on écrit $0$ et on retient $1$) :

\begin{center}
\begin{tabular}{ccccc}
 & \tiny{1} & \tiny{1} & \tiny{1} & \\
 & 1 & 0 & 1 & 1 \\
$+$ & & 1 & 1 & 0 \\ \hline
1 & 0 & 0 & 0 & 1
\end{tabular}
\end{center}

Détail colonne par colonne, de droite à gauche :
\begin{itemize}
\item Colonne 0 (unités) : $1 + 0 = 1$. On écrit $1$, retenue $0$.
\item Colonne 1 (deux) : $1 + 1 = 2 = 1 \times 2 + 0$. On écrit $0$, retenue $1$.
\item Colonne 2 (quatre) : $0 + 1 + 1 = 2 = 1 \times 2 + 0$. On écrit $0$, retenue $1$.
\item Colonne 3 (huit) : $1 + 0 + 1 = 2 = 1 \times 2 + 0$. On écrit $0$, retenue $1$.
\item Colonne 4 (seize) : la retenue $1$ tombe : on écrit $1$.
\end{itemize}
\textbf{Conclusion :} $\overline{1011}^{\,2} + \overline{110}^{\,2} = \overline{10001}^{\,2}$.

(Vérification : $\overline{1011}^{\,2} = 11$, $\overline{110}^{\,2} = 6$, $11+6=17$, et $\overline{10001}^{\,2} = 16+1=17$ \checkmark.)

\textbf{2. Addition en base 5.} On pose l'opération (en base 5, on retient $1$ dès qu'une colonne atteint $5$) :

\begin{center}
\begin{tabular}{cccc}
 & \tiny{1} & \tiny{1} & \\
 & 2 & 4 & 3 \\
$+$ & 1 & 3 & 4 \\ \hline
 & 4 & 3 & 2
\end{tabular}
\end{center}

Détail colonne par colonne :
\begin{itemize}
\item Colonne 0 : $3 + 4 = 7 = 1 \times 5 + 2$. On écrit $2$, retenue $1$.
\item Colonne 1 : $4 + 3 + 1 = 8 = 1 \times 5 + 3$. On écrit $3$, retenue $1$.
\item Colonne 2 : $2 + 1 + 1 = 4 = 0 \times 5 + 4$. On écrit $4$, retenue $0$.
\end{itemize}
Donc $\overline{243}^{\,5} + \overline{134}^{\,5} = \overline{432}^{\,5}$.

\textbf{Vérification en base 10 :}
\begin{align*}
\overline{243}^{\,5} &= 2 \times 25 + 4 \times 5 + 3 = 50 + 20 + 3 = 73, \\
\overline{134}^{\,5} &= 1 \times 25 + 3 \times 5 + 4 = 25 + 15 + 4 = 44, \\
\overline{432}^{\,5} &= 4 \times 25 + 3 \times 5 + 2 = 100 + 15 + 2 = 117.
\end{align*}
Or $73 + 44 = 117$ \checkmark. Le résultat est confirmé : $\overline{243}^{\,5} + \overline{134}^{\,5} = \overline{432}^{\,5}$.
\end{exoresolu}

\courssub{Multiplication en base $a$}

\begin{methode}[Multiplication posée en base $a$]
Pour multiplier deux nombres écrits en base $a$ :
\begin{enumerate}
\item \textbf{Poser} la multiplication comme en base 10 : multiplier le premier nombre par chaque chiffre du second, en décalant d'un rang à chaque ligne (produits partiels).
\item Dans chaque produit partiel, appliquer la règle de la retenue en base $a$ : si un produit de chiffres (augmenté de la retenue) vaut $P$, on effectue la division euclidienne $P = aq + r$, on écrit $r$ et on retient $q$.
\item \textbf{Additionner les produits partiels en base $a$} (avec les retenues en base $a$).
\item \textbf{Vérifier} en convertissant tout en base 10.
\end{enumerate}
\textbf{Astuce :} Construire mentalement la table de multiplication en base $a$ pour les petits produits ; en base $2$, elle est triviale ($0 \times x = 0$, $1 \times x = x$), ce qui rend les multiplications binaires très rapides.
\end{methode}

\begin{exoresolu}[Multiplication en base 2 et en base 3]
\begin{enumerate}
\item Calculer en base $2$ : $\overline{101}^{\,2} \times \overline{11}^{\,2}$.
\item Calculer en base $3$ : $\overline{21}^{\,3} \times \overline{2}^{\,3}$. Vérifier en base 10.
\end{enumerate}
\tcblower
\textbf{1. Multiplication en base 2.} On pose l'opération :

\begin{center}
\begin{tabular}{ccccc}
 & & 1 & 0 & 1 \\
$\times$ & & & 1 & 1 \\ \hline
 & & 1 & 0 & 1 \\
 & 1 & 0 & 1 & \\ \hline
 & 1 & 1 & 1 & 1
\end{tabular}
\end{center}

Détail :
\begin{itemize}
\item Produit par le chiffre des unités ($1$) : $101 \times 1 = 101$.
\item Produit par le chiffre de rang 1 ($1$) : $101 \times 1 = 101$, décalé d'un rang vers la gauche.
\item Addition des produits partiels en base 2 : $101 + 1010$ :
colonne 0 : $1+0=1$ ; colonne 1 : $0+1=1$ ; colonne 2 : $1+0=1$ ; colonne 3 : $0+1=1$. Résultat : $1111$.
\end{itemize}
\textbf{Conclusion :} $\overline{101}^{\,2} \times \overline{11}^{\,2} = \overline{1111}^{\,2}$.

(Vérification : $\overline{101}^{\,2} = 5$, $\overline{11}^{\,2} = 3$, $5 \times 3 = 15$, et $\overline{1111}^{\,2} = 8+4+2+1 = 15$ \checkmark.)

\textbf{2. Multiplication en base 3.} On calcule $\overline{21}^{\,3} \times 2$ chiffre par chiffre, de droite à gauche :
\begin{itemize}
\item $2 \times 1 = 2 = 0 \times 3 + 2$ : on écrit $2$, retenue $0$.
\item $2 \times 2 = 4 = 1 \times 3 + 1$ : on écrit $1$, retenue $1$.
\item La retenue $1$ tombe en tête : on écrit $1$.
\end{itemize}
Donc $\overline{21}^{\,3} \times \overline{2}^{\,3} = \overline{112}^{\,3}$.

\textbf{Vérification en base 10 :}
\[ \overline{21}^{\,3} = 2 \times 3 + 1 = 7 \quad ; \quad 7 \times 2 = 14 \quad ; \quad \overline{112}^{\,3} = 1 \times 9 + 1 \times 3 + 2 = 14. \quad \checkmark \]
\textbf{Conclusion :} $\overline{21}^{\,3} \times \overline{2}^{\,3} = \overline{112}^{\,3}$.
\end{exoresolu}

\begin{attention}[Les 5 erreurs qui coûtent des points au Probatoire]
\begin{enumerate}
\item \textbf{Oublier la condition $0 \leq r < b$.} Écrire $a = bq + r$ sans préciser que $0 \leq r < b$ rend l'affirmation « $q$ est le quotient et $r$ le reste » \textbf{fausse} : par exemple $50 = 7 \times 6 + 8$ est une égalité vraie, mais ce n'est \textbf{pas} la division euclidienne de $50$ par $7$ car $8 \geq 7$ (la bonne écriture est $50 = 7 \times 7 + 1$). Cite et vérifie toujours cette condition.
\item \textbf{Lire les restes dans le mauvais sens.} Dans la conversion de la base 10 vers la base $a$, les restes se lisent \textbf{du dernier obtenu vers le premier} (de bas en haut). Lire dans l'autre sens donne un nombre différent : sécurise-toi en reconvertissant ton résultat en base 10.
\item \textbf{Utiliser un chiffre interdit.} En base $a$, les chiffres appartiennent à $\{0, 1, \ldots, a-1\}$. Écrire $\overline{253}^{\,5}$ ou $\overline{1023}^{\,2}$ est une erreur d'existence, pas de calcul : vérifie toujours la validité de l'écriture \textbf{avant} de convertir.
\item \textbf{Garder les réflexes de la base 10 dans les opérations.} En base $a$, la retenue part dès que la colonne atteint $a$, pas $10$. Erreur fatale typique : écrire en base $2$ que $1 + 1 = 2$ (alors que $1 + 1 = \overline{10}^{\,2}$), ou en base $5$ que $3 + 4 = 7$ (alors que $3 + 4 = \overline{12}^{\,5}$).
\item \textbf{Passer d'une base $a$ à une base $b$ sans transiter par la base 10}, ou mélanger les bases dans un même calcul. Toute conversion entre deux bases différentes de 10 se fait en deux temps : base $a \to$ base 10 (décomposition en puissances) puis base 10 $\to$ base $b$ (divisions successives). Et n'oublie jamais d'\textbf{indiquer la base en exposant} : un nombre sans base précisée est ambigu et la copie est sanctionnée.
\end{enumerate}
\end{attention}

\courssub{Entraînement et évaluation}

\begin{exercice}[Exercices d'application]
\begin{enumerate}
\item Effectuer les divisions euclidiennes suivantes et écrire l'égalité correspondante :
      \begin{itemize}
      \item $3\,456$ par $17$
      \item $10\,000$ par $37$
      \end{itemize}
\item Dans une division euclidienne, le dividende est $589$, le quotient est $23$ et le reste est $r$. On sait que le diviseur $b$ est un entier à deux chiffres. Déterminer $b$ et $r$.
\item Convertir en base 10 :
      \begin{itemize}
      \item $\overline{1101101}^{\,2}$
      \item $\overline{4031}^{\,5}$
      \item $\overline{1202}^{\,3}$ (vérifier la validité d'abord)
      \end{itemize}
\item Écrire les nombres suivants dans la base indiquée :
      \begin{itemize}
      \item $2025$ en base $2$
      \item $500$ en base $6$
      \item $1\,000$ en base $8$
      \end{itemize}
\item Effectuer les opérations suivantes \textbf{sans passer par la base 10} (vérifier ensuite en base 10) :
      \begin{itemize}
      \item $\overline{1101}^{\,2} + \overline{1011}^{\,2}$
      \item $\overline{324}^{\,5} + \overline{231}^{\,5}$
      \item $\overline{101}^{\,2} \times \overline{101}^{\,2}$
      \item $\overline{22}^{\,3} \times \overline{12}^{\,3}$
      \end{itemize}
\item Passer directement :
      \begin{itemize}
      \item $\overline{101101}^{\,2}$ en base $4$
      \item $\overline{3102}^{\,4}$ en base $2$
      \end{itemize}
\end{enumerate}
\end{exercice}

\begin{corrige}[Exercices d'application]
\begin{enumerate}
\item \textbf{Divisions euclidiennes :}
      \begin{itemize}
      \item $3\,456 = 17 \times 203 + 5$ ($q=203, r=5$). Vérif : $17\times203=3451$, $3451+5=3456$.
      \item $10\,000 = 37 \times 270 + 10$ ($q=270, r=10$). Vérif : $37\times270=9990$, $9990+10=10000$.
      \end{itemize}
\item \textbf{Diviseur et reste inconnus :} $589 = 23b + r$ avec $0 \leq r < b$ et $10 \leq b \leq 99$.
      $23b \leq 589 < 23(b+1) \Rightarrow b \leq 25,6 \Rightarrow b \leq 25$. Aussi $589 - 23b < b \Rightarrow 589 < 24b \Rightarrow b > 24,5 \Rightarrow b \geq 25$.
      Donc $b = 25$. Alors $r = 589 - 23\times25 = 589 - 575 = 14$. Vérif : $14 < 25$ \checkmark. $\boxed{b=25, r=14}$.
\item \textbf{Conversions vers base 10 :}
      \begin{itemize}
      \item $\overline{1101101}^{\,2} = 1\times64 + 1\times32 + 0\times16 + 1\times8 + 1\times4 + 0\times2 + 1 = 109$.
      \item $\overline{4031}^{\,5} = 4\times125 + 0\times25 + 3\times5 + 1 = 500+0+15+1 = 516$.
      \item $\overline{1202}^{\,3}$ : chiffres $\{1,2,0,2\}$ tous $<3$, valide. $= 1\times27 + 2\times9 + 0\times3 + 2 = 27+18+2 = 47$.
      \end{itemize}
\item \textbf{Écriture en base $a$ :}
      \begin{itemize}
      \item $2025$ en base $2$ : divisions par 2 $\to$ $\overline{11111101001}^{\,2}$.
      \item $500$ en base $6$ : $500 = 6\times83+2$, $83=6\times13+5$, $13=6\times2+1$, $2=6\times0+2$ $\to$ $\overline{2152}^{\,6}$.
      \item $1000$ en base $8$ : $1000 = 8\times125+0$, $125=8\times15+5$, $15=8\times1+7$, $1=8\times0+1$ $\to$ $\overline{1750}^{\,8}$.
      \end{itemize}
\item \textbf{Opérations en base $a$ :}
      \begin{itemize}
      \item $\overline{1101}^{\,2} + \overline{1011}^{\,2} = \overline{11000}^{\,2}$ (Vérif: $13+11=24$, $\overline{11000}^{\,2}=16+8=24$).
      \item $\overline{324}^{\,5} + \overline{231}^{\,5}$ : Col 0: $4+1=5=\overline{10}^{\,5}$ (0, ret 1). Col 1: $2+3+1=6=\overline{11}^{\,5}$ (1, ret 1). Col 2: $3+2+1=6=\overline{11}^{\,5}$ (1, ret 1). Col 3: ret 1. Rés: $\overline{1110}^{\,5}$. (Vérif: $89+66=155$, $\overline{1110}^{\,5}=125+25+5=155$).
      \item $\overline{101}^{\,2} \times \overline{101}^{\,2} = \overline{11001}^{\,2}$ ($5\times5=25$, $\overline{11001}^{\,2}=16+8+1=25$).
      \item $\overline{22}^{\,3} \times \overline{12}^{\,3}$ : $\overline{22}^{\,3}=8$, $\overline{12}^{\,3}=5$, $8\times5=40$. En base 3: $22\times2 = \overline{121}^{\,3}$ (ret), $22\times1$ (décalé) $= \overline{220}^{\,3}$. Somme: $\overline{121}+\overline{220}=\overline{1111}^{\,3}$ ($1\times27+1\times9+1\times3+1=40$).
      \end{itemize}
\item \textbf{Passage direct :}
      \begin{itemize}
      \item $\overline{101101}^{\,2} = 45_{10}$. $45$ en base $4$: $45=4\times11+1$, $11=4\times2+3$, $2=4\times0+2$ $\to$ $\overline{231}^{\,4}$.
      \item $\overline{3102}^{\,4} = 3\times64 + 1\times16 + 0\times4 + 2 = 210_{10}$. $210$ en base $2$: $210=2\times105+0$, $105=2\times52+1$, $52=2\times26+0$, $26=2\times13+0$, $13=2\times6+1$, $6=2\times3+0$, $3=2\times1+1$, $1=2\times0+1$ $\to$ $\overline{11010010}^{\,2}$.
      \end{itemize}
\end{enumerate}
\end{corrige}

\newpage

\coursec{Exercices}

\courssub{Je vérifie mes connaissances}

\begin{exercice}[QCM : division euclidienne]
Pour chaque question, une seule réponse est exacte. Entoure la bonne lettre.
\begin{enumerate}
\item La division euclidienne de $4573$ par $37$ donne :
\begin{enumerate}
\item $q = 123$ et $r = 22$ ;
\item $q = 124$ et $r = 25$ ;
\item $q = 123$ et $r = 59$ ;
\item $q = 122$ et $r = 59$.
\end{enumerate}
\item Dans la relation $a = bq + r$ de la division euclidienne de $a$ par $b$ ($b \neq 0$), le reste $r$ vérifie :
\begin{enumerate}
\item $0 \le r \le b$ ;
\item $0 < r < b$ ;
\item $0 \le r < b$ ;
\item $0 \le r < q$.
\end{enumerate}
\item L'écriture $(1101)_2$ désigne en base $10$ le nombre :
\begin{enumerate}
\item $11$ ;
\item $13$ ;
\item $14$ ;
\item $1101$.
\end{enumerate}
\item Le nombre $(234)_5$ est égal en base $10$ à :
\begin{enumerate}
\item $69$ ;
\item $67$ ;
\item $234$ ;
\item $59$.
\end{enumerate}
\end{enumerate}
\end{exercice}

\begin{exercice}[Vrai ou faux, justifié]
Réponds par VRAI ou FAUX en justifiant chaque réponse par un calcul ou un contre-exemple.
\begin{enumerate}
\item « Dans la division euclidienne de $515$ par $24$, le quotient est $21$ et le reste est $11$. »
\item « L'écriture $(1023)_4$ est une écriture correcte en base $4$. »
\item « L'écriture $(258)_7$ est une écriture correcte en base $7$. »
\item « Le nombre $(1000)_2$ est égal à $8$ en base $10$. »
\item « En base $2$, on a $(1)_2 + (1)_2 = (2)_2$. »
\end{enumerate}
\end{exercice}

\begin{exercice}[Définitions et vocabulaire]
\begin{enumerate}
\item Énonce le théorème de la division euclidienne dans $\mathbb{N}$, en précisant les conditions portant sur chaque entier intervenant dans la relation.
\item Qu'appelle-t-on \cle{base} d'un système de numération ? Cite les chiffres utilisés en base $2$, en base $5$ et en base $8$.
\item Pourquoi l'écriture $(312)_3$ est-elle impossible ? Justifie.
\item Complète : dans l'écriture $N = \overline{x_n x_{n-1} \ldots x_1 x_0}^{\,a}$, les chiffres $x_i$ vérifient la condition \ldots\ et le chiffre $x_n$ vérifie de plus \ldots
\end{enumerate}
\end{exercice}

\begin{exercice}[Calculs directs]
\begin{enumerate}
\item Détermine le quotient et le reste de la division euclidienne de $4573$ par $37$, puis écris la relation $a = bq + r$ et vérifie que la condition sur le reste est satisfaite.
\item Détermine le quotient et le reste de la division euclidienne de $1000$ par $9$.
\item Convertis en base $10$ : $(110101)_2$ ; $(2102)_3$ ; $(475)_8$.
\item Écris en base $2$ les nombres $7$, $16$ et $31$.
\end{enumerate}
\end{exercice}

\courssub{Je m'entraîne}

\begin{exercice}[Conversions vers la base 2, 3, 5 et 8]
On considère l'entier $N = 156$.
\begin{enumerate}
\item Par la méthode des divisions successives, écris $156$ en base $2$.
\item Écris $156$ en base $3$, puis en base $5$, puis en base $8$.
\item Vérifie chaque résultat en développant l'écriture obtenue en somme de puissances ordonnées de la base, et en retrouvant $156$.
\end{enumerate}
\end{exercice}

\begin{exercice}[Passage direct de la base 2 à la base 8]
\begin{enumerate}
\item Convertis $(101110)_2$ en base $10$.
\item Convertis $(101110)_2$ directement en base $8$ par regroupement des bits en paquets de trois, en partant de la droite.
\item Retrouve le résultat de la question 2 en convertissant en base $8$ le nombre obtenu à la question 1.
\item Par la même méthode de regroupement, écris $(1101011)_2$ en base $8$.
\end{enumerate}
\end{exercice}

\begin{exercice}[Additions et multiplications en base 2]
\begin{enumerate}
\item Pose et effectue en base $2$ l'addition $(1101)_2 + (111)_2$.
\item Pose et effectue en base $2$ la multiplication $(1011)_2 \times (101)_2$.
\item Vérifie chacun des deux résultats en convertissant tous les nombres en base $10$ et en y effectuant les opérations.
\end{enumerate}
\end{exercice}

\begin{exercice}[Additions et multiplications en base 8]
\begin{enumerate}
\item Pose et effectue en base $8$ l'addition $(356)_8 + (477)_8$.
\item Pose et effectue en base $8$ la multiplication $(25)_8 \times (6)_8$.
\item Vérifie chacun des deux résultats par un passage en base $10$.
\end{enumerate}
\end{exercice}

\courssub{Je me prépare au Probatoire}

\begin{exercice}[Les roses de la fête des mères]
Pour la fête des mères, un fleuriste du marché central de Yaoundé dispose de $515$ roses. Il les conditionne en bouquets de $24$ roses chacun.
\begin{enumerate}
\item En effectuant la division euclidienne de $515$ par $24$, détermine le nombre de bouquets complets qu'il peut confectionner et le nombre de roses restantes.
\item Écris la relation $a = bq + r$ correspondant à cette situation et vérifie que $0 \le r < 24$.
\item Le fleuriste vend chaque bouquet à $\SI{2500}{F}$ et chaque rose restante à l'unité à $\SI{150}{F}$. Calcule la recette totale qu'il obtient s'il vend toutes ses roses.
\item Il souhaite n'avoir aucune rose restante. Quel est le plus petit nombre de roses supplémentaires qu'il doit acheter pour pouvoir confectionner un bouquet de plus ?
\end{enumerate}
\end{exercice}

\begin{exercice}[Le code mystérieux]
Un technicien de CAMTEL relève sur un ancien document un nombre noté $(213)_a$, écrit dans une base $a$ inconnue ($a \in \mathbb{N}$, $a \ge 2$). Il sait par ailleurs que ce nombre vaut $58$ en base $10$.
\begin{enumerate}
\item Justifie que la base $a$ vérifie nécessairement $a \ge 4$.
\item Exprime $(213)_a$ en fonction de $a$ et montre que $a$ est solution de l'équation $2a^2 + a + 3 = 58$.
\item Résous dans $\mathbb{R}$ l'équation $2a^2 + a - 55 = 0$.
\item Déduis-en la valeur de la base $a$, en tenant compte de la question 1.
\item Vérifie ta réponse en convertissant $(213)_a$ en base $10$ avec la valeur de $a$ trouvée.
\item Écris le nombre $58$ en base $2$.
\end{enumerate}
\end{exercice}

\begin{exercice}[Divisibilité et dernier chiffre]
On rappelle que tout entier naturel $N$ écrit en base $a$ se décompose sous la forme
\[ N = x_n a^n + x_{n-1} a^{n-1} + \cdots + x_1 a + x_0 \quad \text{avec } 0 \le x_i < a. \]
\textbf{Partie A : le cas de la base 2}
\begin{enumerate}
\item Écris en base $2$ les nombres $12$ et $13$. Que remarques-tu concernant le dernier chiffre ?
\item Soit $N$ un entier dont l'écriture binaire se termine par $0$. En utilisant la décomposition rappelée ci-dessus avec $a = 2$, montre que $N$ est pair.
\item Réciproquement, montre que si l'écriture binaire de $N$ se termine par $1$, alors $N$ est impair.
\end{enumerate}
\textbf{Partie B : généralisation}
\begin{enumerate}
\item Soit $N$ un entier écrit en base $a$. En factorisant $a$ dans la décomposition de $N$, montre que $N = aQ + x_0$ où $Q$ est un entier naturel que l'on précisera.
\item Déduis-en que $N$ est divisible par $a$ si et seulement si son dernier chiffre $x_0$ est nul.
\item Application : sans les convertir en base $10$, indique lesquels des nombres suivants sont divisibles par leur base : $(11010)_2$ ; $(234)_5$ ; $(560)_8$ ; $(123)_4$.
\end{enumerate}
\end{exercice}

\newpage

\coursec{Corrigés des exercices}

\begin{corrige}[Exercice 1 — QCM : division euclidienne]
\begin{enumerate}
\item \textbf{Réponse a.} On calcule : $37 \times 123 = 4551$ et $4573 - 4551 = 22$. Donc $4573 = 37 \times 123 + 22$ avec $0 \le 22 < 37$ : la condition sur le reste est satisfaite. Vérifions que les autres propositions échouent : pour b, $37 \times 124 + 25 = 4588 + 25 = 4613 \neq 4573$ ; pour c et d, le reste $59$ ne vérifie pas la condition $r < 37$ (car $59 \ge 37$).
\item \textbf{Réponse c.} Le reste vérifie toujours $0 \le r < b$ : il peut être nul (division exacte) mais doit être \emph{strictement} inférieur au diviseur.
\item \textbf{Réponse b.} $(1101)_2 = 1\times 2^3 + 1 \times 2^2 + 0 \times 2 + 1 = 8 + 4 + 1 = 13$.
\item \textbf{Réponse a.} $(234)_5 = 2 \times 5^2 + 3 \times 5 + 4 = 50 + 15 + 4 = 69$.
\end{enumerate}
\end{corrige}

\begin{corrige}[Exercice 2 — Vrai ou faux, justifié]
\begin{enumerate}
\item \textbf{VRAI.} $24 \times 21 = 504$ et $515 - 504 = 11$, donc $515 = 24 \times 21 + 11$ avec $0 \le 11 < 24$. Le quotient est bien $21$ et le reste $11$.
\item \textbf{VRAI.} En base $4$, les chiffres autorisés sont $0, 1, 2, 3$. Les chiffres de $(1023)_4$ sont $1, 0, 2, 3$, tous strictement inférieurs à $4$, et le premier chiffre est non nul : l'écriture est correcte.
\item \textbf{FAUX.} En base $7$, les chiffres autorisés sont $0, 1, 2, 3, 4, 5, 6$. Or l'écriture $(258)_7$ contient le chiffre $8$, qui vérifie $8 \ge 7$ : elle est donc impossible.
\item \textbf{VRAI.} $(1000)_2 = 1 \times 2^3 + 0 \times 2^2 + 0 \times 2 + 0 = 8$.
\item \textbf{FAUX.} Le chiffre $2$ n'existe pas en base $2$. En posant l'addition : $1 + 1 = 2 = 1 \times 2 + 0$, on écrit $0$ et on retient $1$, d'où $(1)_2 + (1)_2 = (10)_2$.
\end{enumerate}
\end{corrige}

\begin{corrige}[Exercice 3 — Définitions et vocabulaire]
\begin{enumerate}
\item \textbf{Théorème (division euclidienne dans $\mathbb{N}$).} Soient $a \in \mathbb{N}$ et $b \in \mathbb{N}$ avec $b \neq 0$. Il existe un \emph{unique} couple d'entiers naturels $(q, r)$ tel que
\[ a = bq + r \quad \text{avec} \quad 0 \le r < b. \]
L'entier $q$ s'appelle le \cle{quotient} et $r$ le \cle{reste} de la division euclidienne de $a$ par $b$.
\item La \cle{base} d'un système de numération est l'entier naturel $a \ge 2$ dont les puissances successives servent à décomposer tout entier ; c'est aussi le nombre de chiffres distincts utilisés. En base $2$ : chiffres $0, 1$ ; en base $5$ : chiffres $0, 1, 2, 3, 4$ ; en base $8$ : chiffres $0, 1, 2, 3, 4, 5, 6, 7$.
\item En base $3$, les seuls chiffres autorisés sont $0, 1, 2$. L'écriture $(312)_3$ contient le chiffre $3$, qui vérifie $3 \ge 3$ : elle est donc impossible.
\item Les chiffres $x_i$ vérifient la condition $0 \le x_i < a$ (pour tout $i$), et le chiffre $x_n$ vérifie de plus $x_n \neq 0$ (le premier chiffre d'une écriture n'est jamais nul).
\end{enumerate}
\end{corrige}

\begin{corrige}[Exercice 4 — Calculs directs]
\begin{enumerate}
\item On pose la division : $37 \times 100 = 3700$ ; $4573 - 3700 = 873$ ; $37 \times 20 = 740$ ; $873 - 740 = 133$ ; $37 \times 3 = 111$ ; $133 - 111 = 22$. Donc $q = 100 + 20 + 3 = 123$ et $r = 22$. La relation s'écrit :
\[ 4573 = 37 \times 123 + 22, \quad \text{et} \quad 0 \le 22 < 37 \ \checkmark \]
\item $9 \times 111 = 999$ et $1000 - 999 = 1$, donc $1000 = 9 \times 111 + 1$ avec $0 \le 1 < 9$. Quotient : $111$ ; reste : $1$.
\item Conversions en base $10$ :
\begin{align*}
(110101)_2 &= 2^5 + 2^4 + 2^2 + 1 = 32 + 16 + 4 + 1 = 53 ;\\
(2102)_3 &= 2 \times 3^3 + 1 \times 3^2 + 0 \times 3 + 2 = 54 + 9 + 2 = 65 ;\\
(475)_8 &= 4 \times 8^2 + 7 \times 8 + 5 = 256 + 56 + 5 = 317.
\end{align*}
\item Écritures en base $2$ : $7 = 4 + 2 + 1 = (111)_2$ ; $16 = 2^4 = (10000)_2$ ; $31 = 16 + 8 + 4 + 2 + 1 = (11111)_2$.
\end{enumerate}
\end{corrige}

\begin{corrige}[Exercice 5 — Conversions vers la base 2, 3, 5 et 8]
\begin{enumerate}
\item \textbf{Base 2} — divisions successives par $2$ :
\begin{center}
\begin{tabular}{|c|c|c|}
\hline
\rowcolor{popblueL} Dividende & Quotient & Reste \\
\hline
$156$ & $78$ & $0$ \\
$78$ & $39$ & $0$ \\
$39$ & $19$ & $1$ \\
$19$ & $9$ & $1$ \\
$9$ & $4$ & $1$ \\
$4$ & $2$ & $0$ \\
$2$ & $1$ & $0$ \\
$1$ & $0$ & $1$ \\
\hline
\end{tabular}
\end{center}
On lit les restes de bas en haut : $156 = (10011100)_2$.
\item \textbf{Base 3} : $156 = 3 \times 52 + 0$ ; $52 = 3 \times 17 + 1$ ; $17 = 3 \times 5 + 2$ ; $5 = 3 \times 1 + 2$ ; $1 = 3 \times 0 + 1$. Lecture de bas en haut : $156 = (12210)_3$.

\textbf{Base 5} : $156 = 5 \times 31 + 1$ ; $31 = 5 \times 6 + 1$ ; $6 = 5 \times 1 + 1$ ; $1 = 5 \times 0 + 1$. D'où $156 = (1111)_5$.

\textbf{Base 8} : $156 = 8 \times 19 + 4$ ; $19 = 8 \times 2 + 3$ ; $2 = 8 \times 0 + 2$. D'où $156 = (234)_8$.
\item \textbf{Vérifications} par développement en puissances ordonnées :
\begin{align*}
(10011100)_2 &= 2^7 + 2^4 + 2^3 + 2^2 = 128 + 16 + 8 + 4 = 156\ \checkmark \\
(12210)_3 &= 3^4 + 2 \times 3^3 + 2 \times 3^2 + 1 \times 3 + 0 = 81 + 54 + 18 + 3 = 156\ \checkmark \\
(1111)_5 &= 5^3 + 5^2 + 5 + 1 = 125 + 25 + 5 + 1 = 156\ \checkmark \\
(234)_8 &= 2 \times 8^2 + 3 \times 8 + 4 = 128 + 24 + 4 = 156\ \checkmark
\end{align*}
\end{enumerate}
\end{corrige}

\begin{corrige}[Exercice 6 — Passage direct de la base 2 à la base 8]
\begin{enumerate}
\item $(101110)_2 = 2^5 + 2^3 + 2^2 + 2 = 32 + 8 + 4 + 2 = 46$.
\item On regroupe les bits par paquets de trois en partant de la droite : $(101110)_2 = \underbrace{101}_{5}\ \underbrace{110}_{6}$. Or $(101)_2 = 5$ et $(110)_2 = 6$, donc $(101110)_2 = (56)_8$.
\item Vérification : $46 = 8 \times 5 + 6$, donc $46 = (56)_8$. On retrouve bien le résultat de la question 2. (Contrôle : $(56)_8 = 5 \times 8 + 6 = 46$ \checkmark)
\item On regroupe : $(1101011)_2 = \underbrace{001}_{1}\ \underbrace{101}_{5}\ \underbrace{011}_{3}$ (on complète le paquet de gauche par des zéros pour obtenir trois chiffres). Donc $(1101011)_2 = (153)_8$. Contrôle : $(153)_8 = 64 + 40 + 3 = 107$ et $(1101011)_2 = 64 + 32 + 8 + 2 + 1 = 107$ \checkmark
\end{enumerate}
\textbf{Pourquoi ça marche :} $8 = 2^3$, donc chaque chiffre octal correspond exactement à un bloc de trois chiffres binaires.
\end{corrige}

\begin{corrige}[Exercice 7 — Additions et multiplications en base 2]
\begin{enumerate}
\item Addition posée (retenues sur la ligne du haut) :
\begin{center}
\begin{tabular}{ccccc}
 & \small 1 & \small 1 & \small 1 & \\
 & 1 & 1 & 0 & 1 \\
$+$ & 0 & 1 & 1 & 1 \\
\hline
1 & 0 & 1 & 0 & 0 \\
\end{tabular}
\end{center}
Détail colonne par colonne, de droite à gauche : $1+1 = 2 = (10)_2$ : on écrit $0$, retenue $1$ ; $0+1+1 = 2$ : on écrit $0$, retenue $1$ ; $1+1+1 = 3 = (11)_2$ : on écrit $1$, retenue $1$ ; $1+0+1 = 2$ : on écrit $0$, retenue $1$ que l'on abaisse. Donc $(1101)_2 + (111)_2 = (10100)_2$.
\item Multiplication posée :
\begin{center}
\begin{tabular}{ccccccc}
 & & & 1 & 0 & 1 & 1 \\
$\times$ & & & & 1 & 0 & 1 \\
\hline
 & & & 1 & 0 & 1 & 1 \\
 & & 0 & 0 & 0 & 0 & \\
$+$ & 1 & 0 & 1 & 1 & & \\
\hline
 & 1 & 1 & 0 & 1 & 1 & 1 \\
\end{tabular}
\end{center}
Donc $(1011)_2 \times (101)_2 = (110111)_2$.
\item \textbf{Vérifications en base 10} : $(1101)_2 = 13$, $(111)_2 = 7$, et $13 + 7 = 20 = 16 + 4 = (10100)_2$ \checkmark ; $(1011)_2 = 11$, $(101)_2 = 5$, et $11 \times 5 = 55 = 32 + 16 + 4 + 2 + 1 = (110111)_2$ \checkmark
\end{enumerate}
\end{corrige}

\begin{corrige}[Exercice 8 — Additions et multiplications en base 8]
\begin{enumerate}
\item Addition posée, colonne par colonne (de droite à gauche) :
\begin{itemize}
\item Unités : $6 + 7 = 13 = 1 \times 8 + 5$ : on écrit $5$, retenue $1$ ;
\item « Huitaines » : $5 + 7 + 1 = 13 = 1 \times 8 + 5$ : on écrit $5$, retenue $1$ ;
\item « Soixante-quatraines » : $3 + 4 + 1 = 8 = 1 \times 8 + 0$ : on écrit $0$, retenue $1$ que l'on abaisse.
\end{itemize}
Donc $(356)_8 + (477)_8 = (1055)_8$.
\item Multiplication : $5 \times 6 = 30 = 3 \times 8 + 6$ : on écrit $6$, retenue $3$ ; puis $2 \times 6 + 3 = 15 = 1 \times 8 + 7$ : on écrit $7$, retenue $1$ que l'on abaisse. Donc $(25)_8 \times (6)_8 = (176)_8$.
\item \textbf{Vérifications en base 10} : $(356)_8 = 3 \times 64 + 5 \times 8 + 6 = 238$ et $(477)_8 = 4 \times 64 + 7 \times 8 + 7 = 319$ ; $238 + 319 = 557$ et $(1055)_8 = 512 + 5 \times 8 + 5 = 512 + 40 + 5 = 557$ \checkmark. Ensuite $(25)_8 = 21$, $21 \times 6 = 126$ et $(176)_8 = 64 + 7 \times 8 + 6 = 64 + 56 + 6 = 126$ \checkmark.
\end{enumerate}
\end{corrige}

\begin{corrige}[Exercice 9 — Les roses de la fête des mères]
\begin{enumerate}
\item On effectue la division euclidienne : $24 \times 21 = 504$ et $515 - 504 = 11$, donc
\[ 515 = 24 \times 21 + 11. \]
Le fleuriste peut confectionner \textbf{21 bouquets complets} et il restera \textbf{11 roses}.
\item La relation $a = bq + r$ s'écrit $515 = 24 \times 21 + 11$, avec $a = 515$, $b = 24$, $q = 21$, $r = 11$. On vérifie : $0 \le 11 < 24$ \checkmark.
\item Recette totale :
\[ 21 \times \SI{2500}{F} + 11 \times \SI{150}{F} = \SI{52500}{F} + \SI{1650}{F} = \SI{54150}{F}. \]
\item Pour un bouquet supplémentaire, il faut compléter les $11$ roses restantes jusqu'à $24$ : il manque $24 - 11 = 13$ roses. Le plus petit nombre de roses à acheter est donc $\mathbf{13}$ (il obtiendra alors $22$ bouquets et $0$ rose restante : $515 + 13 = 528 = 24 \times 22$).
\end{enumerate}
\textbf{Barème indicatif (sur 4 pts) :} division euclidienne correcte et interprétation ($1$ pt) ; relation $a = bq + r$ avec vérification de $0 \le r < 24$ ($1$ pt) ; recette $\SI{54150}{F}$ ($1$ pt) ; réponse $13$ roses justifiée ($1$ pt).

\textbf{Points de vigilance :} le nombre de bouquets est le \emph{quotient}, pas le quotient arrondi ; citer explicitement la condition $0 \le r < b$ ; à la question 4, c'est $b - r$ qu'il faut calculer, et non $b$.
\end{corrige}

\begin{corrige}[Exercice 10 — Le code mystérieux]
\begin{enumerate}
\item L'écriture $(213)_a$ contient le chiffre $3$. Or dans une base $a$, tout chiffre $x$ vérifie $0 \le x < a$. Il faut donc $3 < a$, c'est-à-dire $a \ge 4$.
\item Par définition de l'écriture en base $a$ :
\[ (213)_a = 2a^2 + 1 \times a + 3 = 2a^2 + a + 3. \]
Comme ce nombre vaut $58$ en base $10$, $a$ vérifie $2a^2 + a + 3 = 58$.
\item L'équation $2a^2 + a + 3 = 58$ équivaut à $2a^2 + a - 55 = 0$. Discriminant :
\[ \Delta = 1^2 - 4 \times 2 \times (-55) = 1 + 440 = 441 = 21^2. \]
$\Delta > 0$, donc deux solutions réelles :
\[ a_1 = \frac{-1 - 21}{4} = \frac{-22}{4} = -\frac{11}{2} = -5{,}5 \quad ; \quad a_2 = \frac{-1 + 21}{4} = \frac{20}{4} = 5. \]
\item Une base est un entier naturel supérieur ou égal à $2$, et la question 1 impose $a \ge 4$. La solution $-5{,}5$ est à rejeter (négative et non entière). Donc $\mathbf{a = 5}$.
\item Vérification : $(213)_5 = 2 \times 25 + 1 \times 5 + 3 = 50 + 5 + 3 = 58$ \checkmark.
\item Divisions successives de $58$ par $2$ : $58 = 2 \times 29 + 0$ ; $29 = 2 \times 14 + 1$ ; $14 = 2 \times 7 + 0$ ; $7 = 2 \times 3 + 1$ ; $3 = 2 \times 1 + 1$ ; $1 = 2 \times 0 + 1$. Lecture des restes de bas en haut : $58 = (111010)_2$. Contrôle : $32 + 16 + 8 + 2 = 58$ \checkmark.
\end{enumerate}
\textbf{Barème indicatif (sur 6 pts) :} justification $a \ge 4$ ($0{,}5$ pt) ; mise en équation ($1$ pt) ; discriminant et résolution ($1{,}5$ pt) ; choix de $a = 5$ avec rejet justifié ($1$ pt) ; vérification ($1$ pt) ; conversion de $58$ en base $2$ ($1$ pt).

\textbf{Points de vigilance :} la condition « chiffre $<$ base » doit être écrite explicitement ; ne pas oublier de rejeter la racine négative \emph{en justifiant} (une base est un entier $\ge 2$) ; la vérification finale est une étape rédactionnelle attendue au Probatoire.
\end{corrige}

\begin{corrige}[Exercice 11 — Divisibilité et dernier chiffre]
\textbf{Partie A : le cas de la base 2}
\begin{enumerate}
\item $12 = 8 + 4 = (1100)_2$ et $13 = 8 + 4 + 1 = (1101)_2$. On remarque que $12$, qui est pair, se termine par $0$, tandis que $13$, qui est impair, se termine par $1$.
\item Soit $N$ dont l'écriture binaire se termine par $0$, c'est-à-dire $x_0 = 0$. Alors :
\[ N = x_n 2^n + x_{n-1} 2^{n-1} + \cdots + x_1 \times 2 + 0 = 2\underbrace{\left(x_n 2^{n-1} + x_{n-1} 2^{n-2} + \cdots + x_1\right)}_{=\,k\ \in\ \mathbb{N}} = 2k. \]
Donc $N$ est un multiple de $2$ : $N$ est pair.
\item Si l'écriture binaire de $N$ se termine par $1$, alors $x_0 = 1$ et, avec le même entier $k$ que ci-dessus :
\[ N = 2\left(x_n 2^{n-1} + \cdots + x_1\right) + 1 = 2k + 1. \]
Le reste de la division euclidienne de $N$ par $2$ est $1$ (et $0 \le 1 < 2$), donc $N$ est impair.
\end{enumerate}
\textbf{Partie B : généralisation}
\begin{enumerate}
\item On factorise $a$ dans tous les termes sauf le dernier :
\[ N = x_n a^n + x_{n-1} a^{n-1} + \cdots + x_1 a + x_0 = a\underbrace{\left(x_n a^{n-1} + x_{n-1} a^{n-2} + \cdots + x_1\right)}_{=\,Q} + x_0, \]
avec $Q = x_n a^{n-1} + x_{n-1} a^{n-2} + \cdots + x_1 \in \mathbb{N}$. Comme $0 \le x_0 < a$, l'égalité $N = aQ + x_0$ est exactement la division euclidienne de $N$ par $a$ : $x_0$ en est le reste.
\item $N$ est divisible par $a$ si et seulement si le reste de sa division euclidienne par $a$ est nul, c'est-à-dire si et seulement si $x_0 = 0$ : le dernier chiffre de son écriture en base $a$ est nul.
\item Application directe du critère (on regarde uniquement le dernier chiffre) :
\begin{itemize}
\item $(11010)_2$ se termine par $0$ : \textbf{divisible} par $2$ ;
\item $(234)_5$ se termine par $4 \neq 0$ : \textbf{non divisible} par $5$ ;
\item $(560)_8$ se termine par $0$ : \textbf{divisible} par $8$ ;
\item $(123)_4$ se termine par $3 \neq 0$ : \textbf{non divisible} par $4$.
\end{itemize}
\end{enumerate}
\textbf{Barème indicatif (sur 6 pts) :} conversions et remarque ($1$ pt) ; preuve « se termine par $0 \Rightarrow$ pair » ($1$ pt) ; preuve du cas impair ($1$ pt) ; factorisation $N = aQ + x_0$ avec $Q$ précisé ($1{,}5$ pt) ; équivalence de divisibilité ($1$ pt) ; quatre applications correctes ($0{,}5$ pt).

\textbf{Points de vigilance :} dans les démonstrations, il faut \emph{factoriser} proprement et nommer l'entier obtenu ($k$ ou $Q$) ; pour l'équivalence, citer le fait que $N = aQ + x_0$ avec $0 \le x_0 < a$ est la division euclidienne de $N$ par $a$, donc $x_0$ est le reste ; à la question B.3, aucune conversion en base $10$ n'est attendue — le critère suffit.
\end{corrige}

\newpage

\coursec{Fiche bilan}

\begin{popfigure}
\begin{center}\tcbox[colback=popink,colframe=popink,coltext=white,arc=9pt,boxsep=4pt,left=6pt,right=6pt]{\bfseries\small Arithmétique : division euclidienne et numération}\end{center}
\vspace{-2pt}
\begin{tcbraster}[raster columns=3,raster equal height=rows,raster column skip=6pt,raster row skip=6pt,enhanced,blankest]
\begin{tcolorbox}[enhanced,colback=white,colframe=popblue,boxrule=0.9pt,arc=3pt,title={Division euclidienne},fonttitle=\bfseries\scriptsize,coltitle=white,colbacktitle=popblue,left=4pt,right=4pt,top=2pt,bottom=2pt,boxsep=1pt,toptitle=1.5pt,bottomtitle=1.5pt,halign=left]
\begin{itemize}[nosep,leftmargin=*,itemsep=1pt,font=\scriptsize]
\item $a=bq+r$
\item $0\leq r<b$
\end{itemize}
\end{tcolorbox}
\begin{tcolorbox}[enhanced,colback=white,colframe=poporange,boxrule=0.9pt,arc=3pt,title={Base $a$ : principe},fonttitle=\bfseries\scriptsize,coltitle=white,colbacktitle=poporange,left=4pt,right=4pt,top=2pt,bottom=2pt,boxsep=1pt,toptitle=1.5pt,bottomtitle=1.5pt,halign=left]
\begin{itemize}[nosep,leftmargin=*,itemsep=1pt,font=\scriptsize]
\item Chiffres $0$ à $a-1$
\item $N=\sum x_i a^i$
\end{itemize}
\end{tcolorbox}
\begin{tcolorbox}[enhanced,colback=white,colframe=popgreen,boxrule=0.9pt,arc=3pt,title={Conversions},fonttitle=\bfseries\scriptsize,coltitle=white,colbacktitle=popgreen,left=4pt,right=4pt,top=2pt,bottom=2pt,boxsep=1pt,toptitle=1.5pt,bottomtitle=1.5pt,halign=left]
\begin{itemize}[nosep,leftmargin=*,itemsep=1pt,font=\scriptsize]
\item Base $a\to 10$ : développer
\item $10\to$ base $a$ : divisions
\end{itemize}
\end{tcolorbox}
\begin{tcolorbox}[enhanced,colback=white,colframe=poppurple,boxrule=0.9pt,arc=3pt,title={Base à base},fonttitle=\bfseries\scriptsize,coltitle=white,colbacktitle=poppurple,left=4pt,right=4pt,top=2pt,bottom=2pt,boxsep=1pt,toptitle=1.5pt,bottomtitle=1.5pt,halign=left]
\begin{itemize}[nosep,leftmargin=*,itemsep=1pt,font=\scriptsize]
\item Passer par la base 10
\end{itemize}
\end{tcolorbox}
\begin{tcolorbox}[enhanced,colback=white,colframe=poppink,boxrule=0.9pt,arc=3pt,title={Addition en base $a$},fonttitle=\bfseries\scriptsize,coltitle=white,colbacktitle=poppink,left=4pt,right=4pt,top=2pt,bottom=2pt,boxsep=1pt,toptitle=1.5pt,bottomtitle=1.5pt,halign=left]
\begin{itemize}[nosep,leftmargin=*,itemsep=1pt,font=\scriptsize]
\item Retenue $=a$
\end{itemize}
\end{tcolorbox}
\begin{tcolorbox}[enhanced,colback=white,colframe=popgold,boxrule=0.9pt,arc=3pt,title={Multiplication en base $a$},fonttitle=\bfseries\scriptsize,coltitle=white,colbacktitle=popgold,left=4pt,right=4pt,top=2pt,bottom=2pt,boxsep=1pt,toptitle=1.5pt,bottomtitle=1.5pt,halign=left]
\begin{itemize}[nosep,leftmargin=*,itemsep=1pt,font=\scriptsize]
\item Tables en base $a$
\end{itemize}
\end{tcolorbox}
\end{tcbraster}

\legende{Carte mentale de la leçon : les six compétences clés autour de la division euclidienne et de la numération.}
\end{popfigure}

\begin{bilan}
\textbf{1. Division euclidienne dans $\mathbb{N}$}
\begin{itemize}
  \item \cle{Théorème} : pour $a\in\mathbb{N}$ et $b\in\mathbb{N}^*$, il existe un \textbf{unique} couple $(q,r)\in\mathbb{N}^2$ tel que
  \[ a=bq+r \quad\text{avec}\quad 0\leq r<b. \]
  $q$ est le \cle{quotient}, $r$ le \cle{reste}, $b$ le diviseur, $a$ le dividende.
  \item $b$ divise $a$ (on note $b\mid a$) si et seulement si $r=0$.
\end{itemize}

\textbf{2. Numération de base $a$ ($a\in\{2;3;\dots;10\}$)}
\begin{itemize}
  \item Tout entier $N\geq 1$ s'écrit de façon \textbf{unique} :
  \[ N=x_n a^n+x_{n-1}a^{n-1}+\dots+x_1 a+x_0,\quad 0\leq x_i<a,\ x_n\neq 0. \]
  \item On note $N=\overline{x_n x_{n-1}\dots x_1 x_0}^{\,a}$. Les \cle{chiffres} autorisés sont $0,1,\dots,a-1$.
  \item Bases usuelles : $2$ (binaire : $0,1$), $8$ (octal), $10$ (décimal).
\end{itemize}

\textbf{3. Conversions — méthodes express}
\begin{center}
\begin{tabularx}{\linewidth}{|l|X|}\hline
\rowcolor{popblueL} \textbf{Conversion} & \textbf{Méthode} \\ \hline
Base $a$ $\to$ base 10 & Développer : multiplier chaque chiffre par la puissance de $a$ de son rang, puis additionner. \\ \hline
Base 10 $\to$ base $a$ & Divisions euclidiennes successives par $a$ ; lire les restes \textbf{du dernier au premier}. \\ \hline
Base $a$ $\to$ base $b$ & Passer par la base 10 : $a\to 10$ puis $10\to b$. \\ \hline
\end{tabularx}
\end{center}

\textbf{4. Opérations en base $a$}
\begin{itemize}
  \item \cle{Addition} : on additionne colonne par colonne ; dès que la somme atteint $a$, on écrit le reste et on \textbf{retient le quotient} (retenue de $1$, $2$, \dots).
  \item \cle{Multiplication} : mêmes règles qu'en base 10, mais les tables et les retenues se calculent en base $a$.
  \item \cle{Vérification} : toujours contrôler le résultat en convertissant en base 10.
\end{itemize}
\end{bilan}

\begin{aretenir}[Checklist avant le Probatoire]
\begin{itemize}
  \item[$\square$] Je sais énoncer le théorème de la division euclidienne avec la condition $0\leq r<b$.
  \item[$\square$] Je sais déterminer le quotient et le reste de $a$ par $b$ et écrire $a=bq+r$.
  \item[$\square$] Je sais reconnaître si un entier est divisible par un autre (reste nul).
  \item[$\square$] Je connais les chiffres autorisés en base $a$ : de $0$ à $a-1$.
  \item[$\square$] Je sais décomposer $N$ en somme de puissances ordonnées de $a$.
  \item[$\square$] Je sais convertir un nombre de la base $a$ vers la base 10 (développement).
  \item[$\square$] Je sais convertir un nombre de la base 10 vers la base $a$ (divisions successives, restes lus à l'envers).
  \item[$\square$] Je sais passer d'une base $a$ à une base $b$ en transittant par la base 10.
  \item[$\square$] Je sais poser une addition en base $a$ avec la retenue correcte.
  \item[$\square$] Je sais poser une multiplication en base $a$ et vérifier en base 10.
  \item[$\square$] Je n'oublie jamais d'indiquer la base en exposant : $\overline{1011}^{\,2}\neq 1011$.
  \item[$\square$] Je vérifie que chaque chiffre utilisé est strictement inférieur à la base.
\end{itemize}
\end{aretenir}

\findecours

\end{document}
